% Lesson 38 — Grammar: Revision: Direct and indirect speech — Anglais Tle A — Cours signé OnBuch+
% Programme officiel MINESEC. Compiler avec : tectonic EN38-grammar-revision-direct-and-indirect-speech.tex
\newcommand{\DOCMATIERE}{Anglais}
\newcommand{\DOCNIVEAU}{Tle A}
\newcommand{\DOCMODULE}{Chapter 4 : Citizenship and Human Rights}
\newcommand{\DOCLECON}{EN38}
\newcommand{\DOCTITRE}{Lesson 38 — Grammar: Revision: Direct and indirect speech}
% OnBuch+ — préambule « cours signés » ENRICHI (compilé avec tectonic / XeLaTeX)
% Système visuel « Pop » d'OnBuch+ : fond crème (jamais blanc), bordures encre
% épaisses, ombres « dures » décalées, titres Archivo Black en capitales.
% Version MATHÉMATIQUES : graphes (pgfplots/tikz), boîtes pédagogiques
% (définition, propriété, méthode, attention, le savais-tu, exercice résolu,
% exercice, TP, bilan), page de garde et signature OnBuch+ ancrée en bas.
%
% Macros attendues dans main.tex AVANT \input{preamble} :
%   \DOCMATIERE  \DOCNIVEAU  \DOCMODULE  \DOCLECON (numéro)  \DOCTITRE
\documentclass[11pt]{article}

\usepackage{fontspec}
\usepackage[english,french]{babel} % français = langue principale ; l'anglais via \eng{..} / textebox
\usepackage[a4paper,margin=2cm,top=3.4cm,bottom=2.8cm,headheight=1.4cm,footskip=1.3cm]{geometry}
\usepackage{amsmath,amssymb,amsfonts}
\usepackage{mathtools} % \xmapsto, \coloneqq... souvent utilisés par les modèles
\usepackage{cancel} % simplifications barrées (bilans de forces)
% \abs{z} (module, valeur absolue) et \norm{u} (norme), utilisés par les modèles
\providecommand{\abs}{}\let\abs\relax\DeclarePairedDelimiter{\abs}{\lvert}{\rvert}
\providecommand{\norm}{}\let\norm\relax\DeclarePairedDelimiter{\norm}{\lVert}{\rVert}
\usepackage[table]{xcolor} % [table] : fournit aussi \rowcolors (alternance de lignes)
\usepackage{pagecolor}
\usepackage{tikz}
\usetikzlibrary{arrows.meta,positioning,calc,shapes.geometric,shapes.misc,shapes.arrows,decorations.pathmorphing,decorations.pathreplacing,decorations.markings,patterns,shadows,mindmap,trees,fit,backgrounds,matrix,intersections,angles,quotes}
\usepackage{pgfplots}
\usepackage[version=4]{mhchem} % équations nucléaires \ce{^{235}_{92}U}
\usepackage[european,siunitx]{circuitikz} % schémas électriques
\pgfplotsset{compat=1.18}
\usepgfplotslibrary{fillbetween,groupplots} % aires entre courbes (calcul intégral)
\usepackage{enumitem}
\usepackage{fancyhdr}
% [most] charge toutes les bibliothèques tcolorbox (listings, minted, raster,
% documentation...) et gonfle inutilement la mémoire TeX sur un document long
% (~300 boîtes) : on ne charge que ce qui est réellement utilisé.
\usepackage[skins,breakable]{tcolorbox}
\usepackage{graphicx}
\usepackage{colortbl}
\usepackage{booktabs}
\usepackage{tabularx}
\usepackage{diagbox} % case d'en-tête diagonale des tableaux à double entrée
% Colonnes extensibles centrée / gauche / droite (C, L, R), souvent utilisées par les modèles
\newcolumntype{C}{>{\centering\arraybackslash}X}
\newcolumntype{L}{>{\raggedright\arraybackslash}X}
\newcolumntype{R}{>{\raggedleft\arraybackslash}X}
\usepackage{array}
\usepackage{multicol}
\usepackage{multirow}
\usepackage{siunitx}
% parse-numbers=false : accepte aussi \SI{2,5 \times 10^{-3}}{...}
\sisetup{locale=FR,output-decimal-marker={,},group-minimum-digits=4,parse-numbers=false}
\DeclareSIUnit\ohm{\ensuremath{\symup{\Omega}}} % U+2126 absent de la police
\DeclareSIUnit\division{div} % réglages d'oscilloscope (ms/div, V/div)
\DeclareSIUnit\tr{tr} % tours (vitesse de rotation, ex. \SI{1500}{\tr\per\minute})
\DeclareSIUnit\molar{M} % concentration molaire (ex. \SI{3}{\molar}, \SI{1}{\micro\molar})
\DeclareSIUnit\an{an} % année (le modèle confond parfois avec \year, un compteur TeX) — ex. \SI{5}{\kilo\meter\per\an}
\DeclareSIUnit\FCFA{FCFA} % franc CFA (ex. \SI{500}{\FCFA\per\kilo\gram})
\DeclareSIUnit\degreeBrix{\textdegree Brix} % teneur en sucre d'un sirop/jus (réfractométrie)
\DeclareSIUnit\month{mois} % le modèle utilise parfois \month, un compteur TeX, comme unité de durée
\usepackage{qrcode}
% \degree : commande fréquemment utilisée par les modèles pour un angle ou une
% température en dehors de \SI{}{} (non fournie par les paquets chargés).
\providecommand{\degree}{^{\circ}}
% \qed (fin de démonstration), utilisé par les modèles sans amsthm
\providecommand{\qed}{\hfill$\square$}
% \barre : double barre des valeurs interdites dans un tableau de variations (tkz-tab)
\providecommand{\barre}{\ensuremath{\|}}
% environnement proof (amsthm non chargé) utilisé par les modèles
\ifdefined\proof\else\newenvironment{proof}[1][Démonstration]{\par\noindent\textit{#1.}\ }{\qed\par}\fi
\usepackage{lastpage}
\usepackage{hyperref}
\hypersetup{hidelinks}

% ============================================================
% Palette « Pop » OnBuch+ — mêmes hex que la classe P (plus_ui.dart)
% ============================================================
\definecolor{poporange}{HTML}{F59321}
\definecolor{popdark}{HTML}{E07A0C}
\definecolor{popcream}{HTML}{FFF6E8}
\definecolor{popink}{HTML}{1C1714}
\definecolor{poppurple}{HTML}{7A5AE0}
\definecolor{popgreen}{HTML}{2E9E6A}
\definecolor{poppink}{HTML}{E0457B}
\definecolor{popblue}{HTML}{2F7DE1}
\definecolor{popgold}{HTML}{C98A12}
\definecolor{popmuted}{HTML}{8A7F76}
\definecolor{popline}{HTML}{EADFD2}
\definecolor{popgray}{HTML}{8A7F76}
\colorlet{popgrey}{popgray}
% Alias tolérants (le modèle nomme parfois les couleurs différemment)
\colorlet{popwhite}{white}
\colorlet{popblack}{popink}
\colorlet{popred}{poppink}
\colorlet{popyellow}{popgold}
% Teintes claires (fonds de tableaux/encadrés) — le modèle nomme parfois
% les couleurs avec un suffixe L (ex. popblueL) sans les avoir définies.
\colorlet{poporangeL}{poporange!20}
\colorlet{popdarkL}{popdark!20}
\colorlet{poppurpleL}{poppurple!20}
\colorlet{popgreenL}{popgreen!20}
\colorlet{poppinkL}{poppink!20}
\colorlet{popblueL}{popblue!20}
\colorlet{popgoldL}{popgold!20}
\colorlet{popredL}{popred!20}
\colorlet{popgrayL}{popgray!20}
% Teintes claires pour fonds de tableaux / zones
\colorlet{popblueL}{popblue!10!white}
\colorlet{poporangeL}{poporange!14!white}
\colorlet{popgreenL}{popgreen!12!white}
\colorlet{poppurpleL}{poppurple!12!white}
\colorlet{poppinkL}{poppink!12!white}
\colorlet{popgoldL}{popgold!14!white}

\pagecolor{popcream}

% ============================================================
% Typographie
% ============================================================
\defaultfontfeatures{Ligatures=TeX}
\setmainfont{PlusJakartaSans-Medium.ttf}[Path=../../../fonts/,BoldFont=PlusJakartaSans-Bold.ttf]
% Symboles, API et emojis absents de la police principale : repli sur DejaVu Sans (caractères actifs)
\newfontface\symfont{DejaVuSans.ttf}[Path=../../../fonts/]
\makeatletter
\newcommand\symactive[1]{\catcode#1=13 \begingroup\lccode`\~=#1 \lowercase{\endgroup\def~}{{\symfont\char#1}}}
\newcommand\symblank[1]{\catcode#1=13 \begingroup\lccode`\~=#1 \lowercase{\endgroup\def~}{}}
\newcommand\symhyph[1]{\catcode#1=13 \begingroup\lccode`\~=#1 \lowercase{\endgroup\def~}{\mbox{-}}}
\newcount\symc
\newcommand\symrange[2]{\symc=#1 \@whilenum\symc<#2 \do{\expandafter\symactive\expandafter{\the\symc}\advance\symc by 1 }}
\newcommand\symblankrange[2]{\symc=#1 \@whilenum\symc<#2 \do{\expandafter\symblank\expandafter{\the\symc}\advance\symc by 1 }}
\makeatother
\symrange{"0250}{"02FF}\symactive{"2713}\symactive{"2714}\symactive{"2717}\symactive{"2718}\symactive{"2610}\symactive{"2611}\symactive{"2612}
\symhyph{"2011}\symhyph{"2E1A}\symblankrange{"1F300}{"1FAFF}\symblank{"FE0F}
% Maths en sans-serif (Fira Math) avec chiffres et lettres droites de Plus
% Jakarta, pour que les formules se fondent dans le texte.
\usepackage[math-style=ISO,bold-style=ISO]{unicode-math}
\setmathfont{FiraMath-Regular.otf}[Path=../../../fonts/]
\setmathfont{PlusJakartaSans-Medium.ttf}[Path=../../../fonts/,range={up/num,up/{Latin,latin},\mathup}]
\usepackage{icomma} % virgule décimale sans espace en mode math
% Caractères mathématiques Unicode absents des polices de texte (Plus Jakarta,
% Archivo Black) : rendus par leur équivalent mathématique, en texte comme en
% formule — sinon ils s'affichent en carré vide (ex. « Arithmétique dans ℤ »).
\usepackage{newunicodechar}
\newunicodechar{ℝ}{\ensuremath{\mathbb{R}}}
\newunicodechar{ℤ}{\ensuremath{\mathbb{Z}}}
\newunicodechar{ℂ}{\ensuremath{\mathbb{C}}}
\newunicodechar{ℕ}{\ensuremath{\mathbb{N}}}
\newunicodechar{ℚ}{\ensuremath{\mathbb{Q}}}
\newunicodechar{𝔻}{\ensuremath{\mathbb{D}}}
\newunicodechar{α}{\ensuremath{\alpha}}
\newunicodechar{β}{\ensuremath{\beta}}
\newunicodechar{γ}{\ensuremath{\gamma}}
\newunicodechar{δ}{\ensuremath{\delta}}
\newunicodechar{ε}{\ensuremath{\varepsilon}}
\newunicodechar{θ}{\ensuremath{\theta}}
\newunicodechar{λ}{\ensuremath{\lambda}}
\newunicodechar{μ}{\ensuremath{\mu}}
\newunicodechar{σ}{\ensuremath{\sigma}}
\newunicodechar{τ}{\ensuremath{\tau}}
\newunicodechar{φ}{\ensuremath{\varphi}}
\newunicodechar{ψ}{\ensuremath{\psi}}
\newunicodechar{ω}{\ensuremath{\omega}}
\newunicodechar{Σ}{\ensuremath{\Sigma}}
% majuscules produites par \MakeUppercase dans les titres (α-aminés -> Α)
\newunicodechar{Α}{\ensuremath{\alpha}}
\newunicodechar{Β}{\ensuremath{\beta}}
\newunicodechar{Γ}{\ensuremath{\Gamma}}
\newunicodechar{Δ}{\ensuremath{\Delta}}
\newunicodechar{Ε}{\ensuremath{\varepsilon}}
\newunicodechar{Θ}{\ensuremath{\Theta}}
\newunicodechar{Λ}{\ensuremath{\Lambda}}
\newunicodechar{Μ}{\ensuremath{\mu}}
\newunicodechar{Τ}{\ensuremath{\tau}}
\newunicodechar{Φ}{\ensuremath{\Phi}}
\newunicodechar{Ψ}{\ensuremath{\Psi}}
\newunicodechar{Ω}{\ensuremath{\Omega}}
\newunicodechar{↦}{\ensuremath{\mapsto}}
\newunicodechar{∈}{\ensuremath{\in}}
\newunicodechar{∉}{\ensuremath{\notin}}
\newunicodechar{∧}{\ensuremath{\wedge}}
\newunicodechar{⇔}{\ensuremath{\Leftrightarrow}}
\newunicodechar{⟺}{\ensuremath{\Longleftrightarrow}}
\newunicodechar{⇒}{\ensuremath{\Rightarrow}}
\newunicodechar{⟹}{\ensuremath{\Longrightarrow}}
\newunicodechar{∘}{\ensuremath{\circ}}
\newunicodechar{∪}{\ensuremath{\cup}}
\newunicodechar{∩}{\ensuremath{\cap}}
\newunicodechar{≡}{\ensuremath{\equiv}}
\newunicodechar{⊂}{\ensuremath{\subset}}
\newunicodechar{⊥}{\ensuremath{\perp}}
\newunicodechar{∃}{\ensuremath{\exists}}
\newunicodechar{∀}{\ensuremath{\forall}}
\newunicodechar{∛}{\ensuremath{\sqrt[3]{\,}}}
\newunicodechar{∜}{\ensuremath{\sqrt[4]{\,}}}
\newunicodechar{⃗}{\ensuremath{{}^{\to}}}
\newunicodechar{ⁿ}{\ifmmode^{n}\else\textsuperscript{\ensuremath{n}}\fi}
\newunicodechar{ˣ}{\ifmmode^{x}\else\textsuperscript{\ensuremath{x}}\fi}
\newunicodechar{ᵇ}{\ifmmode^{b}\else\textsuperscript{\ensuremath{b}}\fi}
\newunicodechar{ʳ}{\ifmmode^{r}\else\textsuperscript{\ensuremath{r}}\fi}
\newunicodechar{ᵘ}{\ifmmode^{u}\else\textsuperscript{\ensuremath{u}}\fi}
\newunicodechar{ʸ}{\ifmmode^{y}\else\textsuperscript{\ensuremath{y}}\fi}
\newunicodechar{ᵃ}{\ifmmode^{a}\else\textsuperscript{\ensuremath{a}}\fi}
\newunicodechar{ᵉ}{\ifmmode^{e}\else\textsuperscript{\ensuremath{e}}\fi}
\newunicodechar{ᵏ}{\ifmmode^{k}\else\textsuperscript{\ensuremath{k}}\fi}
\newunicodechar{ᵖ}{\ifmmode^{p}\else\textsuperscript{\ensuremath{p}}\fi}
\newunicodechar{ᴺ}{\ifmmode^{N}\else\textsuperscript{\ensuremath{N}}\fi}
\newunicodechar{ᶜ}{\ifmmode^{c}\else\textsuperscript{\ensuremath{c}}\fi}
\newunicodechar{ʰ}{\ifmmode^{h}\else\textsuperscript{\ensuremath{h}}\fi}
\newunicodechar{ᵠ}{\ifmmode^{q}\else\textsuperscript{\ensuremath{q}}\fi}
\newunicodechar{ₙ}{\ifmmode_{n}\else\textsubscript{\ensuremath{n}}\fi}
\newunicodechar{ₐ}{\ifmmode_{a}\else\textsubscript{\ensuremath{a}}\fi}
\newunicodechar{ᵢ}{\ifmmode_{i}\else\textsubscript{\ensuremath{i}}\fi}
\newunicodechar{ᵣ}{\ifmmode_{r}\else\textsubscript{\ensuremath{r}}\fi}
\newunicodechar{ₖ}{\ifmmode_{k}\else\textsubscript{\ensuremath{k}}\fi}
\newunicodechar{ᵦ}{\ifmmode_{\beta}\else\textsubscript{\ensuremath{\beta}}\fi}
\newunicodechar{ₓ}{\ifmmode_{x}\else\textsubscript{\ensuremath{x}}\fi}
\newfontfamily\popdisplay{ArchivoBlack-Regular.ttf}[Path=../../../fonts/]
\newfontfamily\poplogo{SpaceGrotesk-Bold.ttf}[Path=../../../fonts/]
\newfontfamily\popsemi{PlusJakartaSans-SemiBold.ttf}[Path=../../../fonts/]
\newfontfamily\popxbold{PlusJakartaSans-ExtraBold.ttf}[Path=../../../fonts/]
\newfontfamily\popxboldls{PlusJakartaSans-ExtraBold.ttf}[Path=../../../fonts/,LetterSpace=3.0]

\setlist[enumerate]{leftmargin=1.7em,itemsep=5pt,topsep=4pt}
\setlist[itemize]{leftmargin=1.7em,itemsep=4pt,topsep=4pt,label={\color{poporange}\textbullet}}
\setlength{\parindent}{0pt}
\setlength{\parskip}{5pt}
\renewcommand{\arraystretch}{1.25}

% pgfplots : style Pop par défaut
\pgfplotsset{
  every axis/.append style={axis line style={popink,line width=1pt},tick style={popink},
    grid style={popline,line width=0.5pt},label style={font=\small},tick label style={font=\footnotesize}},
  popcurve/.style={poporange,line width=1.8pt,no markers,smooth},
}
% Même style disponible dans un \draw TikZ simple (hors pgfplots)
\tikzset{popcurve/.style={poporange,line width=1.8pt}}
\tikzset{popgrid/.style={popline,very thin}}
\colorlet{popcurve}{poporange} % le nom du style est parfois utilisé comme couleur

% ============================================================
% Pill / squiggle
% ============================================================
\newcommand{\poppill}[3]{%
  \tikz[baseline=-0.55ex]\node[draw=popink,line width=1.2pt,rounded corners=2.6pt,%
    fill=#2,inner xsep=6.5pt,inner ysep=3pt]{%
    \popxboldls\fontsize{7}{7}\selectfont\color{#3}\MakeUppercase{#1}};%
}
\newcommand{\popsquiggle}[1]{%
  \tikz[baseline=-0.4ex]{
    \draw[#1,line width=2.6pt,line cap=round]
      plot [smooth] coordinates {(0,0.09) (0.34,0.01) (0.68,0.21) (1.02,0.01) (1.36,0.10)};
  }%
}
% Mot-clé surligné (marqueur orange sous le texte)
\newcommand{\cle}[1]{{\popxbold\color{popdark}#1}}

% ============================================================
% En-tête / pied
% ============================================================
\pagestyle{fancy}
\fancyhf{}
\renewcommand{\headrulewidth}{0pt}
\renewcommand{\footrulewidth}{0pt}
\lhead{\raisebox{-0.15ex}{\poplogo\fontsize{13}{13}\selectfont\color{popink}OnBuch\color{poporange}+}}
\rhead{\poppill{\DOCMATIERE\ \textbullet\ \DOCNIVEAU\ \textbullet\ Leçon \DOCLECON}{white}{popink}}
\lfoot{\popxbold\fontsize{7.6}{7.6}\selectfont\color{popmuted}\MakeUppercase{Cours signé OnBuch+}\ \textbullet\ {\popsemi\DOCTITRE}}
\rfoot{\tikz[baseline=-0.6ex]\node[rounded corners=8.5pt,fill=popink,minimum height=17pt,minimum width=17pt,inner xsep=5pt,inner sep=0]{\popxbold\fontsize{8}{8}\selectfont\color{popcream}\thepage\,/\,\pageref*{LastPage}};}

% ============================================================
% Titres
% ============================================================
\newcommand{\chaptitre}[2]{%
  \begin{tcolorbox}[enhanced,colback=poporange,colframe=popink,boxrule=2.6pt,arc=11pt,
      drop shadow=popink,left=15pt,right=15pt,top=12pt,bottom=11pt,before skip=3pt,after skip=18pt]
    \poppill{#2}{popink}{popcream}\par\vspace{9pt}
    {\raggedright\hyphenpenalty=10000\exhyphenpenalty=10000\popdisplay\fontsize{22}{24}\selectfont\color{popcream}\MakeUppercase{#1}\par}
  \end{tcolorbox}
}
% Section numérotée : pastille orange avec le numéro + titre Archivo
\newcounter{popsec}
\newcommand{\coursec}[1]{%
  \stepcounter{popsec}\setcounter{popsub}{0}%
  \par\vspace{16pt}\noindent
  \begin{minipage}{\linewidth}
  \tikz[baseline=-0.7ex]\node[draw=popink,line width=1.6pt,fill=poporange,rounded corners=4pt,minimum size=22pt,inner sep=0]{\popdisplay\fontsize{12}{12}\selectfont\color{popcream}\thepopsec};%
  \hspace{8pt}{\popdisplay\fontsize{15}{17}\selectfont\color{popink}\MakeUppercase{#1}}\par
  \vspace{4pt}\noindent{\color{poporange}\rule{36pt}{3.6pt}}
  \end{minipage}\par\nopagebreak\vspace{8pt}
}
\newcounter{popsub}
\newcommand{\courssub}[1]{%
  \stepcounter{popsub}%
  \par\vspace{9pt}\noindent{\popxbold\fontsize{11.5}{13}\selectfont\color{popink}{\color{poporange}\thepopsec.\thepopsub}\hspace{6pt}#1}\par\nopagebreak\vspace{4pt}
}

% ============================================================
% Boîtes pédagogiques — même recette : fond blanc, bordure épaisse colorée,
% ombre dure encre, badge plein. Titre optionnel [#1].
% ============================================================
\tcbset{popbase/.style={breakable,enhanced,colback=white,boxrule=2.3pt,arc=9pt,
  shadow={3pt}{-3pt}{0pt}{popink},left=14pt,right=14pt,top=11pt,bottom=11pt,before skip=11pt,after skip=15pt,
  fontupper=\popsemi\fontsize{10.6}{14.5}\selectfont\color{popink},
  fontlower=\popsemi\fontsize{10.6}{14.5}\selectfont\color{popink}}}
% Coche dessinée (le glyphe ✓ n'existe pas dans les polices du document)
\newcommand{\popcheck}[1]{\tikz[baseline=-0.5ex]\draw[#1,line width=1.6pt,line cap=round,line join=round](-0.13,0.0)--(-0.04,-0.09)--(0.13,0.1);}
\providecommand{\coche}{\popcheck{popgreen}}
\providecommand{\resultat}[1]{\par\smallskip\noindent\fcolorbox{popgreen}{popgreenL}{\parbox{\dimexpr\linewidth-2\fboxsep-2\fboxrule\relax}{\textbf{Résultat.} #1}}\par\smallskip}
\newcommand{\popboxhead}[3]{\poppill{#1}{#2}{white}\ifx\relax#3\relax\else\hspace{6pt}{\popxbold\fontsize{10.6}{12}\selectfont\color{popink}#3}\fi\par\vspace{7pt}}

\newtcolorbox{aretenir}[1][]{popbase,colframe=popgreen,before upper={\popboxhead{À retenir}{popgreen}{#1}}}
% Texte en anglais (document de lecture, dialogue, extrait) : cadre
% d'étude. Un titre optionnel (source, auteur) s'affiche dans l'en-tête.
\newtcolorbox{textebox}[1][]{popbase,colframe=popblue,before upper={\popboxhead{Text}{popblue}{#1}\begin{otherlanguage}{english}},after upper={\end{otherlanguage}}}
% Marques ✓ / ✗ (forme correcte / fautive) souvent utilisées par les modèles
\providecommand{\cross}{\textcolor{poppink}{\ensuremath{\times}}}
\providecommand{\xmark}{\cross}\providecommand{\crossmark}{\cross}\providecommand{\cmark}{\ensuremath{\checkmark}}
% Ligne de réponse / justification à compléter (inventée par les modèles)
\providecommand{\justification}[1]{\par\noindent\textit{Justification~:} \dotfill\par}
% \textipa (package tipa non chargé) : la police ne couvre pas l'API -> simple passage
\providecommand{\textipa}[1]{#1}
% Soulignement (ulem non chargé) : \uline, \ul
\providecommand{\uline}[1]{\underline{#1}}\providecommand{\ul}[1]{\underline{#1}}
% \glossaire{\item ...} inventée par les modèles : liste de glossaire
\providecommand{\glossaire}[1]{\begin{itemize}#1\end{itemize}}
% \alert (beamer) inventée par les modèles : texte mis en évidence
\providecommand{\alert}[1]{\textbf{\textcolor{poppink}{#1}}}
% variantes ASCII d'ordinaux écrites en macro
\providecommand{\iere}{\textsuperscript{re}}\providecommand{\ieme}{\textsuperscript{e}}\providecommand{\ere}{\textsuperscript{re}}\providecommand{\eme}{\textsuperscript{e}}\providecommand{\er}{\textsuperscript{er}}\providecommand{\ie}{\textsuperscript{e}}
% Ordinaux français écrits en macro par les modèles : 1\ière, 3\ième
\newcommand{\ière}{\textsuperscript{re}}\newcommand{\ième}{\textsuperscript{e}}
% \exemple (inventée par les modèles) : étiquette « Exemple : »
\providecommand{\exemple}{\textbf{Exemple~:}\ }
% \square utilisable aussi hors mode math (cases à cocher)
\AtBeginDocument{\let\squareMath\square\renewcommand{\square}{\ensuremath{\squareMath}}}
% Mot ou phrase en anglais à l'intérieur d'un paragraphe français
\newcommand{\eng}[1]{\ifmmode\text{\textit{#1}}\else\begingroup\selectlanguage{english}\itshape #1\endgroup\fi}
\let\esp\eng % alias toléré
% flèches d'intonation inventées par les modèles
\providecommand{\textsearrow}{}\renewcommand{\textsearrow}{\ensuremath{\searrow}}\providecommand{\textnearrow}{}\renewcommand{\textnearrow}{\ensuremath{\nearrow}}
% \fr{..} : mot français (inventé par les modèles)
\providecommand{\fr}[1]{\textit{#1}}
\newtcolorbox{exemplebox}[1][]{popbase,colframe=poporange,before upper={\popboxhead{Exemple}{poporange}{#1}}}
\newtcolorbox{definition}[1][]{popbase,colframe=popblue,before upper={\popboxhead{Définition}{popblue}{#1}}}
\newtcolorbox{propriete}[1][]{popbase,colframe=poppurple,before upper={\popboxhead{Propriété}{poppurple}{#1}}}
\newtcolorbox{methode}[1][]{popbase,colframe=popgold,before upper={\popboxhead{Méthode}{popgold}{#1}}}
\newtcolorbox{attention}[1][]{popbase,colframe=poppink,before upper={\popboxhead{Attention}{poppink}{#1}}}
\newtcolorbox{experience}[1][]{popbase,colframe=popblue,colback=popblueL,before upper={\popboxhead{Expérience}{popblue}{#1}}}
% « Le savais-tu ? » : carte violette pleine (contexte, culture, Cameroun)
\newtcolorbox{savaistu}[1][]{popbase,colback=poppurple,colframe=popink,
  fontupper=\popsemi\fontsize{10.4}{14.2}\selectfont\color{white},
  before upper={\poppill{Le savais-tu\,?}{white}{poppurple}\ifx\relax#1\relax\else\hspace{6pt}{\popxbold\color{white}#1}\fi\par\vspace{7pt}}}
% Exercice résolu : énoncé en haut, \tcblower puis la solution sur fond crème
\newcounter{popexo}\newcounter{popexr}
\newtcolorbox{exoresolu}[1][]{popbase,colframe=poporange,colbacklower=popcream,
  segmentation style={popink,line width=1.2pt,dashed},
  before upper={\stepcounter{popexr}\popboxhead{Exercice résolu \thepopexr}{poporange}{#1}},
  before lower={\poppill{Solution}{popink}{popcream}\par\vspace{6pt}}}
% Exercice (non résolu en place) — la correction est dans la partie Corrigés
\newtcolorbox{exercice}[1][]{popbase,colframe=popink,
  before upper={\stepcounter{popexo}\popboxhead{Exercice \thepopexo}{popink}{#1}}}
% Corrigé
\newtcolorbox{corrige}[1][]{popbase,colframe=popgreen,colback=popgreenL,
  before upper={\popboxhead{Corrigé}{popgreen}{#1}}}
% Objectifs / prérequis / bilan
\newtcolorbox{objectifs}{popbase,colframe=popink,colback=poporangeL,
  before upper={\popboxhead{Objectifs de la leçon}{popink}{}}}
\newtcolorbox{prerequis}{popbase,colframe=popink,colback=popblueL,
  before upper={\popboxhead{Prérequis}{popblue}{}}}
\newtcolorbox{bilan}{popbase,colframe=popink,colback=white,boxrule=2.8pt,
  before upper={\popboxhead{Fiche bilan}{poporange}{L'essentiel à connaître pour l'examen}}}
% Figure encadrée avec légende
\newtcolorbox{popfigure}[1][]{enhanced,breakable=false,colback=white,colframe=popline,boxrule=1.4pt,arc=8pt,
  left=10pt,right=10pt,top=10pt,bottom=8pt,before skip=10pt,after skip=12pt,halign=center,
  lower separated=false,halign lower=center,
  fontlower=\popsemi\fontsize{9}{11.5}\selectfont\color{popmuted},
  #1}
\newcommand{\legende}[1]{\par\vspace{6pt}{\popsemi\fontsize{9}{11.5}\selectfont\color{popmuted}#1}}

% ============================================================
% Signature OnBuch+
% ============================================================
\newcommand{\poplogoinline}[1]{{\poplogo\fontsize{#1}{#1}\selectfont\color{popink}OnBuch\color{poporange}+}}
\newcommand{\signatureonbuch}{%
  \begin{tcolorbox}[enhanced,colback=popink,colframe=popink,boxrule=2.2pt,arc=10pt,
      drop shadow=poporange,left=14pt,right=14pt,top=10pt,bottom=10pt,before skip=0pt,after skip=0pt]
    \tikz[baseline=-0.7ex]\node[circle,fill=popgreen,draw=popcream,line width=1.3pt,minimum size=22pt,inner sep=0]{\popcheck{white}};%
    \hspace{9pt}\begin{minipage}[c]{0.62\linewidth}
      {\popxbold\fontsize{12}{13}\selectfont\color{popcream}Signé\ \textbullet\ {\poplogo OnBuch\color{poporange}+}}\par\vspace{2pt}
      {\raggedright\popsemi\fontsize{8}{10.5}\selectfont\color{popline}\MakeUppercase{Équipe pédagogique OnBuch+}\\\MakeUppercase{Conforme au programme officiel MINESEC}\par}
    \end{minipage}\hfill
    {\poplogo\fontsize{22}{22}\selectfont\color{popcream}OnBuch\color{poporange}+}
  \end{tcolorbox}%
}

% ============================================================
% Page de garde
% #1 titre · #2 matière · #3 niveau · #4 module · #5 n° leçon · #6 durée ·
% #7 description / objectifs (texte ou itemize)
% ============================================================
\newcommand{\pagegarde}[7]{%
  \thispagestyle{empty}
  \newgeometry{a4paper,margin=1.7cm,top=1.6cm,bottom=1.4cm}%
  \begin{tikzpicture}[remember picture,overlay]
    \fill[poporange!18] ([xshift=-3.2cm,yshift=-3.6cm]current page.north east) circle (4.6cm);
    \fill[poppurple!14] ([xshift=2.4cm,yshift=7.6cm]current page.south west) circle (3.8cm);
  \end{tikzpicture}%
  {\poplogo\fontsize{30}{30}\selectfont\color{popink}OnBuch\color{poporange}+}\hfill
  \raisebox{0.6ex}{\poppill{Cours signé \textbullet\ Premium}{popink}{popcream}}\par
  \vspace{1pt}\popsquiggle{poppurple}\par
  \vspace{8pt}
  \begin{tcolorbox}[enhanced,colback=poporange,colframe=popink,boxrule=2.8pt,arc=13pt,
      drop shadow=popink,left=18pt,right=18pt,top=12pt,bottom=13pt,after skip=10pt]
    \poppill{#2 \textbullet\ #3}{popink}{popcream}\hspace{5pt}\poppill{#4}{white}{popink}\par\vspace{8pt}
    {\popdisplay\fontsize{14}{14}\selectfont\color{popink}LEÇON #5}\par\vspace{4pt}
    {\raggedright\hyphenpenalty=10000\exhyphenpenalty=10000\popdisplay\fontsize{25}{27}\selectfont\color{popcream}\MakeUppercase{#1}\par}
    \vspace{8pt}
    {\popxbold\fontsize{9.5}{9.5}\selectfont\color{popink}\MakeUppercase{Durée conseillée : #6}}
  \end{tcolorbox}
  \begin{tcolorbox}[enhanced,colback=white,colframe=popink,boxrule=2.2pt,arc=10pt,
      drop shadow=popink,left=16pt,right=16pt,top=10pt,bottom=10pt,after skip=8pt]
    \poppill{Dans cette leçon}{poppurple}{white}\par\vspace{5pt}
    \popsemi\fontsize{9.3}{12.6}\selectfont\color{popink}#7
  \end{tcolorbox}
  \noindent\begin{minipage}[t]{0.487\linewidth}
    \begin{tcolorbox}[enhanced,colback=popgreen,colframe=popink,boxrule=2.2pt,arc=10pt,
        drop shadow=popink,left=11pt,right=11pt,top=10pt,bottom=10pt,height=3.1cm,valign=center]
      \tcbox[colback=white,colframe=popink,boxrule=1.3pt,arc=5pt,boxsep=0pt,nobeforeafter,
        left=3pt,right=3pt,top=3pt,bottom=3pt,valign=center]{\qrcode[height=1.9cm]{https://play.google.com/store/apps/details?id=com.luvvix.onbuch&hl=fr}}%
      \hspace{8pt}\begin{minipage}[c]{0.5\linewidth}
        {\popdisplay\fontsize{11}{12.5}\selectfont\color{white}\MakeUppercase{Télécharge OnBuch}}\par\vspace{4pt}
        {\raggedright\popsemi\fontsize{8.4}{11}\selectfont\color{white}Gratuit \textbullet\ Google Play\\Tuteur IA, annales, concours, résultats\par}
      \end{minipage}
    \end{tcolorbox}
  \end{minipage}\hfill
  \begin{minipage}[t]{0.487\linewidth}
    \begin{tcolorbox}[enhanced,colback=poppurple,colframe=popink,boxrule=2.2pt,arc=10pt,
        drop shadow=popink,left=11pt,right=11pt,top=10pt,bottom=10pt,height=3.1cm,valign=center]
      \tikz[baseline=-0.7ex]\node[circle,fill=white,draw=popink,line width=1.3pt,minimum size=1.6cm,inner sep=0]{\popdisplay\fontsize{24}{24}\selectfont\color{poppurple}+};%
      \hspace{8pt}\begin{minipage}[c]{0.6\linewidth}
        {\popdisplay\fontsize{11}{12.5}\selectfont\color{white}\MakeUppercase{Passe à OnBuch+}}\par\vspace{4pt}
        {\raggedright\popsemi\fontsize{8.4}{11}\selectfont\color{white}Tous les cours signés, Tuteur IA illimité, fiches PDF — onglet OnBuch+ de l'app\par}
      \end{minipage}
    \end{tcolorbox}
  \end{minipage}
  \vspace{8pt}
  \popcontenu
  \vfill
  \signatureonbuch
  \restoregeometry
  \newpage
}

% Rangée « Ce cours contient » (page de garde)
\newcommand{\poptile}[3]{% #1 couleur #2 gros texte #3 légende
  \begin{tcolorbox}[enhanced,colback=#1,colframe=popink,boxrule=2pt,arc=9pt,shadow={2.5pt}{-2.5pt}{0pt}{popink},
      left=6pt,right=6pt,top=9pt,bottom=9pt,halign=center,valign=center,height=1.75cm,width=\linewidth,nobeforeafter]
    {\popdisplay\fontsize{12.5}{13}\selectfont\color{white}\MakeUppercase{#2}}\par\vspace{3pt}
    {\centering\popsemi\fontsize{7.8}{9.5}\selectfont\color{white}#3\par}
  \end{tcolorbox}}
\newcommand{\popcontenu}{%
  \par\noindent{\popxboldls\fontsize{7.5}{7.5}\selectfont\color{popmuted}CE COURS SIGNÉ CONTIENT}\par\vspace{7pt}
  \noindent\begin{minipage}[t]{0.235\linewidth}\poptile{popblue}{Cours}{complet, illustré, pas à pas}\end{minipage}\hfill
  \begin{minipage}[t]{0.235\linewidth}\poptile{poporange}{Méthodes}{et exercices résolus}\end{minipage}\hfill
  \begin{minipage}[t]{0.235\linewidth}\poptile{poppink}{Exercices}{type examen + corrigés}\end{minipage}\hfill
  \begin{minipage}[t]{0.235\linewidth}\poptile{popgreen}{Fiche bilan}{l'essentiel à retenir}\end{minipage}\par}

% Sommaire
\newtcolorbox{sommaire}{enhanced,colback=white,colframe=popink,boxrule=2.2pt,arc=10pt,
  drop shadow=popink,left=18pt,right=18pt,top=13pt,bottom=13pt,before skip=6pt,after skip=20pt,
  before upper={\poppill{Sommaire}{poppurple}{white}\par\vspace{9pt}\popsemi\fontsize{10.6}{14.5}\selectfont\color{popink}}}

% Signature de fin de cours — poussée tout en bas de la dernière page
\newcommand{\findecours}{%
  \par\vfill\nopagebreak
  \begin{center}\popsquiggle{poporange}\end{center}\vspace{4pt}
  \signatureonbuch
}
\AtBeginDocument{\def\llbracket{\mathopen{[\mkern-3.5mu[}}\def\rrbracket{\mathclose{]\mkern-3.5mu]}}} % intervalles d'entiers (glyphe ⟦ absent de la police)
% Glyphes absents de Fira Math : redéfinis à partir de glyphes disponibles
\AtBeginDocument{%
  \def\setminus{\mathbin{\backslash}}\let\smallsetminus\setminus
  \def\vdots{\mathord{\vbox{\baselineskip3.2pt\lineskiplimit0pt\kern4pt\hbox{.}\hbox{.}\hbox{.}}}}%
  \def\ddots{\mathinner{\mkern1mu\raise6.5pt\hbox{.}\mkern2mu\raise3.6pt\hbox{.}\mkern2mu\raise0.7pt\hbox{.}\mkern1mu}}%
  \def\triangle{\mathord{\scalebox{0.9}{$\symup{\Delta}$}}}%
  \def\nabla{\mathord{\raisebox{\depth}{\scalebox{1}[-1]{$\symup{\Delta}$}}}}%
  \def\checkmark{\popcheck{popdark}}%
  \def\star{\mathbin{\ast}}\def\bigstar{\mathord{\ast}}%
  \def\diamond{\mathbin{\scalebox{0.6}{$\Diamond$}}}%
  \def\bigcup{\mathop{\mathchoice{\raisebox{-0.3em}{\scalebox{1.7}{$\cup$}}}{\scalebox{1.25}{$\cup$}}{\cup}{\cup}}\displaylimits}%
  \def\bigcap{\mathop{\mathchoice{\raisebox{-0.3em}{\scalebox{1.7}{$\cap$}}}{\scalebox{1.25}{$\cap$}}{\cap}{\cap}}\displaylimits}%
  \def\lesseqqgtr{\mathrel{\lessgtr}}\def\bigoplus{\mathop{\oplus}\displaylimits}%
}
\providecommand{\ssi}{si et seulement si} % abréviation fréquente
\AtBeginDocument{\providecommand{\wideparen}{\overparen}} % arc de cercle

%% FIGFIX-BEGIN v5 — correctifs globaux des figures (docs/onbuch-generation/tools/figfix)
\usepackage{adjustbox}\usepackage{etoolbox}\tcbuselibrary{raster}
% toute figure TikZ/pgfplots plus large que la ligne est réduite à la largeur disponible
\BeforeBeginEnvironment{tikzpicture}{\begin{adjustbox}{max width=\linewidth}}
\AfterEndEnvironment{tikzpicture}{\end{adjustbox}}
% légendes pgfplots lisibles par-dessus les courbes
\pgfplotsset{every axis/.append style={legend style={fill=white,fill opacity=0.92,text opacity=1,draw=black!15,inner sep=3pt}}}
% étiquettes d'arêtes (syntaxe edge["..."]) avec fond blanc
\tikzset{every edge quotes/.append style={fill=white,inner sep=1.2pt}}
%% FIGFIX-END
\begin{document}

\pagegarde{Lesson 38 — Grammar: Revision: Direct and indirect speech}{Anglais}{Tle A}{Chapter 4 : Citizenship and Human Rights}{EN38}{2 h}{Cette leçon de révision systématise le passage du discours direct au discours indirect : transformations des temps, des pronoms, des marqueurs spatio-temporels et des verbes introducteurs. Elle entraîne l'élève à rapporter déclarations, questions, ordres et demandes dans le contexte du chapitre « Citizenship and Human Rights ».
\vspace{4pt}
\begin{multicols}{2}\small
\begin{itemize}
\item À la fin de cette leçon, je sais distinguer le discours direct du discours indirect et identifier le verbe introducteur (reporting verb).
\item À la fin de cette leçon, je sais appliquer le recul des temps (backshift) : present simple → past simple, present perfect → past perfect, will → would, can → could.
\item À la fin de cette leçon, je sais transformer les pronoms, possessifs et marqueurs spatio-temporels (now → then, today → that day, here → there).
\item À la fin de cette leçon, je sais rapporter une question fermée avec if/whether et une question ouverte avec le mot interrogatif suivi de l'ordre sujet-verbe.
\item À la fin de cette leçon, je sais rapporter un ordre, une demande ou un conseil avec tell/ask/advise + complément + to-infinitive.
\item À la fin de cette leçon, je sais choisir le verbe introducteur adapté au sens : say, tell, ask, promise, warn, suggest, admit, deny.
\end{itemize}\end{multicols}}

\begin{sommaire}
\begin{multicols}{2}
\begin{enumerate}
\item Observation et découverte
\item Le recul des temps (backshift)
\item Pronoms et marqueurs spatio-temporels
\item Rapporter questions, ordres et demandes
\item Verbes introducteurs et comparaison avec le français
\item Entraînement guidé et pièges récapitulatifs
\item Activité d'intégration
\item Méthodes et exercices résolus
\item Exercices
\item Corrigés des exercices
\item Fiche bilan
\end{enumerate}
\end{multicols}
\end{sommaire}

\begin{objectifs}
\begin{itemize}
\item À la fin de cette leçon, je sais distinguer le discours direct du discours indirect et identifier le verbe introducteur (reporting verb).
\item À la fin de cette leçon, je sais appliquer le recul des temps (backshift) : present simple → past simple, present perfect → past perfect, will → would, can → could.
\item À la fin de cette leçon, je sais transformer les pronoms, possessifs et marqueurs spatio-temporels (now → then, today → that day, here → there).
\item À la fin de cette leçon, je sais rapporter une question fermée avec if/whether et une question ouverte avec le mot interrogatif suivi de l'ordre sujet-verbe.
\item À la fin de cette leçon, je sais rapporter un ordre, une demande ou un conseil avec tell/ask/advise + complément + to-infinitive.
\item À la fin de cette leçon, je sais choisir le verbe introducteur adapté au sens : say, tell, ask, promise, warn, suggest, admit, deny.
\item À la fin de cette leçon, je sais reconnaître les cas où le temps ne change pas (vérités générales, faits encore vrais, reporting verb au présent).
\item À la fin de cette leçon, je sais éviter les erreurs typiques des francophones (inversion maintenue dans les questions rapportées, oubli du backshift, confusion say/tell).
\item À la fin de cette leçon, je sais rapporter un court dialogue ou une déclaration de presse en un paragraphe cohérent au discours indirect.
\end{itemize}
\end{objectifs}

\begin{prerequis}
\begin{itemize}
\item Les temps principaux de l'indicatif anglais (present simple, present continuous, past simple, present perfect, future avec will)
\item Les pronoms personnels, possessifs et démonstratifs
\item La structure de la phrase interrogative directe (inversion sujet-auxiliaire)
\item Les modaux can, must, may et leurs équivalents au passé
\item La négation et la place des auxiliaires (notions de Première)
\end{itemize}
\end{prerequis}

\newpage

\coursec{Observation et découverte}

Pendant l'Assemblée générale des élèves du lycée de Bafoussam, le délégué annonce : \eng{The principal said that the inter-class debate on human rights would take place on Friday.} Un élève distrait lève la main : « Qu'a-t-il dit exactement ? » Pour répondre, il faut savoir rapporter fidèlement les paroles de quelqu'un : c'est tout l'art du \cle{discours indirect}. Au Cameroun anglophone comme au Nigeria ou au Ghana, journalistes, avocats et témoins utilisent chaque jour le \cle{reported speech} pour relater déclarations, promesses et accusations. Maîtriser cette structure, c'est pouvoir rapporter sans déformer — une compétence de citoyen et un classique du Baccalauréat.

\textbf{Problématique :} comment passer des paroles exactes prononcées par quelqu'un (\eng{direct speech}) à un compte rendu fidèle de ces paroles (\eng{indirect} ou \eng{reported speech}) ? Quels changements faut-il opérer, et pourquoi ?

\courssub{Lire et observer : un article de presse}

\begin{textebox}[Cameroon Tribune — “Minister Addresses Human Rights Day”]
Yaoundé, December 10. The Minister of Human Rights addressed journalists yesterday at a press conference in the Hilton Hotel. “We will protect every citizen,” the Minister said. “I am proud of our young voters, and we are working hard to guarantee free elections.”

The Minister said that they would protect every citizen. He added that he was proud of the young voters of the country and that they were working hard to guarantee free elections.

A journalist from the BBC then asked the Minister whether the government would invite international observers. The Minister replied that invitations had already been sent. Another reporter asked him when the new electoral law would be published. He answered that it would be ready the following month.

Finally, a human rights activist told the Minister to publish the names of all detainees. The Minister asked the activist to submit a written request to his office. “Justice is the foundation of peace,” he concluded. The audience applauded, and the journalists hurried to report his words in the evening news.
\end{textebox}

\textbf{Glossaire :}

\begin{center}
\begin{tabularx}{\linewidth}{|l|l|X|}
\hline
\rowcolor{popblueL} English & Nature & Français \\
\hline
to address journalists & verbe & s'adresser aux journalistes \\
\hline
a press conference & nom & une conférence de presse \\
\hline
a voter & nom & un électeur \\
\hline
whether & conjonction & si (interrogation indirecte) \\
\hline
a detainee & nom & un détenu \\
\hline
to submit a request & verbe & soumettre une requête \\
\hline
the foundation & nom & le fondement, la base \\
\hline
to applaud & verbe & applaudir \\
\hline
an observer & nom & un observateur \\
\hline
to guarantee & verbe & garantir \\
\hline
a polling station & nom & un bureau de vote \\
\hline
a ballot box & nom & une urne \\
\hline
\end{tabularx}
\end{center}

\textbf{Questions de compréhension (oral) :}
\begin{enumerate}
\item Where did the press conference take place?
\item What did the Minister promise?
\item What did the BBC journalist ask?
\item What did the activist tell the Minister to do?
\end{enumerate}

\courssub{Observer les deux façons de rapporter les paroles}

Reprenons les deux premiers paragraphes de l'article. Le premier cite les paroles exactes du ministre ; le second les rapporte.

\begin{exemplebox}[Direct speech : les paroles exactes]
\eng{“We will protect every citizen,” the Minister said.}

\eng{“I am proud of our young voters, and we are working hard to guarantee free elections.”}

On observe : des \cle{guillemets} (“ … ”), un \cle{verbe introducteur} (\eng{said}), les temps du moment de la parole (\eng{will protect}, \eng{am}, \eng{are working}), et les pronoms du locuteur (\eng{we}, \eng{I}, \eng{our}).
\end{exemplebox}

\begin{exemplebox}[Indirect speech : les paroles rapportées]
\eng{The Minister said (that) they would protect every citizen.}

\eng{He added (that) he was proud of the young voters and (that) they were working hard to guarantee free elections.}

On observe : plus de guillemets ; apparition du mot \cle{that} ; recul des temps (\eng{will} $\to$ \eng{would}, \eng{am} $\to$ \eng{was}, \eng{are working} $\to$ \eng{were working}) ; changement des pronoms (\eng{we} $\to$ \eng{they}, \eng{I} $\to$ \eng{he}, \eng{our} $\to$ \eng{their/the}).
\end{exemplebox}

\begin{definition}[Direct speech et indirect speech]
Le \cle{discours direct} (\eng{direct speech}) reproduit les paroles exactes de quelqu'un, entre guillemets, telles qu'elles ont été prononcées : \eng{“I know my rights,” she said.}

Le \cle{discours indirect} (\eng{indirect} ou \eng{reported speech}) rapporte les paroles de quelqu'un \textbf{sans les citer mot à mot} et \textbf{sans guillemets} ; il entraîne des changements de \textbf{temps verbaux}, de \textbf{pronoms} et de \textbf{repères spatio-temporels} : \eng{She said (that) she knew her rights.}
\end{definition}

\begin{exemplebox}[Un exemple traduit, mot à mot]
\eng{“I know my rights,” she said.} (« Je connais mes droits », a-t-elle dit.)

$\to$ \eng{She said (that) she knew her rights.} (Elle a dit qu'elle connaissait ses droits.)

Remarquez que le français opère la même transformation : \emph{je connais} $\to$ \emph{elle connaissait}. Le mécanisme est donc familier ; seules les formes anglaises doivent être apprises.
\end{exemplebox}

\courssub{Le mot that : facultatif après say et tell}

\begin{propriete}[La conjonction that]
Après les verbes introducteurs \eng{say} et \eng{tell}, la conjonction \cle{that} est \textbf{facultative} :

\eng{The Minister said (that) they would protect every citizen.} = \eng{The Minister said they would protect every citizen.}

\begin{itemize}
\item \eng{say something (to somebody)} : \eng{He said (to me) that he was tired.} (Il m'a dit qu'il était fatigué.)
\item \eng{tell somebody something} : \eng{He told me that he was tired.} — attention, \eng{tell} exige un \textbf{complément d'objet indirect} (la personne à qui l'on parle).
\end{itemize}

Dans la presse écrite et à l'examen, on garde souvent \eng{that} pour la clarté ; à l'oral, on l'omet fréquemment.
\end{propriete}

\begin{attention}[say ou tell ? L'erreur classique du francophone]
\begin{itemize}
\item \eng{He said me that...} est \textbf{faux}. On dit \eng{He said that...} ou \eng{He told me that...}
\item \eng{He told that he was tired.} est \textbf{faux} : \eng{tell} doit être suivi d'une personne : \eng{He told \textbf{me} that he was tired.}
\item Ne confondez pas \eng{say} (dire quelque chose) et \eng{tell} (dire à quelqu'un). Prononciation : \eng{said} se dit « sèd » (/sɛd/) et non « saïd » !
\end{itemize}
\end{attention}

\courssub{Trois types d'énoncés à rapporter}

Dans notre article, le journaliste n'a pas rapporté que des déclarations. Observons :

\begin{center}
\begin{tabularx}{\linewidth}{|X|X|X|}
\hline
\rowcolor{popblueL} Type d'énoncé & Direct speech & Reported speech \\
\hline
Déclaration (\eng{statement}) & \eng{“We will protect every citizen.”} & \eng{He said (that) they would protect every citizen.} \\
\hline
Question fermée (oui/non) & \eng{“Will you invite observers?”} & \eng{A journalist asked \textbf{whether/if} the government would invite observers.} \\
\hline
Question ouverte (\eng{wh-}) & \eng{“When will the law be published?”} & \eng{A reporter asked him \textbf{when} the law would be published.} \\
\hline
Ordre (\eng{command}) & \eng{“Publish the names!”} & \eng{She told the Minister \textbf{to publish} the names.} \\
\hline
Demande (\eng{request}) & \eng{“Please submit a request.”} & \eng{He asked the activist \textbf{to submit} a written request.} \\
\hline
\end{tabularx}
\end{center}

\begin{propriete}[Les trois structures du discours rapporté]
\begin{enumerate}
\item \textbf{Déclarations} : \eng{say / tell + (that) + phrase} — l'ordre des mots ne change pas (sujet + verbe).
\item \textbf{Questions} : \eng{ask + if/whether} (question fermée) ou \eng{ask + mot interrogatif} (\eng{when, where, why, how...}) + \textbf{ordre sujet–verbe} (on supprime l'inversion et l'auxiliaire \eng{do/does/did}).
\item \textbf{Ordres et demandes} : \eng{tell / ask + complément de personne + to + base verbale} ; à la forme négative : \eng{not to} (\eng{He told me \textbf{not to} be late.}).
\end{enumerate}
\end{propriete}

\begin{popfigure}
\begin{center}\tcbox[colback=popblue,colframe=popblue,coltext=white,arc=9pt,boxsep=4pt,left=6pt,right=6pt]{\bfseries\small Reported speech}\end{center}
\vspace{-2pt}
\begin{tcbraster}[raster columns=2,raster equal height=rows,raster column skip=6pt,raster row skip=6pt,enhanced,blankest]
\begin{tcolorbox}[enhanced,colback=white,colframe=popgreen,boxrule=0.9pt,arc=3pt,title={Statements},fonttitle=\bfseries\scriptsize,coltitle=white,colbacktitle=popgreen,left=4pt,right=4pt,top=2pt,bottom=2pt,boxsep=1pt,toptitle=1.5pt,bottomtitle=1.5pt,halign=left]
\begin{itemize}[nosep,leftmargin=*,itemsep=1pt,font=\scriptsize]
\item say / tell + (that)
\end{itemize}
\end{tcolorbox}
\begin{tcolorbox}[enhanced,colback=white,colframe=poporange,boxrule=0.9pt,arc=3pt,title={Questions},fonttitle=\bfseries\scriptsize,coltitle=white,colbacktitle=poporange,left=4pt,right=4pt,top=2pt,bottom=2pt,boxsep=1pt,toptitle=1.5pt,bottomtitle=1.5pt,halign=left]
\begin{itemize}[nosep,leftmargin=*,itemsep=1pt,font=\scriptsize]
\item ask + if / whether
\item ask + wh- word
\end{itemize}
\end{tcolorbox}
\begin{tcolorbox}[enhanced,colback=white,colframe=poppurple,boxrule=0.9pt,arc=3pt,title={Commands / Requests},fonttitle=\bfseries\scriptsize,coltitle=white,colbacktitle=poppurple,left=4pt,right=4pt,top=2pt,bottom=2pt,boxsep=1pt,toptitle=1.5pt,bottomtitle=1.5pt,halign=left]
\begin{itemize}[nosep,leftmargin=*,itemsep=1pt,font=\scriptsize]
\item tell / ask + sb + to-inf.
\end{itemize}
\end{tcolorbox}
\end{tcbraster}

\legende{Carte mentale : les trois types d'énoncés rapportés et leurs verbes introducteurs.}
\end{popfigure}

\courssub{Le mécanisme général de la transformation}

Comment passe-t-on du discours direct au discours indirect ? Trois opérations se combinent, que nous détaillerons dans les sections suivantes :

\begin{enumerate}
\item un \textbf{verbe introducteur} au passé (\eng{said, told, asked, added, replied...}) ;
\item le \textbf{recul des temps} (\eng{backshift}) : \eng{will} $\to$ \eng{would}, \eng{am} $\to$ \eng{was}, etc. ;
\item l'\textbf{adaptation des pronoms et des marqueurs spatio-temporels} : \eng{I} $\to$ \eng{he/she}, \eng{now} $\to$ \eng{then}, \eng{today} $\to$ \eng{that day}, \eng{next year} $\to$ \eng{the following year}.
\end{enumerate}

\begin{popfigure}
\begin{tikzpicture}[node distance=0.6cm]
\node[draw=popgreen, fill=popgreenL, rounded corners, align=center, text width=3.6cm, font=\small] (direct) {\textbf{Direct speech}\\ “I am proud of our\\ young voters.”};
\node[draw=popblue, fill=popblueL, rounded corners, align=center, text width=4.6cm, font=\small, right=1.4cm of direct] (boite) {\textbf{Transformation}\\ verbe introducteur\\ + backshift\\ + pronoms et marqueurs};
\node[draw=poppurple, fill=poppurpleL, rounded corners, align=center, text width=3.9cm, font=\small, right=1.4cm of boite] (indirect) {\textbf{Reported speech}\\ He said (that) he was\\ proud of their young voters.};
\draw[-{Stealth[length=3mm]}, line width=1.2pt, popdark] (direct) -- (boite);
\draw[-{Stealth[length=3mm]}, line width=1.2pt, popdark] (boite) -- (indirect);
\end{tikzpicture}
\legende{Schéma de la transformation : du discours direct au discours rapporté.}
\end{popfigure}

\begin{textebox}[Direct $\to$ reported : un exemple complet à retenir]
\eng{“I am proud of our young voters,” said the mayor. “We will organise free elections next year.”}

$\to$ \eng{The mayor said (that) he was proud of their young voters and added that they would organise free elections the following year.}

Observez chaque changement : \eng{I} $\to$ \eng{he} ; \eng{am} $\to$ \eng{was} ; \eng{our} $\to$ \eng{their} ; \eng{we} $\to$ \eng{they} ; \eng{will} $\to$ \eng{would} ; \eng{next year} $\to$ \eng{the following year}. Le verbe \eng{added} permet d'enchaîner élégamment la deuxième déclaration.
\end{textebox}

\begin{savaistu}[Le discours indirect, langue de la presse]
Dans la presse anglophone (\eng{BBC}, \eng{The Guardian}, \eng{Cameroon Tribune}), le discours indirect est \textbf{la norme} pour rapporter les déclarations officielles : il permet au journaliste de condenser un long discours en quelques lignes. Les guillemets ne servent que pour les citations \textbf{mot à mot}, courtes et percutantes, censées frapper le lecteur. C'est pourquoi un même article alterne souvent les deux : une phrase entre guillemets pour l'émotion, du reported speech pour l'information. Au procès, le témoin qui rapporte des propos utilise exactement la même grammaire : \eng{He told the court that the accused had threatened him.}
\end{savaistu}

\begin{exoresolu}[Repérer direct et indirect dans l'article]
Énoncé — Dans les phrases suivantes tirées de l'article, indiquez s'il s'agit de discours direct (D) ou rapporté (R), puis soulignez le verbe introducteur.

a) \eng{“We will protect every citizen,” the Minister said.}

b) \eng{He added that he was proud of the young voters.}

c) \eng{A journalist asked the Minister whether the government would invite observers.}

d) \eng{“Justice is the foundation of peace,” he concluded.}

\tcblower
Solution détaillée :

a) \textbf{D} — guillemets + verbe introducteur \eng{said} : paroles exactes citées.

b) \textbf{R} — pas de guillemets ; verbe introducteur \eng{added} + \eng{that} ; backshift (\eng{am} $\to$ \eng{was}).

c) \textbf{R} — question rapportée introduite par \eng{asked} + \eng{whether} ; ordre sujet–verbe rétabli (\eng{the government would invite}).

d) \textbf{D} — citation mot à mot entre guillemets ; verbe introducteur \eng{concluded}.

Méthode : cherchez d'abord les \textbf{guillemets} (présents = direct ; absents = rapporté), puis identifiez le \textbf{verbe introducteur} et, au discours rapporté, vérifiez le \textbf{recul du temps}.
\end{exoresolu}

\courssub{Le recul des temps (Backshift) : règle complète et exceptions}

C'est le cœur mécanique de la transformation. Lorsque le verbe introducteur est au \textbf{passé} (\eng{said, told, asked, wondered...}), les temps de la proposition rapportée reculent d'un cran vers le passé.

\begin{propriete}[Tableau du recul des temps (Backshift)]
\begin{center}
\begin{tabularx}{\linewidth}{|X|X|X|}
\hline
\rowcolor{popblueL} Direct Speech (Temps original) & Reported Speech (Après backshift) & Exemple \\
\hline
Present Simple (\eng{I vote}) & Past Simple (\eng{he voted}) & \eng{“I vote in Bamenda.”} $\to$ \eng{He said he voted in Bamenda.} \\
\hline
Present Continuous (\eng{I am voting}) & Past Continuous (\eng{he was voting}) & \eng{“I am voting now.”} $\to$ \eng{He said he was voting then.} \\
\hline
Present Perfect (\eng{I have voted}) & Past Perfect (\eng{he had voted}) & \eng{“I have voted.”} $\to$ \eng{She said she had voted.} \\
\hline
Present Perfect Continuous & Past Perfect Continuous & \eng{“I have been waiting.”} $\to$ \eng{He said he had been waiting.} \\
\hline
Past Simple (\eng{I voted}) & Past Perfect (\eng{he had voted}) & \eng{“I voted yesterday.”} $\to$ \eng{He said he had voted the day before.} \\
\hline
Past Continuous (\eng{I was voting}) & Past Perfect Continuous (\eng{he had been voting}) & \eng{“I was voting.”} $\to$ \eng{She said she had been voting.} \\
\hline
\eng{will} (\eng{I will vote}) & \eng{would} (\eng{he would vote}) & \eng{“I will vote.”} $\to$ \eng{He said he would vote.} \\
\hline
\eng{can} (\eng{I can vote}) & \eng{could} (\eng{he could vote}) & \eng{“I can vote.”} $\to$ \eng{She said she could vote.} \\
\hline
\eng{may} (\eng{It may rain}) & \eng{might} (\eng{it might rain}) & \eng{“It may rain.”} $\to$ \eng{He said it might rain.} \\
\hline
\eng{must} (\eng{I must go}) & \eng{had to} (\eng{he had to go}) & \eng{“I must go.”} $\to$ \eng{He said he had to go.} \\
\hline
\eng{shall} (\eng{Shall I go?}) & \eng{should} (\eng{he should go}) & \eng{“Shall I go?”} $\to$ \eng{He asked if he should go.} \\
\hline
\end{tabularx}
\end{center}

\textbf{Modaux qui ne changent pas :} \eng{would, could, might, should, ought to, used to}.
\eng{“I would help if I could.”} $\to$ \eng{He said he would help if he could.}
\end{propriete}

\begin{attention}[Exceptions au backshift : quand le temps ne recule PAS]
Le backshift n'est \textbf{pas obligatoire} (mais reste correct) dans trois cas :
\begin{enumerate}
\item \textbf{Vérité générale / fait scientifique immuable :} \eng{The teacher said (that) the Earth \textbf{revolves} around the Sun.} (Pas \eng{revolved}).
\item \textbf{Report immédiat / situation encore actuelle :} \eng{“I am hungry,” she says.} $\to$ \eng{She \textbf{says} she \textbf{is} hungry.} (Verbe introducteur au présent) ; ou \eng{She \textbf{said} she \textbf{is} hungry.} (Si on rapporte tout de suite et qu'elle a encore faim).
\item \textbf{Verbe au conditionnel / modaux invariables :} voir tableau ci-dessus (\eng{would, could, might, should}).
\end{enumerate}
\textbf{Conseil Bac :} En cas de doute à l'écrit, \textbf{appliquez le backshift} : c'est toujours correct grammaticalement et attendu par les correcteurs.
\end{attention}

\begin{exemplebox}[Backshift et aspect progressif / perfectif]
\eng{“We are building a new school,” the mayor said.}
$\to$ \eng{The mayor said (that) they \textbf{were building} a new school.}

\eng{“The road has been repaired,” the engineer said.}
$\to$ \eng{The engineer said (that) the road \textbf{had been repaired}.} (Passif : \eng{has been} $\to$ \eng{had been}).

\eng{“I will have finished by 5 pm,” the student said.}
$\to$ \eng{The student said (that) he \textbf{would have finished} by 5 pm.}
\end{exemplebox}

\courssub{Pronoms et marqueurs spatio-temporels : le décalage de perspective}

Le discours indirect déplace le point de vue : le « je » du locuteur devient « il/elle » du rapporteur ; le « ici » devient « là-bas » ; le « maintenant » devient « alors ».

\begin{propriete}[Tableau de conversion des pronoms et déterminants]
\begin{center}
\begin{tabularx}{\linewidth}{|X|X|X|}
\hline
\rowcolor{popblueL} Direct Speech & Reported Speech & Remarque \\
\hline
\eng{I, me, my, mine} & \eng{he/she, him/her, his/her, his/hers} & Selon le genre du locuteur \\
\hline
\eng{we, us, our, ours} & \eng{they, them, their, theirs} & \\
\hline
\eng{you, your, yours} & \eng{I/we, my/our, mine/ours} OU \eng{he/she/they...} & Dépend de qui rapporte à qui \\
\hline
\eng{this, these} & \eng{that, those} & Proximité $\to$ distance \\
\hline
\end{tabularx}
\end{center}
\end{propriete}

\begin{propriete}[Tableau de conversion des marqueurs spatio-temporels]
\begin{center}
\begin{tabularx}{\linewidth}{|X|X|}
\hline
\rowcolor{popblueL} Direct Speech & Reported Speech \\
\hline
\eng{now} & \eng{then} / \eng{at that moment} \\
\hline
\eng{today} & \eng{that day} \\
\hline
\eng{tonight} & \eng{that night} \\
\hline
\eng{yesterday} & \eng{the day before} / \eng{the previous day} \\
\hline
\eng{tomorrow} & \eng{the next day} / \eng{the following day} \\
\hline
\eng{last week / month / year} & \eng{the week / month / year before} / \eng{the previous week...} \\
\hline
\eng{next week / month / year} & \eng{the following week / month / year} \\
\hline
\eng{ago} & \eng{before} \\
\hline
\eng{here} & \eng{there} \\
\hline
\eng{this (adjectif)} & \eng{that} / \eng{the} \\
\hline
\end{tabularx}
\end{center}
\end{propriete}

\begin{attention}[Pièges fréquents : \eng{this/that} et \eng{here/there}]
\begin{itemize}
\item \eng{“I like \textbf{this} book.”} $\to$ \eng{He said he liked \textbf{that} book.} (Pas \eng{this}).
\item \eng{“Come \textbf{here}.”} $\to$ \eng{He told me to go \textbf{there}.} (Le verbe de mouvement change souvent : \eng{come} $\to$ \eng{go}).
\item \eng{“I’ll do it \textbf{now}.”} $\to$ \eng{He said he would do it \textbf{then}.} (Jamais \eng{now} au discours rapporté passé).
\end{itemize}
\end{attention}

\begin{exemplebox}[Application combinée : pronoms + temps + marqueurs]
\eng{“I will submit \textbf{my} report \textbf{tomorrow},” the official said.}
$\to$ \eng{The official said (that) \textbf{he} would submit \textbf{his} report \textbf{the following day}.}

\eng{“We have been waiting \textbf{here} for two hours,” the voters complained.}
$\to$ \eng{The voters complained (that) \textbf{they} had been waiting \textbf{there} for two hours.}
\end{exemplebox}

\courssub{Rapporter les questions, les ordres et les demandes : cas particuliers}

\courssub{Les questions (Interrogation indirecte)}

La règle d'or : \textbf{ordre affirmatif (sujet + verbe)}. On supprime l'inversion et l'auxiliaire \eng{do/does/did}.

\begin{propriete}[Questions fermées (Yes/No) : \eng{if} ou \eng{whether}]
\eng{“Do you know your rights?”} $\to$ \eng{He asked \textbf{if / whether} I knew my rights.}
\eng{“Will you come?”} $\to$ \eng{She asked \textbf{whether} I would come.}
\begin{itemize}
\item \eng{whether} est plus formel, souvent préféré après préposition (\eng{asked about whether...}) ou avant infinitif (\eng{asked whether to go}).
\item \eng{if} est plus courant à l'oral.
\end{itemize}
\end{propriete}

\begin{propriete}[Questions ouvertes (Wh-) : on garde le mot interrogatif]
\eng{“Where is the polling station?”} $\to$ \eng{He asked \textbf{where} the polling station \textbf{was}.}
\eng{“Why did you vote?”} $\to$ \eng{She asked \textbf{why} I \textbf{had voted}.}
\eng{“How can we help?”} $\to$ \eng{They asked \textbf{how} they \textbf{could} help.}
\begin{itemize}
\item Pas d'auxiliaire \eng{do/does/did} : \eng{“Where \textbf{does} he live?”} $\to$ \eng{She asked where he \textbf{lived}.} (Pas \eng{did live}).
\item Si le mot interrogatif est le \textbf{sujet} : pas de changement d'ordre. \eng{“Who called?”} $\to$ \eng{He asked who \textbf{called}.} (Pas \eng{who did call}).
\end{itemize}
\end{propriete}

\begin{attention}[Erreur majeure : conserver l'inversion ou l'auxiliaire \eng{do}]
\begin{itemize}
\item \textbf{Faux :} \eng{He asked where was the station.} / \eng{He asked where did I vote.}
\item \textbf{Juste :} \eng{He asked where the station was.} / \eng{He asked where I voted.}
\end{itemize}
\end{attention}

\courssub{Les ordres et les demandes (Infinitif)}

\begin{propriete}[Structure : \eng{tell/ask/order/warn/advise/beg + personne + (not) to + BV}]
\begin{center}
\begin{tabularx}{\linewidth}{|X|X|}
\hline
\rowcolor{popblueL} Direct (Impératif) & Reported (Infinitif) \\
\hline
\eng{“Sit down!”} & \eng{He told me \textbf{to sit} down.} \\
\hline
\eng{“Don't move!”} & \eng{He told me \textbf{not to move}.} \\
\hline
\eng{“Please, help me.”} & \eng{She asked me \textbf{to help} her.} \\
\hline
\eng{“Let's go!”} & \eng{He suggested \textbf{going} / \textbf{that we should go}.} \\
\hline
\end{tabularx}
\end{center}
\end{propriete}

\begin{attention}[Verbes introducteurs spécifiques pour ordres/demandes]
\begin{itemize}
\item \eng{tell} = ordre fort (\eng{He told me to leave.}).
\item \eng{ask} = demande polie (\eng{He asked me to help.}).
\item \eng{order, command} = ordre officiel/militaire.
\item \eng{warn} = avertir (\eng{He warned me not to go.}).
\item \eng{advise} = conseiller (\eng{She advised me to study.}).
\item \eng{beg, implore} = supplier.
\item \eng{remind} = rappeler (\eng{He reminded me to vote.}).
\end{itemize}
\textbf{Attention :} \eng{suggest} et \eng{recommend} ne s'utilisent \textbf{PAS} avec \eng{to}-infinitif.
\eng{“Let's vote,” he said.} $\to$ \eng{He suggested \textbf{voting} / \textbf{that we (should) vote}.} (Jamais \eng{suggested us to vote}).
\end{attention}

\courssub{Verbes introducteurs : classification et comparaison avec le français}

Au-delà de \eng{say, tell, ask}, la richesse de l'anglais réside dans le choix du verbe introducteur qui nuance l'intention (promettre, menacer, nier, proposer...). Chaque verbe impose une structure grammaticale précise.

\begin{propriete}[Classification des verbes introducteurs par structure]
\begin{center}
\begin{tabularx}{\linewidth}{|l|X|X|}
\hline
\rowcolor{popblueL} Structure & Verbes principaux & Exemple (Direct $\to$ Reported) \\
\hline
\eng{V + (that) + clause} & \eng{say, tell, state, declare, claim, admit, deny, explain, complain, promise, threaten, warn, agree, refuse, decide, hope, expect} & \eng{“I will pay,” he said.} $\to$ \eng{He \textbf{promised} (that) he would pay.} \\
\hline
\eng{V + obj + (that) + clause} & \eng{tell, inform, convince, assure, remind, notify, persuade} & \eng{“The meeting is at 5,” she told me.} $\to$ \eng{She \textbf{informed} me (that) the meeting was at 5.} \\
\hline
\eng{V + to-inf} & \eng{promise, agree, refuse, offer, threaten, decide, hope, claim, swear} & \eng{“I’ll help you,” he said.} $\to$ \eng{He \textbf{offered to help} me.} \\
\hline
\eng{V + obj + to-inf} & \eng{ask, tell, order, advise, warn, beg, invite, encourage, remind, force, allow, enable, persuade} & \eng{“Please stay,” she said.} $\to$ \eng{She \textbf{asked} me \textbf{to stay}.} \\
\hline
\eng{V + -ing} & \eng{suggest, recommend, propose, admit, deny, regret, imagine, consider, risk, avoid} & \eng{“Let's strike,” they said.} $\to$ \eng{They \textbf{suggested striking} / \textbf{that they (should) strike}.} \\
\hline
\eng{V + prep + -ing} & \eng{apologize for, insist on, accuse of, blame for, congratulate on, complain about, object to} & \eng{“Sorry I'm late,” he said.} $\to$ \eng{He \textbf{apologized for being} late.} \\
\hline
\end{tabularx}
\end{center}
\end{propriete}

\begin{methode}[Comment choisir le bon verbe et la bonne structure ?]
\begin{enumerate}
\item Identifier l'\textbf{intention} du locuteur (promettre ? ordonner ? s'excuser ? nier ?).
\item Choisir le verbe anglais correspondant (\eng{promise, order, apologize, deny...}).
\item Appliquer \textbf{la structure imposée par ce verbe} (voir tableau ci-dessus).
\item Effectuer le \textbf{backshift} et les changements de pronoms/marqueurs si le verbe introducteur est au passé.
\end{enumerate}
\end{methode}

\begin{attention}[Faux amis et pièges de traduction (Français $\to$ Anglais)]
\begin{center}
\begin{tabularx}{\linewidth}{|X|X|X|}
\hline
\rowcolor{popblueL} Français & Erreur fréquente (Traduction littérale) & Anglais correct \\
\hline
Il a \textbf{dit} qu'il viendrait. & \eng{He said he would come.} & \eng{He said he would come.} (OK) \\
Il m'a \textbf{dit} qu'il viendrait. & \eng{He said me he would come.} & \eng{He \textbf{told} me he would come.} \\
Il a \textbf{demandé} de partir. & \eng{He asked to leave.} (Sujet = lui) & \eng{He asked \textbf{me} to leave.} (Si c'est à moi) \\
Il a \textbf{demandé} si je venais. & \eng{He asked if I come.} & \eng{He asked \textbf{whether/if} I \textbf{came}.} \\
Il a \textbf{suggéré} d'aller voter. & \eng{He suggested to go vote.} & \eng{He suggested \textbf{going} to vote / \textbf{that we (should) go} vote.} \\
Il a \textbf{prétendu} être innocent. & \eng{He pretended to be innocent.} & \eng{He \textbf{claimed} to be innocent.} (\eng{pretend} = faire semblant) \\
Il a \textbf{nié} avoir volé. & \eng{He denied to have stolen.} & \eng{He denied \textbf{stealing} / \textbf{having stolen}.} \\
\hline
\end{tabularx}
\end{center}
\end{attention}

\begin{savaistu}[Subjonctif français vs Infinitif/\eng{that}-clause anglais]
Le français utilise massivement le subjonctif après les verbes de volonté, doute, émotion (\emph{il veut que je parte, il faut que tu saches}). L'anglais utilise :
\begin{itemize}
\item L'\textbf{infinitif} (\eng{He wants me \textbf{to leave}.} / \eng{It is necessary \textbf{to know}.})
\item La construction \eng{that + should + BV} (formel) ou \eng{that + présent/passé} (\eng{He insists that I \textbf{leave / should leave}.}).
\end{itemize}
Ne traduisez jamais \emph{que} + subjonctif par \eng{that} + subjonctif anglais (qui n'existe pas sous cette forme).
\end{savaistu}

\courssub{Entraînement guidé et pièges récapitulatifs}

\begin{exercice}[Transformation complète : Déclarations et Questions]
Transformez les phrases suivantes au discours rapporté. Le verbe introducteur est au passé.

1. \eng{“The new law protects minorities,” the deputy declared.}
2. \eng{“We have been campaigning for weeks,” the activists said.}
3. \eng{“Did you register to vote?” the officer asked the young man.}
4. \eng{“Why hasn't the report been published?” the journalist asked.}
5. \eng{“I cannot guarantee the result,” the candidate admitted.}
6. \eng{“Where did you put the ballot papers?” the supervisor asked.}
\end{exercice}

\begin{exercice}[Ordres, demandes et verbes introducteurs variés]
Rapportez les énoncés en choisissant le verbe introducteur le plus approprié parmi la liste : \eng{advise, warn, order, beg, suggest, remind, promise, threaten, deny, apologize for}.

1. \eng{“Don't forget to bring your ID card,” the monitor said to the students.}
2. \eng{“I will certainly resign if I lose,” the mayor said.}
3. \eng{“You should see a doctor about that cough,” the nurse said to me.}
4. \eng{“If you strike, I will fire you all,” the manager said to the workers.}
5. \eng{“Let's organize a peaceful march,” the leader said.}
6. \eng{“I didn't take the money,” the suspect said.}
7. \eng{“Please, please, don't arrest my brother,” the woman said to the police.}
8. \eng{“Stand up straight!” the sergeant said to the recruits.}
9. \eng{“Don't go near the border zone,” the guide said to the tourists.}
10. \eng{“I'm sorry I interrupted you,” the student said to the teacher.}
\end{exercice}

\begin{exercice}[Correction d'erreurs : Chasse aux fautes francophones]
Chaque phrase contient une erreur (grammaire, ordre des mots, choix du verbe, backshift). Corrigez-la.

1. \eng{He said me that he was tired.}
2. \eng{She asked me where did I live.}
3. \eng{The teacher told to the students to be quiet.}
4. \eng{He suggested me to apply for the job.}
5. \eng{They asked if they can come tomorrow.}
6. \eng{He promised that he will come back.}
7. \eng{She denied to have seen the document.}
8. \eng{The minister said that the elections are free.} (Contexte : rapport d'hier, élections terminées).
9. \eng{He told that he would come.}
10. \eng{She asked me to not make noise.}
\end{exercice}

\begin{exercice}[Traduction : Français $\to$ Anglais (Style Bac)]
Traduisez en anglais en utilisant le discours indirect.

1. Le ministre a déclaré que le gouvernement protégerait les droits de l'homme.
2. Le journaliste a demandé si des observateurs internationaux seraient invités.
3. L'activiste a ordonné au ministre de publier la liste des détenus.
4. L'officier a conseillé aux électeurs de ne pas toucher aux urnes.
5. Le candidat a promis qu'il organiserait des élections libres l'année suivante.
6. L'élève a demandé où se trouvait le bureau de vote.
7. Le témoin a nié avoir vu l'accusé sur les lieux.
8. Le guide a averti les touristes de ne pas aller près de la frontière.
\end{exercice}

\begin{corrige}[Exercices « Entraînement guidé »]

\textbf{Exercice 1 (Transformation) :}
1. \eng{The deputy declared (that) the new law protected minorities.}
2. \eng{The activists said (that) they had been campaigning for weeks.}
3. \eng{The officer asked the young man if/whether he had registered to vote.}
4. \eng{The journalist asked why the report had not been published.} (Sujet \eng{the report} $\to$ pas d'auxiliaire \eng{did}).
5. \eng{The candidate admitted (that) he could not guarantee the result.}
6. \eng{The supervisor asked where I had put the ballot papers.} (\eng{did put} $\to$ \eng{had put}).

\textbf{Exercice 2 (Verbes introducteurs) :}
1. \eng{The monitor reminded the students to bring their ID cards.}
2. \eng{The mayor promised that he would resign if he lost.} (Ou \eng{promised to resign}).
3. \eng{The nurse advised me to see a doctor about that cough.}
4. \eng{The manager threatened to fire the workers if they struck.} (Ou \eng{threatened that he would fire...}).
5. \eng{The leader suggested organizing a peaceful march / that they (should) organize a peaceful march.}
6. \eng{The suspect denied stealing / having stolen the money.}
7. \eng{The woman begged the police not to arrest her brother.}
8. \eng{The sergeant ordered the recruits to stand up straight.}
9. \eng{The guide warned the tourists not to go near the border zone.}
10. \eng{The student apologized for interrupting the teacher.}

\textbf{Exercice 3 (Correction) :}
1. \eng{He said (that) he was tired.} / \eng{He told me (that) he was tired.}
2. \eng{She asked me where I lived.}
3. \eng{The teacher told the students to be quiet.}
4. \eng{He suggested that I (should) apply for the job.} / \eng{He suggested applying for the job.}
5. \eng{They asked if they could come the next day / the following day.}
6. \eng{He promised that he would come back.} (Ou \eng{promised to come back}).
7. \eng{She denied seeing / having seen the document.}
8. \eng{The minister said that the elections were free.} (Backshift obligatoire : contexte passé).
9. \eng{He told me that he would come.} / \eng{He said that he would come.}
10. \eng{She asked me not to make noise.} (\eng{not to} avant l'infinitif).

\textbf{Exercice 4 (Traduction) :}
1. \eng{The Minister declared (that) the government would protect human rights.}
2. \eng{The journalist asked whether international observers would be invited.}
3. \eng{The activist ordered the Minister to publish the list of detainees.}
4. \eng{The officer advised the voters not to touch the ballot boxes.}
5. \eng{The candidate promised that he would organize free elections the following year.}
6. \eng{The student asked where the polling station was.}
7. \eng{The witness denied seeing / having seen the accused at the scene.}
8. \eng{The guide warned the tourists not to go near the border.}
\end{corrige}

\begin{aretenir}[L'essentiel de la leçon — Fiche de révision Bac]
\begin{itemize}
\item \textbf{Direct speech} = guillemets, paroles exactes. \textbf{Reported speech} = pas de guillemets, verbe introducteur au passé, adaptations.
\item \textbf{Backshift (Règle d'or) :} Présent $\to$ Prétérit ; Prétérit/Perfect $\to$ Plus-que-parfait ; \eng{will/can/may/must} $\to$ \eng{would/could/might/had to}. Modaux \eng{would/could/might/should} inchangés.
\item \textbf{Exceptions :} Vérités générales, report immédiat (verbe introducteur présent), modaux invariables.
\item \textbf{Pronoms/Marqueurs :} \eng{I/we} $\to$ \eng{he/she/they} ; \eng{now} $\to$ \eng{then} ; \eng{today} $\to$ \eng{that day} ; \eng{tomorrow} $\to$ \eng{the next/following day} ; \eng{yesterday} $\to$ \eng{the day before} ; \eng{here} $\to$ \eng{there} ; \eng{this/these} $\to$ \eng{that/those}.
\item \textbf{Questions :} \eng{ask + if/whether} (fermée) ou \eng{ask + wh-} (ouverte) + \textbf{ordre SUJET + VERBE} (pas d'inversion, pas de \eng{do}).
\item \textbf{Ordres/Demandes :} \eng{tell/ask/order/warn/advise + personne + (not) to + BV}. \eng{Let's} $\to$ \eng{suggest + -ing / that + should}.
\item \textbf{Verbes introducteurs :} Maîtriser les 6 structures (\eng{V+that}, \eng{V+obj+that}, \eng{V+to-inf}, \eng{V+obj+to-inf}, \eng{V+-ing}, \eng{V+prep+-ing}).
\item \textbf{Pièges francophones :} \eng{said me} (non), \eng{told that} (non), \eng{asked where did I go} (non), \eng{suggested to go} (non), \eng{denied to do} (non), \eng{not to} vs \eng{to not} (\eng{not to}).
\end{itemize}
\end{aretenir}

\coursec{Le recul des temps (backshift)}

Dans la section précédente, nous avons observ\'e un dialogue de campagne citoyenne et rep\'er\'e que les paroles rapport\'ees ne ressemblent pas mot pour mot aux paroles prononc\'ees~: les temps des verbes ont \guillemotleft\ gliss\'e\guillemotright\ vers le pass\'e. Ce ph\'enom\`ene a un nom~: le \cle{backshift}, c'est-\`a-\`dire le \emph{recul des temps}. C'est le m\'ecanisme central du discours indirect en anglais, et c'est pr\'ecis\'ement l\`a que les francophones commettent le plus d'erreurs, car le fran\c{c}ais fonctionne diff\'eremment. \'Etudions-le m\'ethodiquement.

\courssub{Observation~: le temps recule d'un cran}

\begin{textebox}[At the community meeting --- before and after]
Yesterday evening, the neighbourhood committee met in the school hall. The chairwoman opened the session and announced: \eng{``We \textbf{have collected} enough signatures, and the town council \textbf{will examine} our petition next month.''} A young farmer stood up and declared: \eng{``I \textbf{am working} on a water project, but I \textbf{cannot finish} it alone.''} An old man added: \eng{``We \textbf{waited} too long last year; we \textbf{must act} now.''}

This morning, a journalist reported the meeting on the radio: \eng{``The chairwoman said they \textbf{had collected} enough signatures and that the town council \textbf{would examine} their petition the following month. A young farmer said he \textbf{was working} on a water project but that he \textbf{could not finish} it alone. An old man added that they \textbf{had waited} too long the previous year and that they \textbf{had to act} immediately.''}
\end{textebox}

\textbf{Glossaire~:} \emph{chairwoman} (n.)~: pr\'esidente de s\'eance~; \emph{petition} (n.)~: p\'etition~; \emph{to collect} (v.)~: collecter, r\'eunir~; \emph{to act} (v.)~: agir.

Observez les paires de phrases. \`A chaque fois que le verbe introducteur est au pass\'e (\eng{said}, \eng{announced}, \eng{added}), le temps du verbe rapport\'e \textbf{recule d'un cran vers le pass\'e}~:

\begin{center}
\begin{tabularx}{\linewidth}{|X|X|}
\hline
\rowcolor{popblueL} \textbf{Paroles directes} & \textbf{Paroles rapport\'ees} \\
\hline
We \textbf{have collected} (present perfect) & they \textbf{had collected} (past perfect) \\
\hline
the council \textbf{will examine} (future) & the council \textbf{would examine} (conditionnel) \\
\hline
I \textbf{am working} (present continuous) & he \textbf{was working} (past continuous) \\
\hline
I \textbf{cannot finish} (can) & he \textbf{could not finish} (could) \\
\hline
We \textbf{waited} (past simple) & they \textbf{had waited} (past perfect) \\
\hline
We \textbf{must act} (must) & they \textbf{had to act} (had to) \\
\hline
\end{tabularx}
\end{center}

\courssub{La r\`egle du backshift}

\begin{propriete}[R\`egle du backshift (recul des temps)]
Quand le verbe introducteur est au \textbf{pass\'e} (\eng{said}, \eng{told}, \eng{asked}, \eng{explained}, \eng{announced}...), chaque temps de la phrase rapport\'ee \textbf{recule d'un cran vers le pass\'e}~:
\begin{itemize}
\item \eng{am / is / are} $\rightarrow$ \eng{was / were}
\item \eng{have / has} $\rightarrow$ \eng{had} (auxiliaire du present perfect)
\item \eng{do / does / did} $\rightarrow$ past simple $\rightarrow$ \eng{had} + participe pass\'e
\item \eng{will} $\rightarrow$ \eng{would}~; \eng{can} $\rightarrow$ \eng{could}~; \eng{may} $\rightarrow$ \eng{might}
\end{itemize}
Le past perfect, lui, ne peut pas reculer davantage~: il reste \textbf{inchang\'e}.
\end{propriete}

\begin{popfigure}
\begin{center}
\begin{tikzpicture}[font=\small]
\% Axe du temps
\draw[-{Stealth[length=3mm]}, thick, popdark] (-0.5,0) -- (11.5,0);
\node[below, popmuted] at (11.2,0) {temps};
\% Jalons
\foreach \x/\t in {1.5/{past perfect}, 5/{past}, 8.5/{present}} {
  \fill[popblue] (\x,0) circle (2.5pt);
  \node[above=2pt, popdark] at (\x,0) {\t};
}
\node[above=2pt, popdark] at (10.8,0) {future};
\fill[popblue] (10.8,0) circle (2.5pt);
\% Fl\`eches de backshift
\draw[-{Stealth[length=2.5mm]}, very thick, poporange] (10.8,-0.6) to[bend right=25] node[below, poporange]{will $\rightarrow$ would} (8.5,-0.6);
\draw[-{Stealth[length=2.5mm]}, very thick, poporange] (8.5,-1.5) to[bend right=25] node[below, poporange]{present $\rightarrow$ past} (5,-1.5);
\draw[-{Stealth[length=2.5mm]}, very thick, poporange] (5,-2.4) to[bend right=25] node[below, poporange]{past $\rightarrow$ past perfect} (1.5,-2.4);
\% \'Etiquette
\node[align=center, poppurple, font=\bfseries] at (5.5,1.1) {BACKSHIFT~: chaque temps recule d'un cran};
\end{tikzpicture}
\end{center}
\legende{Le recul des temps~: apr\`es un verbe introducteur au pass\'e, chaque temps glisse d'un cran vers le pass\'e.}
\end{popfigure}

\courssub{Tableau complet des transformations}

\begin{center}
\begin{tabularx}{\linewidth}{|l|X|X|}
\hline
\rowcolor{popblueL} \textbf{Temps (direct)} & \textbf{Temps (rapport\'e)} & \textbf{Exemple} \\
\hline
present simple & past simple & \eng{``I \textbf{vote} every year.''} $\rightarrow$ \eng{He said he \textbf{voted} every year.} \\
\hline
present continuous & past continuous & \eng{``We \textbf{are building} a school.''} $\rightarrow$ \eng{They said they \textbf{were building} a school.} \\
\hline
present perfect & past perfect & \eng{``She \textbf{has signed}.''} $\rightarrow$ \eng{He said she \textbf{had signed}.} \\
\hline
past simple & past perfect & \eng{``We \textbf{marched}.''} $\rightarrow$ \eng{They said they \textbf{had marched}.} \\
\hline
past perfect & past perfect (inchang\'e) & \eng{``I \textbf{had voted}.''} $\rightarrow$ \eng{She said she \textbf{had voted}.} \\
\hline
will & would & \eng{``I \textbf{will help}.''} $\rightarrow$ \eng{He said he \textbf{would help}.} \\
\hline
can & could & \eng{``We \textbf{can win}.''} $\rightarrow$ \eng{They said they \textbf{could win}.} \\
\hline
may & might & \eng{``It \textbf{may rain}.''} $\rightarrow$ \eng{She said it \textbf{might rain}.} \\
\hline
must (obligation) & had to & \eng{``You \textbf{must vote}.''} $\rightarrow$ \eng{He said we \textbf{had to vote}.} \\
\hline
\end{tabularx}
\end{center}

\begin{exemplebox}[Exemples traduits]
\eng{``We have signed the charter,'' they said.} $\rightarrow$ \eng{They said they had signed the charter.}\\
\guillemotleft\ Nous avons sign\'e la charte\guillemotright, ont-ils dit. $\rightarrow$ Ils ont dit qu'ils avaient sign\'e la charte.

\eng{``I am registering new voters,'' the official said.} $\rightarrow$ \eng{The official said he was registering new voters.}\\
\guillemotleft\ J'inscris de nouveaux \'electeurs\guillemotright, a dit le fonctionnaire. $\rightarrow$ Le fonctionnaire a dit qu'il inscrivait de nouveaux \'electeurs.

\eng{``The results will be published tomorrow,'' the mayor announced.} $\rightarrow$ \eng{The mayor announced that the results would be published the next day.}\\
\guillemotleft\ Les r\'esultats seront publi\'es demain\guillemotright, a annonc\'e le maire. $\rightarrow$ Le maire a annonc\'e que les r\'esultats seraient publi\'es le lendemain.

\eng{``We demonstrated peacefully in Bamenda,'' the students said.} $\rightarrow$ \eng{The students said they had demonstrated peacefully in Bamenda.}\\
\guillemotleft\ Nous avons manifest\'e pacifiquement \`a Bamenda\guillemotright, ont dit les \'etudiants. $\rightarrow$ Les \'etudiants ont dit qu'ils avaient manifest\'e pacifiquement \`a Bamenda.
\end{exemplebox}

\courssub{Les modaux invariables et le cas de must}

\begin{propriete}[Modaux qui ne changent jamais]
Certains auxiliaires modaux sont d\'ej\`a \guillemotleft\ au pass\'e\guillemotright\ ou n'ont pas de forme pass\'ee~: ils restent \textbf{inchang\'es} au discours rapport~:
\begin{itemize}
\item \eng{could} $\rightarrow$ \eng{could}~: \eng{``I could not vote in 2018,'' she said.} $\rightarrow$ \eng{She said she could not vote in 2018.}
\item \eng{would} $\rightarrow$ \eng{would}~: \eng{``I would prefer a debate,'' he said.} $\rightarrow$ \eng{He said he would prefer a debate.}
\item \eng{should} $\rightarrow$ \eng{should}~: \eng{``You should read the constitution,'' the teacher said.} $\rightarrow$ \eng{The teacher said we should read the constitution.}
\item \eng{might} $\rightarrow$ \eng{might}~; \eng{ought to} $\rightarrow$ \eng{ought to}.
\end{itemize}
\end{propriete}

\begin{attention}[Must~: obligation ou d\'eduction~?]
Le modal \eng{must} a deux comportements diff\'erents selon son sens~:
\begin{itemize}
\item \textbf{Obligation}~: \eng{must} $\rightarrow$ \eng{had to}. \eng{``You must vote,'' he said.} $\rightarrow$ \eng{He said we had to vote.} (\guillemotleft\ Vous devez voter\guillemotright\ $\rightarrow$ Il a dit que nous devions voter.)
\item \textbf{D\'eduction, certitude logique}~: \eng{must} est \textbf{conserv\'e}. \eng{``It must be late,'' she said.} $\rightarrow$ \eng{She said it must be late.} (\guillemotleft\ Il doit \^etre tard\guillemotright\ $\rightarrow$ Elle a dit qu'il devait \^etre tard.)
\item \textbf{Interdiction}~: \eng{mustn't} est \textbf{conserv\'e}. \eng{``You mustn't cheat,'' the invigilator said.} $\rightarrow$ \eng{The invigilator said we mustn't cheat.}
\end{itemize}
Ne transformez donc jamais m\'ecaniquement \eng{must} en \eng{had to}~: interrogez d'abord son sens.
\end{attention}

\courssub{Les trois exceptions au backshift}

\begin{propriete}[Quand le backshift ne s'applique pas]
\begin{enumerate}
\item \textbf{Verbe introducteur au pr\'esent}~: si l'on rapporte avec \eng{says}, \eng{tells}, \eng{reports}, aucun temps ne change. \eng{He says he is tired.} (Il dit qu'il est fatigu\'e.)
\item \textbf{V\'erit\'es g\'en\'erales et scientifiques}~: un fait toujours vrai garde son pr\'esent. \eng{The teacher said that water boils at 100 degrees.} (Le professeur a dit que l'eau bout \`a 100 degr\'es.)
\item \textbf{Fait encore vrai au moment o\`u l'on parle}~: le backshift est \textbf{facultatif}. \eng{He said he lives in Yaound\'e.} ou \eng{He said he lived in Yaound\'e.} (Il a dit qu'il habite / habitait \`a Yaound\'e --- et il y habite toujours.)
\end{enumerate}
\end{propriete}

\begin{exemplebox}[Exceptions illustr\'ees]
\eng{The journalist reports that the strike continues.} --- Le journaliste rapporte que la gr\`eve continue. (verbe introducteur au pr\'esent~: pas de recul)

\eng{Our geography teacher told us that the earth goes round the sun.} --- Notre professeur de g\'eographie nous a dit que la terre tourne autour du soleil. (v\'erit\'e scientifique~: pr\'esent conserv\'e)

\eng{She said she works for an NGO in Douala.} --- Elle a dit qu'elle travaille pour une ONG \`a Douala. (c'est encore vrai aujourd'hui~: \eng{works} ou \eng{worked} sont tous deux corrects)
\end{exemplebox}

\courssub{Comparaison avec le fran\c{c}ais et erreurs fr\'equentes}

\begin{center}
\begin{tabularx}{\linewidth}{|X|X|X|}
\hline
\rowcolor{popblueL} \textbf{Fran\c{c}ais} & \textbf{English} & \textbf{Remarque} \\
\hline
Il a dit qu'il viendra. & He said he \textbf{would} come. & Le fran\c{c}ais garde le futur~; l'anglais exige \eng{would}. \\
\hline
Elle a dit qu'elle peut nous aider. & She said she \textbf{could} help us. & \eng{can} recule en \eng{could}. \\
\hline
Il a dit qu'il a sign\'e la charte. & He said he \textbf{had signed} the charter. & Le pass\'e compos\'e fran\c{c}ais correspond au \eng{past perfect} anglais. \\
\hline
Elle dit qu'elle est fatigu\'ee. & She says she \textbf{is} tired. & Verbe introducteur au pr\'esent~: pas de backshift, comme en fran\c{c}ais. \\
\hline
\end{tabularx}
\end{center}

\begin{attention}[L'erreur n\textdegree1 des francophones~: le calque du futur]
En fran\c{c}ais, on dit naturellement \guillemotleft\ il a dit qu'il \textbf{viendra}\guillemotright\. Le francophone \'ecrit alors~: \eng{He said he will come.} C'est une \textbf{erreur} en anglais~: apr\`es un verbe introducteur au pass\'e, \eng{will} doit reculer en \eng{would}~: \eng{He said he \textbf{would} come.} M\^eme vigilance pour \eng{can} $\rightarrow$ \eng{could}~: on n'\'ecrira jamais \eng{She said she can come}, mais \eng{She said she \textbf{could} come.}
\end{attention}

\begin{methode}[V\'erifier un discours rapport\'e en trois \'etapes]
\begin{enumerate}
\item \textbf{Rep\'erer le temps du verbe introducteur.} S'il est au pr\'esent (\eng{says}, \eng{tells}), ne rien changer. S'il est au pass\'e (\eng{said}, \eng{told}), passer \`a l'\'etape 2.
\item \textbf{Appliquer syst\'ematiquement le tableau de backshift} \`a chaque verbe de la phrase rapport\'ee~: present $\rightarrow$ past, past $\rightarrow$ past perfect, \eng{will} $\rightarrow$ \eng{would}, \eng{can} $\rightarrow$ \eng{could}.
\item \textbf{Tester les exceptions avant d'\'ecrire}~: s'agit-il d'une v\'erit\'e g\'en\'erale, d'un fait encore vrai, d'un modal invariable (\eng{could}, \eng{should}, \eng{would}, \eng{might}, \eng{ought to}) ou d'un \eng{must} de d\'eduction~? Si oui, conserver la forme d'origine.
\end{enumerate}
\end{methode}

\begin{exoresolu}[Transformation guid\'ee]
\'Enonc\'e~: rapportez les paroles suivantes apr\`es \eng{The activist said that...}~: \eng{``I have organised three workshops, and I will train new volunteers next month. The government must protect every citizen. Young people can change this country.''}
\tcblower
Solution d\'etaill\'ee~:
\begin{enumerate}
\item Le verbe introducteur \eng{said} est au \textbf{pass\'e}~: le backshift s'applique.
\item \eng{have organised} (present perfect) $\rightarrow$ \eng{had organised} (past perfect).
\item \eng{will train} $\rightarrow$ \eng{would train}~; \eng{next month} $\rightarrow$ \eng{the following month}.
\item \eng{must protect} exprime une \textbf{obligation} $\rightarrow$ \eng{had to protect}.
\item \eng{can change} $\rightarrow$ \eng{could change}.
\end{enumerate}
Phrase compl\`ete~: \eng{The activist said that she had organised three workshops and that she would train new volunteers the following month. She added that the government had to protect every citizen and that young people could change the country.}\\
Traduction~: \guillemotleft\ J'ai organis\'e trois ateliers et je formerai de nouveaux volontaires le mois prochain. Le gouvernement doit prot\'eger chaque citoyen. Les jeunes peuvent changer ce pays.\guillemotright\ $\rightarrow$ L'activiste a dit qu'elle avait organis\'e trois ateliers et qu'elle formerait de nouveaux volontaires le mois suivant. Elle a ajout\'e que le gouvernement devait prot\'eger chaque citoyen et que les jeunes pouvaient changer le pays.
\end{exoresolu}

\begin{exercice}[\`A vous de jouer]
Rapportez chaque phrase apr\`es \eng{He said...} ou \eng{She said...}~:
\begin{enumerate}
\item \eng{``I am proud of my country.''}
\item \eng{``We will defend our rights.''}
\item \eng{``The court has delivered its judgment.''}
\item \eng{``Citizens must respect the law.''} (obligation)
\item \eng{``It must be midnight.''} (d\'eduction)
\item \eng{``Honesty is the best policy.''} (v\'erit\'e g\'en\'erale)
\end{enumerate}
\end{exercice}

\begin{corrige}[Exercice~: \`A vous de jouer]
\begin{enumerate}
\item \eng{He said he was proud of his country.} (present simple $\rightarrow$ past simple)
\item \eng{He said they would defend their rights.} (\eng{will} $\rightarrow$ \eng{would})
\item \eng{She said the court had delivered its judgment.} (present perfect $\rightarrow$ past perfect)
\item \eng{He said citizens had to respect the law.} (\eng{must} d'obligation $\rightarrow$ \eng{had to})
\item \eng{She said it must be midnight.} (\eng{must} de d\'eduction~: conserv\'e)
\item \eng{He said honesty is the best policy.} (v\'erit\'e g\'en\'erale~: pas de backshift~; \eng{was} serait aussi tol\'er\'e mais moins naturel)
\end{enumerate}
\end{corrige}

\begin{aretenir}[L'essentiel du backshift]
Apr\`es un verbe introducteur au \textbf{pass\'e}, chaque temps recule d'un cran~: present $\rightarrow$ past, past $\rightarrow$ past perfect, \eng{will} $\rightarrow$ \eng{would}, \eng{can} $\rightarrow$ \eng{could}, \eng{may} $\rightarrow$ \eng{might}. Trois garde-fous~: pas de recul si le verbe introducteur est au pr\'esent, pas de recul pour les v\'erit\'es g\'en\'erales, recul facultatif si le fait est encore vrai. Et n'oubliez jamais~: le fran\c{c}ais dit \guillemotleft\ il a dit qu'il viendra\guillemotright, l'anglais dit \eng{He said he \textbf{would} come}.
\end{aretenir}

\coursec{Pronoms et marqueurs spatio-temporels}

Dans la section précédente, nous avons vu que lorsqu'on rapporte une parole, les temps verbaux reculent d'un cran vers le passé : c'est le \cle{backshift}. Mais transformer une phrase ne s'arrête pas aux verbes. Imaginez la scène : John dit à Marie, au marché de Mokolo, \eng{“I will meet you here tomorrow.”} Le lendemain, Marie raconte la conversation à une amie, au téléphone, depuis chez elle à Nkolbisson. Elle ne peut évidemment pas répéter mot à mot : \eng{John told me “I will meet you here tomorrow”}. Le \eng{I} désignait John, le \eng{you} désignait Marie, le \eng{here} désignait le marché, et le \eng{tomorrow} désignait un jour qui est maintenant passé. Tous ces petits mots — pronoms, possessifs, démonstratifs, marqueurs de temps et de lieu — doivent être \cle{adaptés au nouveau contexte}. C'est l'objet de cette section.

\courssub{Observation : un dialogue rapporté}

\begin{textebox}[A conversation at the market — reported]
Yesterday afternoon, John met Mary at Mokolo market. He said to her: “I am very busy today. My brother and I are selling our cocoa here, but we will finish early. I will meet you here tomorrow at ten o'clock, in front of this shop. These beans are the best in the market, so come quickly!”

The next day, Mary told her friend Grace about the meeting: “John told me that he was very busy that day. He said that his brother and he were selling their cocoa there, but that they would finish early. He told me he would meet me there the following day at ten o'clock, in front of that shop. He also said those beans were the best in the market and asked me to come quickly.”
\end{textebox}

\textbf{Glossaire.} \emph{cocoa} (n.) : cacao ; \emph{beans} (n. plur.) : fèves ; \emph{the following day} : le lendemain ; \emph{that day} : ce jour-là.

\textbf{Questions de compréhension.}
\begin{enumerate}
\item What did John say about his brother? \emph{(Que disait John de son frère ?)}
\item How does Mary report “I will meet you here tomorrow”? \emph{(Comment Marie rapporte-t-elle « I will meet you here tomorrow » ?)}
\item Which words change between John's direct words and Mary's report? \emph{(Quels mots changent entre les paroles directes de John et le compte rendu de Marie ?)}
\end{enumerate}

Observons les transformations en détail :

\begin{center}
\begin{tabularx}{\linewidth}{|X|X|}
\hline
\rowcolor{popblueL} \textbf{Direct speech (John)} & \textbf{Indirect speech (Mary)} \\
\hline
\eng{I am very busy today.} & \eng{He said he was very busy that day.} \\
\hline
\eng{My brother and I are selling our cocoa here.} & \eng{He said his brother and he were selling their cocoa there.} \\
\hline
\eng{I will meet you here tomorrow.} & \eng{He told me he would meet me there the following day.} \\
\hline
\eng{These beans are the best.} & \eng{He said those beans were the best.} \\
\hline
\end{tabularx}
\end{center}

On constate trois familles de changements : les \cle{pronoms et possessifs} (I → he, my → his, our → their, you → me), les \cle{démonstratifs} (this → that, these → those) et les \cle{marqueurs spatio-temporels} (today → that day, here → there, tomorrow → the following day).

\courssub{Les pronoms personnels et possessifs : tout dépend du contexte}

\begin{propriete}[Règle d'adaptation des pronoms]
Au discours indirect, les pronoms personnels et possessifs s'adaptent à la \cle{nouvelle situation d'énonciation} : on se demande \textbf{qui rapporte}, \textbf{de qui on parle} et \textbf{à qui l'on s'adresse}. Il n'existe pas de correspondance mécanique unique : \eng{I} devient \eng{he}, \eng{she} ou reste \eng{I} selon le cas.
\end{propriete}

Reprenons la phrase de John : \eng{“I will meet you here tomorrow.”}

\begin{itemize}
\item Si \textbf{Mary} rapporte : \eng{John told me (that) he would meet me there the following day.} — \eng{I} (John) → \eng{he} ; \eng{you} (Mary) → \eng{me}.
\item Si \textbf{John lui-même} rapporte plus tard : \eng{I told Mary (that) I would meet her there the following day.} — \eng{I} reste \eng{I} ; \eng{you} (Mary) → \eng{her}.
\item Si \textbf{un témoin} rapporte : \eng{John told Mary (that) he would meet her there the following day.} — \eng{I} → \eng{he} ; \eng{you} → \eng{her}.
\end{itemize}

\begin{popfigure}
\begin{center}
\begin{tikzpicture}[node distance=1.1cm and 2.2cm, >=Stealth]
\node[draw=popblue, fill=popblueL, rounded corners, align=center, minimum width=2.6cm] (A) {\textbf{A — John}\\ \eng{I} / \eng{my}};
\node[draw=popgreen, fill=popgreenL, rounded corners, align=center, minimum width=2.6cm, right=of A] (B) {\textbf{B — Mary}\\ \eng{you} / \eng{your}};
\coordinate (mid) at ($(A)!0.5!(B)$);
\node[draw=poporange, fill=poporangeL, rounded corners, align=center, minimum width=3.2cm, below=of mid] (C) {\textbf{C — le rapporteur}\\ (Grace)};
\draw[->, thick, popblue] (A) -- node[above, align=center, font=\small]{\eng{“I will meet you}\\ \eng{here tomorrow.”}} (B);
\draw[->, thick, poporange] (A.south) -- node[left, align=center, font=\small]{\eng{I} $\rightarrow$ \eng{he}\\ \eng{my} $\rightarrow$ \eng{his}} (C.north west);
\draw[->, thick, popgreen] (B.south) -- node[right, align=center, font=\small]{\eng{you} $\rightarrow$ \eng{me/her}\\ \eng{your} $\rightarrow$ \eng{my/her}} (C.north east);
\end{tikzpicture}
\end{center}
\legende{Mini-schéma de situation : A parle à B ; C rapporte. Le choix du pronom dépend de qui parle et de qui est rapporté.}
\end{popfigure}

\begin{center}
\begin{tabularx}{\linewidth}{|X|X|X|}
\hline
\rowcolor{popblueL} \textbf{Discours direct} & \textbf{Discours indirect} & \textbf{Remarque} \\
\hline
\eng{I} & \eng{he / she} (ou \eng{I}) & selon qui rapporte \\
\hline
\eng{we} & \eng{they} (ou \eng{we}) & \eng{“We are ready,” they said.} → \eng{They said they were ready.} \\
\hline
\eng{me} & \eng{him / her} & complément d'objet \\
\hline
\eng{you} & \eng{me / him / her / us / them} & selon le destinataire \\
\hline
\eng{my} & \eng{his / her} & adjectif possessif \\
\hline
\eng{our} & \eng{their} & \eng{our cocoa} → \eng{their cocoa} \\
\hline
\eng{mine} & \eng{his / hers} & pronom possessif : \eng{“The book is mine.”} → \eng{She said the book was hers.} \\
\hline
\eng{yours} & \eng{mine / his / hers} & selon le contexte \\
\hline
\end{tabularx}
\end{center}

\begin{exemplebox}[Exemples traduits]
\eng{“I have lost my keys,” Peter said.} → \eng{Peter said (that) he had lost his keys.} (« J'ai perdu mes clés », a dit Peter. → Peter a dit qu'il avait perdu ses clés.)

\eng{“We will bring our books,” the students promised.} → \eng{The students promised (that) they would bring their books.} (« Nous apporterons nos livres », ont promis les élèves. → Les élèves ont promis qu'ils apporteraient leurs livres.)

\eng{“This victory is mine!” the champion shouted.} → \eng{The champion shouted that that victory was his.} (« Cette victoire est à moi ! » a crié le champion. → Le champion a crié que cette victoire-là était à lui.)
\end{exemplebox}

\courssub{Les démonstratifs : de la proximité à la distance}

\begin{propriete}[Règle des démonstratifs]
Les démonstratifs de \cle{proximité} deviennent des démonstratifs de \cle{distance} au discours indirect, car l'objet n'est généralement plus sous les yeux du rapporteur :
\eng{this} → \eng{that} ; \eng{these} → \eng{those} ; \eng{this place} → \eng{that place}.
\end{propriete}

\begin{exemplebox}[Exemples traduits]
\eng{“I like this shirt,” she said.} → \eng{She said she liked that shirt.} (« J'aime cette chemise-ci », a-t-elle dit. → Elle a dit qu'elle aimait cette chemise-là.)

\eng{“These mangoes are sweet,” the seller said.} → \eng{The seller said those mangoes were sweet.} (« Ces mangues-ci sont sucrées », a dit le vendeur. → Le vendeur a dit que ces mangues-là étaient sucrées.)
\end{exemplebox}

\courssub{Les marqueurs temporels et spatiaux : le tableau à connaître par cœur}

\begin{propriete}[Règle des marqueurs spatio-temporels]
Tout repère de temps ou de lieu lié au \cle{moment de la parole} se décale au discours indirect, car le moment du rapport est différent : \eng{now} → \eng{then}, \eng{today} → \eng{that day}, \eng{tomorrow} → \eng{the following day}, \eng{here} → \eng{there}.
\end{propriete}

\begin{center}
\begin{tabularx}{\linewidth}{|X|X|X|}
\hline
\rowcolor{popblueL} \textbf{Direct speech} & \textbf{Indirect speech} & \textbf{Français} \\
\hline
\eng{now} & \eng{then} & maintenant → alors / à ce moment-là \\
\hline
\eng{today} & \eng{that day} & aujourd'hui → ce jour-là \\
\hline
\eng{yesterday} & \eng{the day before / the previous day} & hier → la veille \\
\hline
\eng{tomorrow} & \eng{the next day / the following day} & demain → le lendemain \\
\hline
\eng{ago} & \eng{before / earlier} & il y a (deux jours) → (deux jours) auparavant \\
\hline
\eng{last week / last night} & \eng{the week before / the previous week ; the night before} & la semaine dernière → la semaine précédente \\
\hline
\eng{next week / next year} & \eng{the following week / the following year} & la semaine prochaine → la semaine suivante \\
\hline
\eng{tonight} & \eng{that night} & ce soir → ce soir-là \\
\hline
\eng{here} & \eng{there} & ici → là \\
\hline
\eng{this place} & \eng{that place} & cet endroit-ci → cet endroit-là \\
\hline
\end{tabularx}
\end{center}

\begin{popfigure}
\begin{center}
\begin{tikzpicture}[node distance=0.3cm and 1.6cm, >=Stealth]
\node[draw=popblue, fill=popblue, text=white, rounded corners, minimum width=4.2cm, font=\bfseries] (h1) {Direct};
\node[draw=popgreen, fill=popgreen, text=white, rounded corners, minimum width=4.2cm, right=of h1] (h2) {Indirect};

\node[draw=popline, fill=popblueL, rounded corners, minimum width=4.2cm, below=of h1] (l1) {\eng{now}};
\node[draw=popline, fill=popgreenL, rounded corners, minimum width=4.2cm, right=of l1] (r1) {\eng{then}};
\draw[->, thick, poporange] (l1) -- (r1);

\node[draw=popline, fill=popblueL, rounded corners, minimum width=4.2cm, below=of l1] (l2) {\eng{today}};
\node[draw=popline, fill=popgreenL, rounded corners, minimum width=4.2cm, right=of l2] (r2) {\eng{that day}};
\draw[->, thick, poporange] (l2) -- (r2);

\node[draw=popline, fill=popblueL, rounded corners, minimum width=4.2cm, below=of l2] (l3) {\eng{here}};
\node[draw=popline, fill=popgreenL, rounded corners, minimum width=4.2cm, right=of l3] (r3) {\eng{there}};
\draw[->, thick, poporange] (l3) -- (r3);

\node[draw=popline, fill=popblueL, rounded corners, minimum width=4.2cm, below=of l3] (l4) {\eng{this}};
\node[draw=popline, fill=popgreenL, rounded corners, minimum width=4.2cm, right=of l4] (r4) {\eng{that}};
\draw[->, thick, poporange] (l4) -- (r4);
\end{tikzpicture}
\end{center}
\legende{Tableau de transformation : les marqueurs de proximité (temps et lieu) deviennent des marqueurs de distance au discours indirect.}
\end{popfigure}

\begin{exemplebox}[Exemples traduits]
\eng{“I will meet you here tomorrow,” John told me.} → \eng{John told me he would meet me there the following day.} (« Je te retrouverai ici demain », m'a dit John. → John m'a dit qu'il me retrouverait là le lendemain.)

\eng{“I saw the match yesterday,” Eric said.} → \eng{Eric said he had seen the match the day before.} (« J'ai vu le match hier », a dit Eric. → Eric a dit qu'il avait vu le match la veille.)

\eng{“We are leaving now,” they announced.} → \eng{They announced that they were leaving then.} (« Nous partons maintenant », ont-ils annoncé. → Ils ont annoncé qu'ils partaient à ce moment-là.)

\eng{“I arrived two days ago,” she explained.} → \eng{She explained that she had arrived two days before.} (« Je suis arrivée il y a deux jours », a-t-elle expliqué. → Elle a expliqué qu'elle était arrivée deux jours auparavant.)
\end{exemplebox}

\begin{attention}[Exception : quand le marqueur ne change pas]
Si l'on rapporte les paroles \cle{le jour même} et \cle{au même endroit}, les marqueurs de temps et de lieu peuvent rester identiques :
\eng{“I will see you today,” he said.} → (le même jour) \eng{He said he would see me today.} — correct et naturel.
De même, si l'on est encore sur place : \eng{She said she would wait here.} Le changement n'est obligatoire que lorsque le contexte a réellement changé. Au Bac, en l'absence de précision, on applique la transformation complète : c'est la réponse attendue.
\end{attention}

\begin{methode}[Transformer en trois étapes]
\begin{enumerate}
\item \textbf{Verbe introducteur} : choisir \eng{say} / \eng{tell} (+ objet) et placer \eng{that} (facultatif).
\item \textbf{Backshift des temps} : faire reculer chaque temps verbal d'un cran (will → would, am → was, have seen → had seen).
\item \textbf{Pronoms et marqueurs} : adapter pronoms, possessifs, démonstratifs, temps et lieu au nouveau contexte.
\end{enumerate}
Relire ensuite la phrase rapportée entièrement : chaque mot du discours direct a-t-il été vérifié ?
\end{methode}

\begin{exoresolu}[Exercice de transformation]
Énoncé. Mettez au discours indirect, en commençant par \eng{She told me (that)...} : \eng{“I bought these shoes here yesterday, and I will wear them to the party tomorrow,” Ngono said to me.}
\tcblower
Solution détaillée. Étape 1 : verbe introducteur → \eng{She told me (that)}. Étape 2 : backshift → \eng{bought} (past simple) devient \eng{had bought} ; \eng{will wear} devient \eng{would wear}. Étape 3 : pronoms et marqueurs → \eng{I} devient \eng{she} ; \eng{these shoes} devient \eng{those shoes} ; \eng{here} devient \eng{there} ; \eng{yesterday} devient \eng{the day before} ; \eng{tomorrow} devient \eng{the following day}.
Réponse complète : \eng{She told me (that) she had bought those shoes there the day before, and that she would wear them to the party the following day.} (Ngono m'a dit qu'elle avait acheté ces chaussures-là là-bas la veille, et qu'elle les porterait à la fête le lendemain.)
\end{exoresolu}

\begin{attention}[Erreurs fréquentes des francophones]
\begin{itemize}
\item \textbf{Conserver \eng{I} tel quel} : les francophones oublient souvent de transformer les pronoms. Le \eng{I} de la citation devient \eng{he} ou \eng{she} selon le locuteur, \textbf{jamais} conservé tel quel (sauf si le locuteur rapporte ses propres paroles). Faux : \eng{John said I was tired.} (si John parlait de lui-même) ; juste : \eng{John said he was tired.}
\item \textbf{Oublier \eng{the} devant les marqueurs} : on écrit \eng{the day before}, \eng{the following day}, \eng{the previous week} — jamais \eng{day before} seul.
\item \textbf{Traduire « la veille » par \eng{the eve}} : \eng{the eve} signifie « la veille » au sens de « la veille de Noël » (\eng{Christmas Eve}), pas « le jour précédent ».
\item \textbf{Confondre \eng{ago} et \eng{before}} : \eng{ago} se compte depuis maintenant ; au discours rapporté, on utilise \eng{before} : \eng{two days ago} → \eng{two days before}.
\end{itemize}
\end{attention}

\begin{savaistu}[« La veille » et « le lendemain » au Bac]
\eng{The day before} et \eng{the following day} sont les équivalents les plus fréquents de « la veille » et « le lendemain » dans les épreuves de transformation au Baccalauréat. Les apprendre par cœur fait gagner des points précieux : chaque année, des candidats perdent des points en écrivant \eng{the day after} (possible, mais moins attendu) ou en oubliant l'article \eng{the}. Retenez le trio magique : \eng{yesterday → the day before}, \eng{tomorrow → the following day}, \eng{today → that day}.
\end{savaistu}

\begin{exercice}[À vous de jouer]
Mettez les phrases suivantes au discours indirect :
\begin{enumerate}
\item \eng{“I am waiting for my sister here,” Divine said.}
\item \eng{“We played football in this field last week,” the boys said.}
\item \eng{“I will visit my grandmother tomorrow,” Amina told me.}
\item \eng{“These documents are mine,” the lawyer said.}
\end{enumerate}
\end{exercice}

\begin{corrige}[Exercice « À vous de jouer »]
\begin{enumerate}
\item \eng{Divine said (that) she was waiting for her sister there.} (I → she ; my → her ; am → was ; here → there.)
\item \eng{The boys said (that) they had played football in that field the week before.} (we → they ; played → had played ; this → that ; last week → the week before.)
\item \eng{Amina told me (that) she would visit her grandmother the following day.} (I → she ; will → would ; my → her ; tomorrow → the following day.)
\item \eng{The lawyer said (that) those documents were his.} (these → those ; are → were ; mine → his.)
\end{enumerate}
\end{corrige}

\begin{aretenir}[L'essentiel de la section]
Au discours indirect, tout ce qui était « proche » du locuteur devient « distant » : les personnes (I → he/she, my → his/her), les objets (this → that, these → those), le temps (now → then, today → that day, tomorrow → the following day) et le lieu (here → there). Le choix exact du pronom dépend toujours du contexte : qui rapporte, de qui l'on parle, à qui l'on s'adresse. Exception : si le rapport a lieu le jour même et au même endroit, les marqueurs peuvent rester inchangés.
\end{aretenir}

\coursec{Rapporter questions, ordres et demandes}

Dans la section précédente, nous avons vu comment les pronoms et les marqueurs spatio-temporels changent lorsqu'on rapporte des paroles (\eng{I} devient \eng{he/she}, \eng{today} devient \eng{that day}, etc.). Mais jusqu'ici, nous n'avons rapporté que des \emph{affirmations}. Or, dans la vie quotidienne — au commissariat, au tribunal, lors d'une interview ou d'une réunion de quartier — on rapporte aussi des \emph{questions}, des \emph{ordres} et des \emph{demandes}. C'est précisément l'objet de cette section : comment dire en anglais « il m'a demandé si… », « elle m'a demandé où… », « le professeur nous a dit de… » ?

\courssub{Observation : un dialogue au bureau des élections}

\begin{textebox}[At the electoral office — a dialogue]
A journalist is interviewing Mrs. Ngono, an election observer in Yaoundé, after the vote.

\textbf{Journalist:} “Did you see any problems during the vote?”

\textbf{Mrs. Ngono:} “No, everything was calm.”

\textbf{Journalist:} “Where were you at six o'clock?”

\textbf{Mrs. Ngono:} “I was at the polling station near the market.”

\textbf{Journalist:} “Why did some voters leave early?”

\textbf{Mrs. Ngono:} “Because the queue was very long.”

\textbf{Journalist:} “Can you describe the counting?”

\textbf{Mrs. Ngono:} “Yes. It was transparent.”

Later, the journalist reports the interview to his editor:

“I asked Mrs. Ngono \textbf{if she had seen} any problems during the vote. She answered that everything had been calm. I asked her \textbf{where she had been} at six o'clock, and she replied that she had been at the polling station near the market. I also asked \textbf{why some voters had left} early. Finally, I asked her \textbf{to describe} the counting, and she said it had been transparent.”
\end{textebox}

\textbf{Glossaire :} \eng{polling station} (n.) : bureau de vote ; \eng{queue} (n.) : file d'attente ; \eng{counting} (n.) : dépouillement ; \eng{transparent} (adj.) : transparent ; \eng{editor} (n.) : rédacteur en chef.

\textbf{Questions de compréhension :}
\begin{enumerate}
\item What did the journalist ask about the counting?
\item Where was Mrs. Ngono at six o'clock?
\item In the reported version, what happens to the word order of the questions?
\end{enumerate}

\textbf{Observation guidée.} Comparez les questions directes et les questions rapportées :

\begin{center}
\begin{tabularx}{\linewidth}{|X|X|}
\hline
\rowcolor{popblueL} \textbf{Question directe} & \textbf{Question rapportée} \\
\hline
“Did you see any problems?” & He asked if she \textbf{had seen} any problems. \\
\hline
“Where were you at six o'clock?” & He asked her where she \textbf{had been} at six o'clock. \\
\hline
“Why did some voters leave early?” & He asked why some voters \textbf{had left} early. \\
\hline
“Can you describe the counting?” & He asked her \textbf{to describe} the counting. \\
\hline
\end{tabularx}
\end{center}

Que remarquez-vous ? Dans les questions rapportées, l'inversion sujet–verbe a disparu, l'auxiliaire \eng{do/does/did} a disparu, et il n'y a plus de point d'interrogation. C'est la règle fondamentale de cette section.

\courssub{Les questions fermées (oui/non) : ask + if/whether}

\begin{propriete}[Règle : rapporter une question fermée]
Une \cle{question fermée} est une question à laquelle on répond par \eng{yes} ou \eng{no}. Pour la rapporter :
\begin{center}
\textbf{ask (+ complément) + if / whether + sujet + verbe (avec backshift)}
\end{center}
\begin{itemize}
\item \eng{if} et \eng{whether} signifient tous les deux « si » et sont interchangeables ici ; \eng{whether} est un peu plus soutenu et s'impose devant \eng{or not} placé juste après : \eng{He asked whether or not I would vote.}
\item On retrouve l'\textbf{ordre de la phrase affirmative} : sujet + verbe, \textbf{sans inversion}.
\item L'auxiliaire \eng{do/does/did} \textbf{disparaît} (le verbe principal porte alors la marque du temps).
\item Il n'y a \textbf{pas de point d'interrogation} : la question rapportée est une subordonnée, pas une vraie question.
\item Le backshift s'applique normalement : present simple $\rightarrow$ past simple, present perfect $\rightarrow$ past perfect, \eng{will} $\rightarrow$ \eng{would}, \eng{can} $\rightarrow$ \eng{could}.
\item Autres verbes introducteurs possibles : \eng{wonder} (se demander), \eng{want to know} (vouloir savoir) : \eng{I wondered if she had voted.} (« Je me demandais si elle avait voté. »)
\end{itemize}
\end{propriete}

\begin{exemplebox}[Questions fermées rapportées — exemples traduits]
\eng{“Do you know your rights?”} $\rightarrow$ \eng{He asked me if I knew my rights.} (« Connais-tu tes droits ? » $\rightarrow$ Il m'a demandé si je connaissais mes droits.)

\eng{“Have you registered to vote?”} $\rightarrow$ \eng{She asked me whether I had registered to vote.} (« T'es-tu inscrit pour voter ? » $\rightarrow$ Elle m'a demandé si je m'étais inscrit pour voter.)

\eng{“Will you attend the meeting?”} $\rightarrow$ \eng{They asked us if we would attend the meeting.} (« Assisterez-vous à la réunion ? » $\rightarrow$ Ils nous ont demandé si nous assisterions à la réunion.)

\eng{“Can you speak English?”} $\rightarrow$ \eng{The officer asked him if he could speak English.} (« Pouvez-vous parler anglais ? » $\rightarrow$ L'officier lui a demandé s'il pouvait parler anglais.)
\end{exemplebox}

\begin{attention}[Faux ami : « si »]
Le « si » français se traduit par \eng{if} ou \eng{whether} dans les questions rapportées, mais attention : \eng{if} ne se traduit pas toujours par « si » ! Dans \eng{He asked me if I knew}, \eng{if} = « si » (interrogation indirecte). Ne confondez pas avec le « si » de réponse (\eng{yes}) ni avec le « si » hypothétique, qui suit les mêmes règles de forme mais un sens différent.
\end{attention}

\courssub{Les questions ouvertes (wh-) : ask + mot interrogatif + sujet + verbe}

\begin{propriete}[Règle : rapporter une question ouverte]
Une \cle{question ouverte} commence par un mot interrogatif : \eng{where, when, why, who, what, which, how, how many, how old}… Pour la rapporter :
\begin{center}
\textbf{ask (+ complément) + mot interrogatif + sujet + verbe (avec backshift)}
\end{center}
Le mot interrogatif joue le rôle de connecteur ; on n'ajoute \textbf{pas} \eng{if}. Là encore : ordre sujet + verbe, pas d'inversion, pas d'auxiliaire \eng{do/does/did}, pas de point d'interrogation.
\end{propriete}

\begin{exemplebox}[Questions ouvertes rapportées — exemples traduits]
\eng{“Where do you live?”} $\rightarrow$ \eng{She asked me where I lived.} (« Où habites-tu ? » $\rightarrow$ Elle m'a demandé où j'habitais.)

\eng{“Why did you protest?” the journalist asked the students.} $\rightarrow$ \eng{The journalist asked the students why they had protested.} (« Pourquoi avez-vous manifesté ? » a demandé le journaliste aux étudiants. $\rightarrow$ Le journaliste a demandé aux étudiants pourquoi ils avaient manifesté.)

\eng{“When will the results be published?”} $\rightarrow$ \eng{People asked when the results would be published.} (« Quand les résultats seront-ils publiés ? » $\rightarrow$ Les gens ont demandé quand les résultats seraient publiés.)

\eng{“How old are you?”} $\rightarrow$ \eng{He asked me how old I was.} (« Quel âge as-tu ? » $\rightarrow$ Il m'a demandé quel âge j'avais.)
\end{exemplebox}

\begin{attention}[Le piège majeur des francophones]
Par calque de la question directe, beaucoup d'élèves écrivent : \emph{He asked me where did I live.} C'est \textbf{faux} ! Dans la question rapportée, il n'y a ni inversion ni auxiliaire \eng{did}. La bonne forme est : \eng{He asked me where I lived.} De même, on n'écrit jamais \emph{She asked me what was I doing} mais \eng{She asked me what I was doing}. Retenez : \textbf{une question rapportée n'a pas la forme d'une question}.
\end{attention}

\begin{methode}[Transformer une question directe en question rapportée]
\begin{enumerate}
\item Repérez le type de question : fermée (\eng{if/whether}) ou ouverte (garder le mot interrogatif).
\item \textbf{Barrez mentalement} l'auxiliaire \eng{do/does/did} : il disparaît.
\item \textbf{Remettez le sujet avant le verbe} (ordre de la phrase affirmative).
\item Appliquez le \textbf{backshift} : le verbe recule d'un temps dans le passé.
\item Adaptez les pronoms et les marqueurs spatio-temporels (\eng{you} $\rightarrow$ \eng{I/me}, \eng{here} $\rightarrow$ \eng{there}…).
\item Terminez par un \textbf{point}, jamais un point d'interrogation.
\end{enumerate}
Exemple filé : “Where did you go?” $\rightarrow$ (2) on barre \eng{did} $\rightarrow$ (3) \eng{you went} $\rightarrow$ (4) backshift : \eng{you had gone} $\rightarrow$ (5) \eng{I had gone} $\rightarrow$ \eng{He asked me where I had gone.}
\end{methode}

\begin{popfigure}
\begin{tikzpicture}[
  node distance=0.55cm and 0.9cm,
  q/.style={diamond, draw=poporange, fill=poporangeL, align=center, inner sep=2pt, aspect=2, font=\small},
  box/.style={rectangle, rounded corners, draw=popblue, fill=popblueL, align=center, inner sep=4pt, font=\small},
  lab/.style={font=\scriptsize\itshape, text=popdark},
  arr/.style={-{Stealth[length=2.5mm]}, thick, popdark}
]
\node[q, align=center] (start){Énoncé à\\rapporter ?};
\node[box, below left=0.8cm and 1.6cm of start, align=center] (st){Affirmation\\say / tell + \textbf{that}\\He said (that) he voted.};
\node[box, below=0.8cm of start, align=center] (yesno){Question oui/non\\ask + \textbf{if/whether} + S + V\\He asked if I voted.};
\node[box, below right=0.8cm and 0.4cm of start, align=center] (wh){Question wh-\\ask + \textbf{wh-} + S + V\\He asked where I lived.};
\node[box, right=0.9cm of wh, align=center] (order){Ordre / demande\\tell / ask + \textbf{to}-infinitive\\He told me to vote.};
\draw[arr] (start) -- node[lab, above left]{affirmation} (st);
\draw[arr] (start) -- node[lab, right]{oui / non ?} (yesno);
\draw[arr] (start) -- node[lab, above right]{wh- ?} (wh);
\draw[arr] (start.east) -| node[lab, above, pos=0.25]{ordre} (order.north);
\end{tikzpicture}
\legende{Organigramme de décision : quel verbe introducteur et quelle structure choisir selon le type d'énoncé à rapporter.}
\end{popfigure}

\courssub{Rapporter les ordres : tell/order + complément + to-infinitive}

\begin{propriete}[Règle : rapporter un ordre ou une interdiction]
Un ordre (impératif) se rapporte avec la structure :
\begin{center}
\textbf{tell / order / warn + complément d'objet + to-infinitive}
\end{center}
Pour un ordre \textbf{négatif}, la négation porte sur l'infinitif :
\begin{center}
\textbf{tell / order / warn + complément + \textbf{not to}-infinitive}
\end{center}
\begin{itemize}
\item \eng{tell} = dire à quelqu'un de (neutre) ; \eng{order} = ordonner (autorité) ; \eng{warn} = avertir, mettre en garde.
\item Le complément d'objet est \textbf{obligatoire} : on dit à \emph{quelqu'un} de faire quelque chose (\eng{He told \textbf{me} to sit down}).
\item Pas de backshift ici : l'infinitif est invariable.
\end{itemize}
\end{propriete}

\begin{exemplebox}[Ordres rapportés — exemples traduits]
\eng{“Sit down!”} $\rightarrow$ \eng{The teacher told us to sit down.} (« Asseyez-vous ! » $\rightarrow$ Le professeur nous a dit de nous asseoir.)

\eng{“Show me your identity card!”} $\rightarrow$ \eng{The policeman ordered him to show his identity card.} (« Montrez-moi votre carte d'identité ! » $\rightarrow$ Le policier lui a ordonné de montrer sa carte d'identité.)

\eng{“Don't touch the ballot boxes!”} $\rightarrow$ \eng{The officer warned us not to touch the ballot boxes.} (« Ne touchez pas aux urnes ! » $\rightarrow$ L'officier nous a avertis de ne pas toucher aux urnes.)

\eng{“Don't cheat!”} $\rightarrow$ \eng{He told us not to cheat.} (« Ne trichez pas ! » $\rightarrow$ Il nous a dit de ne pas tricher.)
\end{exemplebox}

\begin{attention}[Place de la négation]
En français, on dit « il m'a dit de ne pas partir » ; en anglais, \eng{not} se place \textbf{avant} le \eng{to} : \eng{He told me \textbf{not to} leave}, et jamais \emph{He told me to not leave} (forme à éviter au niveau scolaire) ni \emph{He told me don't leave}.
\end{attention}

\courssub{Demandes polies, conseils et suggestions}

\begin{propriete}[Règle : demandes, conseils, suggestions]
\begin{itemize}
\item \textbf{Demande polie} : \textbf{ask + complément + to-infinitive}. \eng{“Can you help me?”} $\rightarrow$ \eng{She asked me to help her.} (« Peux-tu m'aider ? » $\rightarrow$ Elle m'a demandé de l'aider.)
\item \textbf{Conseil} : \textbf{advise + complément + to-infinitive}. \eng{“You should vote.”} $\rightarrow$ \eng{He advised me to vote.} (« Tu devrais voter. » $\rightarrow$ Il m'a conseillé de voter.)
\item \textbf{Suggestion} : \textbf{suggest + -ing} ou \textbf{suggest + that + sujet + should + base verbale}. \eng{“Let's organize a debate.”} $\rightarrow$ \eng{She suggested organizing a debate.} / \eng{She suggested that we should organize a debate.} (« Organisons un débat. » $\rightarrow$ Elle a suggéré d'organiser un débat.)
\end{itemize}
\end{propriete}

\begin{attention}[Deux constructions à ne pas confondre]
On dit \eng{He advised \textbf{me} to vote} (avec complément) mais jamais \emph{He suggested me to vote} : \eng{suggest} ne prend \textbf{pas} de complément d'objet personne suivi de l'infinitif. Dites : \eng{He suggested voting} ou \eng{He suggested that I (should) vote}. De même, pour une demande polie, la structure attendue au Baccalauréat est \eng{She asked me to help her}, et non une subordonnée en \eng{that}.
\end{attention}

\begin{popfigure}
\begin{tikzpicture}[node distance=0.5cm, font=\small]
\node[rectangle, rounded corners, draw=popgreen, fill=popgreenL, align=center, inner sep=5pt] (direct) {Question directe\\[2pt]\eng{“Do you vote?”}};
\node[rectangle, rounded corners, draw=popblue, fill=popblueL, align=center, inner sep=5pt, right=3.4cm of direct] (indirect) {Question rapportée\\[2pt]\eng{He asked if I \textbf{voted}.}};
\draw[-{Stealth[length=3mm]}, very thick, poporange] (direct) -- node[above, align=center, font=\scriptsize]{if + sujet + verbe\\backshift} (indirect);
\node[draw=poppink, thick, inner sep=2pt, circle, font=\footnotesize\bfseries, text=poppink] at ($(direct.east)+(0.15,0.35)$) {X};
\node[font=\scriptsize, text=poppink, align=center] at ($(direct.east)+(0.15,-0.55)$) {pas de\\“do”};
\node[font=\scriptsize, text=poppink, align=center] at ($(indirect.west)+(-0.2,-0.75)$) {pas d'inversion\\pas de “?”};
\end{tikzpicture}
\legende{Schéma de transformation : l'auxiliaire \eng{do} et l'inversion disparaissent ; le verbe recule d'un temps (\eng{vote} $\rightarrow$ \eng{voted}).}
\end{popfigure}

\courssub{Tableau récapitulatif et entraînement}

\begin{center}
\begin{tabularx}{\linewidth}{|l|X|X|}
\hline
\rowcolor{popblueL} \textbf{Type d'énoncé} & \textbf{Structure} & \textbf{Exemple} \\
\hline
Question oui/non & ask + if/whether + S + V & He asked me if I knew my rights. \\
\hline
Question wh- & ask + wh- + S + V & She asked me where I lived. \\
\hline
Ordre & tell/order + obj. + to-inf. & The teacher told us to sit down. \\
\hline
Interdiction & tell/warn + obj. + not to-inf. & He warned us not to touch the boxes. \\
\hline
Demande polie & ask + obj. + to-inf. & She asked me to help her. \\
\hline
Conseil & advise + obj. + to-inf. & He advised me to vote. \\
\hline
Suggestion & suggest + -ing / that + should & She suggested organizing a debate. \\
\hline
\end{tabularx}
\end{center}

\begin{exoresolu}[Transformez les énoncés suivants en discours indirect]
Énoncé. Reportez : 1. “Do you understand the law?” (the lawyer / the villagers) — 2. “Why did you write to the minister?” (the journalist / the activist) — 3. “Don't shout in the courtroom!” (the judge / the public) — 4. “Please sign here.” (the clerk / me) — 5. “Let's form an association.” (my friend).
\tcblower
Solution détaillée.

1. Question fermée $\rightarrow$ \eng{if} + ordre affirmatif + backshift (\eng{understand} $\rightarrow$ \eng{understood}) : \eng{The lawyer asked the villagers if they understood the law.} (« L'avocat a demandé aux villageois s'ils comprenaient la loi. »)

2. Question ouverte $\rightarrow$ on garde \eng{why}, on supprime \eng{did}, backshift (\eng{write} $\rightarrow$ \eng{had written}) : \eng{The journalist asked the activist why he had written to the minister.} (« Le journaliste a demandé au militant pourquoi il avait écrit au ministre. »)

3. Ordre négatif $\rightarrow$ \eng{not to}-infinitive : \eng{The judge told the public not to shout in the courtroom.} (« Le juge a dit au public de ne pas crier dans la salle d'audience. »)

4. Demande polie $\rightarrow$ \eng{ask + me + to}-infinitive : \eng{The clerk asked me to sign there.} (« L'employé m'a demandé de signer là. » — attention : \eng{here} devient \eng{there}.)

5. Suggestion $\rightarrow$ \eng{suggest + -ing} ou \eng{that + should} : \eng{My friend suggested forming an association.} / \eng{My friend suggested that we should form an association.} (« Mon ami a suggéré de créer une association. »)
\end{exoresolu}

\begin{exercice}[À vous de jouer]
Mettez les énoncés suivants au discours indirect :
\begin{enumerate}
\item “Have you ever voted?” (the teacher / the students)
\item “What time does the meeting start?” (a citizen / me)
\item “Don't forget your voter's card!” (my mother / me)
\item “Could you explain this article of the Constitution?” (a pupil / the teacher)
\item “You should read the Universal Declaration of Human Rights.” (the teacher / us)
\end{enumerate}
\end{exercice}

\begin{corrige}[Exercice « À vous de jouer »]
\begin{enumerate}
\item \eng{The teacher asked the students if they had ever voted.}
\item \eng{A citizen asked me what time the meeting started.}
\item \eng{My mother told me not to forget my voter's card.}
\item \eng{A pupil asked the teacher to explain that article of the Constitution.}
\item \eng{The teacher advised us to read the Universal Declaration of Human Rights.}
\end{enumerate}
\end{corrige}

\begin{aretenir}[L'essentiel de la section]
\begin{itemize}
\item Question fermée : \eng{ask + if/whether + sujet + verbe} ; question ouverte : \eng{ask + wh- + sujet + verbe}.
\item Dans toute question rapportée : \textbf{ordre affirmatif}, pas d'auxiliaire \eng{do/does/did}, pas de point d'interrogation, backshift obligatoire.
\item Ordre : \eng{tell/order + complément + to-infinitive} ; interdiction : \eng{not to}-infinitive ; demande polie : \eng{ask + complément + to-infinitive}.
\item Conseil : \eng{advise + complément + to-infinitive} ; suggestion : \eng{suggest + -ing} ou \eng{suggest + that + should}.
\item Piège n°1 : \emph{He asked me where did I live} est faux ; la bonne forme est \eng{He asked me where I lived}.
\end{itemize}
\end{aretenir}

\coursec{Verbes introducteurs et comparaison avec le français}

Dans la section précédente, nous avons appris à rapporter les questions (\eng{ask}, \eng{wonder}, \eng{want to know}), les ordres (\eng{tell}, \eng{order}) et les demandes (\eng{ask}, \eng{beg}). Mais la vie d'une phrase rapportée ne se résume pas à ces quelques verbes. Un bon élève de Terminale~— et un futur bachelier~— doit savoir \cle{varier les verbes introducteurs} pour rendre compte avec précision de l'intention du locuteur. Dire \eng{He said he didn't do it} est correct~; écrire \eng{He denied doing it} est meilleur, plus court, plus élégant. C'est cette richesse que nous allons conquérir maintenant.

\courssub{Observation : un même discours, plusieurs verbes}

\begin{textebox}[A student council meeting in Bamenda]
At yesterday's student council meeting, the prefect \textbf{announced} that the school would organise a debate on human rights. The principal \textbf{advised} all candidates to prepare their arguments carefully. She also \textbf{reminded} the students to register before Friday and \textbf{warned} them not to copy texts from the internet. One student \textbf{suggested} holding the debate in both English and French, and another \textbf{offered} to design the posters. The principal \textbf{agreed} to provide a small budget and \textbf{promised} that the winners would receive books. However, a teacher later \textbf{denied} having approved the date, and \textbf{insisted} that the timetable should be checked again.
\end{textebox}

\begin{definition}[Glossaire]
\begin{itemize}
\item \eng{prefect} (n.) : élève chargé de la discipline (préfet / censeur-élève)
\item \eng{to announce} (v.) : annoncer
\item \eng{to register} (v.) : s'inscrire
\item \eng{to hold a debate} (v.) : organiser / tenir un débat
\item \eng{budget} (n.) : budget
\item \eng{timetable} (n.) : emploi du temps
\end{itemize}
\end{definition}

Observons~: chaque verbe introducteur (\eng{announce, advise, remind, warn, suggest, offer, agree, promise, deny, insist}) n'est pas interchangeable. Chacun porte une \cle{intention} précise~: informer, conseiller, avertir, proposer, promettre, nier. De plus, chacun impose une \cle{construction grammaticale} particulière~: \eng{that}, \eng{to}-infinitif ou \eng{-ing}. C'est le cœur de cette section.

\courssub{Say et tell : la règle d'or}

\begin{propriete}[La règle say / tell]
\begin{itemize}
\item \cle{Say} est suivi directement de la parole rapportée~; la personne, si elle est mentionnée, est introduite par \eng{to}~: \eng{He said (that) he was tired.} / \eng{He said to me that he was tired.}
\item \cle{Tell} exige \textbf{obligatoirement} un complément de personne~: \eng{He told me (that) he was tired.}
\item Interdits~: \eng{He told that he was tired.} (il manque la personne) et \eng{He said me that he was tired.} (\eng{say} ne prend jamais de complément de personne direct).
\end{itemize}
\end{propriete}

\begin{exemplebox}[Say ou tell ?]
\begin{itemize}
\item \eng{The mayor said that the bridge would be repaired.} «~Le maire a dit que le pont serait réparé.~»
\item \eng{The mayor told the journalists that the bridge would be repaired.} «~Le maire a dit aux journalistes que le pont serait réparé.~»
\item \eng{She said to me, “I am proud of you.”} «~Elle m'a dit : “Je suis fière de toi.”~» (style direct~: \eng{say to} est possible)
\item \eng{She told me she was proud of me.} «~Elle m'a dit qu'elle était fière de moi.~»
\end{itemize}
\end{exemplebox}

\begin{attention}[Piège n°~1 des francophones]
Le français dit «~il m'a dit de venir~», ce qui pousse à écrire \eng{He said me to come}. C'est \textbf{faux}. L'anglais exige~: \eng{He told me to come} ou \eng{He asked me to come}. Retenez~: \eng{say} ne se construit \textbf{jamais} avec une personne directement après lui. De même, on ne dit pas \eng{tell to me}~: après \eng{tell}, la personne vient \textbf{sans} \eng{to}.
\end{attention}

\courssub{Les trois grandes constructions des verbes introducteurs}

\begin{popfigure}
\begin{center}\tcbox[colback=popblue,colframe=popblue,coltext=white,arc=9pt,boxsep=4pt,left=6pt,right=6pt]{\bfseries\small Reporting verbs}\end{center}
\vspace{-2pt}
\begin{tcbraster}[raster columns=2,raster equal height=rows,raster column skip=6pt,raster row skip=6pt,enhanced,blankest]
\begin{tcolorbox}[enhanced,colback=white,colframe=poporange,boxrule=0.9pt,arc=3pt,title={+ that (clause)},fonttitle=\bfseries\scriptsize,coltitle=white,colbacktitle=poporange,left=4pt,right=4pt,top=2pt,bottom=2pt,boxsep=1pt,toptitle=1.5pt,bottomtitle=1.5pt,halign=left]
\begin{itemize}[nosep,leftmargin=*,itemsep=1pt,font=\scriptsize]
\item say, tell, admit, claim, insist, promise, agree
\end{itemize}
\end{tcolorbox}
\begin{tcolorbox}[enhanced,colback=white,colframe=popgreen,boxrule=0.9pt,arc=3pt,title={+ to-infinitive},fonttitle=\bfseries\scriptsize,coltitle=white,colbacktitle=popgreen,left=4pt,right=4pt,top=2pt,bottom=2pt,boxsep=1pt,toptitle=1.5pt,bottomtitle=1.5pt,halign=left]
\begin{itemize}[nosep,leftmargin=*,itemsep=1pt,font=\scriptsize]
\item advise, ask, warn, order, remind, invite
\item promise, refuse, agree, offer (sans personne)
\end{itemize}
\end{tcolorbox}
\begin{tcolorbox}[enhanced,colback=white,colframe=poppurple,boxrule=0.9pt,arc=3pt,title={+ -ing},fonttitle=\bfseries\scriptsize,coltitle=white,colbacktitle=poppurple,left=4pt,right=4pt,top=2pt,bottom=2pt,boxsep=1pt,toptitle=1.5pt,bottomtitle=1.5pt,halign=left]
\begin{itemize}[nosep,leftmargin=*,itemsep=1pt,font=\scriptsize]
\item deny, admit, suggest, regret
\end{itemize}
\end{tcolorbox}
\end{tcbraster}

\legende{Carte mentale : les trois constructions des verbes introducteurs.}
\end{popfigure}

\textbf{Verbes + to-infinitive.}\par
Ces verbes servent à rapporter des conseils, des ordres, des invitations, des rappels. La construction est~: verbe + \cle{personne} + \eng{to}-infinitif.

\begin{exemplebox}[Verbe + personne + to-infinitive]
\begin{itemize}
\item \eng{“You should see a doctor,” she said.} → \eng{She advised me to see a doctor.} «~Elle m'a conseillé de voir un médecin.~»
\item \eng{“Don't forget your identity card,” the officer said.} → \eng{The officer reminded me to take my identity card.} «~L'agent m'a rappelé de prendre ma carte d'identité.~»
\item \eng{“Would you like to join our club?”} → \eng{They invited me to join their club.} «~Ils m'ont invité à rejoindre leur club.~»
\item \eng{“Be careful with that dog!”} → \eng{He warned the children to be careful with the dog.} «~Il a averti les enfants de se méfier du chien.~»
\end{itemize}
\end{exemplebox}

Quatre verbes de cette famille se construisent \textbf{sans} complément de personne~: \eng{promise, refuse, agree, offer}~: \eng{He promised to help me.} «~Il a promis de m'aider.~»~; \eng{She refused to answer.} «~Elle a refusé de répondre.~»~; \eng{They agreed to meet on Saturday.} «~Ils ont accepté de se retrouver samedi.~»~; \eng{He offered to carry my bag.} «~Il a proposé de porter mon sac.~»

\textbf{Verbes + -ing.}\par
Ces verbes rapportent des aveux, des dénis, des suggestions, des regrets. La construction est~: verbe + \eng{-ing} (sans complément de personne).

\begin{exemplebox}[Verbe + -ing]
\begin{itemize}
\item \eng{“Yes, I stole the money,” he said.} → \eng{He admitted stealing the money.} «~Il a reconnu avoir volé l'argent.~»
\item \eng{“I didn't take the money,” he said.} → \eng{He denied taking the money.} «~“Je n'ai pas pris l'argent”, a-t-il dit. → Il a nié avoir pris l'argent.~»
\item \eng{“Why don't we vote early?”} → \eng{She suggested voting early.} «~Elle a suggéré de voter tôt.~»
\item \eng{“I am sorry I shouted.”} → \eng{He regretted shouting.} «~Il a regretté d'avoir crié.~»
\end{itemize}
\end{exemplebox}

Deux verbes à prépositions très fréquents~: \eng{accuse somebody of -ing} et \eng{apologise for -ing}~: \eng{They accused him of cheating.} «~Ils l'ont accusé d'avoir triché.~»~; \eng{She apologised for being late.} «~Elle s'est excusée d'être en retard.~»

\textbf{Verbes + that + should.}\par
Après \eng{suggest}, \eng{recommend} et \eng{demand}, on trouve une construction formelle~: \eng{that} + sujet + \eng{should} + base verbale~: \eng{He suggested that we should vote early.} «~Il a suggéré que nous votions tôt.~»~; \eng{The committee recommended that the law should be changed.} «~Le comité a recommandé que la loi soit modifiée.~» Remarquez que le français utilise ici le \emph{subjonctif}, absent de l'anglais, qui le remplace par \eng{should}.

\courssub{Tableau récapitulatif des constructions}

\begin{center}
\begin{tabularx}{\linewidth}{|l|X|X|}
\hline
\rowcolor{popblueL} \textbf{Construction} & \textbf{Verbes} & \textbf{Exemple} \\
\hline
+ (that) + phrase & say, tell, admit, claim, insist, announce & \eng{He claimed (that) he was innocent.} \\
\hline
+ personne + to-inf. & advise, ask, encourage, invite, order, remind, tell, warn & \eng{She warned me not to go out.} \\
\hline
+ to-inf. (sans personne) & promise, refuse, agree, offer & \eng{He offered to pay for the taxi.} \\
\hline
+ -ing & admit, deny, suggest, regret & \eng{She denied breaking the window.} \\
\hline
+ that + should & suggest, recommend, demand & \eng{They demanded that he should resign.} \\
\hline
+ prép. + -ing & accuse (of), apologise (for) & \eng{He apologised for lying.} \\
\hline
\end{tabularx}
\end{center}

\courssub{Comparaison avec le français}

\begin{center}
\begin{tabularx}{\linewidth}{|X|X|X|}
\hline
\rowcolor{popblueL} \textbf{Français} & \textbf{English} & \textbf{Remarque} \\
\hline
Il m'a dit de venir. & \eng{He told me to come.} & Jamais \eng{He said me to come}. \\
\hline
Il a dit qu'il viendrait. & \eng{He said he would come.} & Le conditionnel français = \eng{would}. \\
\hline
Il a dit qu'il viendra. & \eng{He said he will come.} & Style direct conservé : futur possible si le fait est encore à venir. \\
\hline
Il a nié avoir pris l'argent. & \eng{He denied taking the money.} & Infinitif passé français = \eng{-ing} anglais. \\
\hline
Il a suggéré que nous votions tôt. & \eng{He suggested that we should vote early.} & Subjonctif français = \eng{should} + base verbale. \\
\hline
Elle m'a accusé d'avoir triché. & \eng{She accused me of cheating.} & Préposition \eng{of} obligatoire. \\
\hline
\end{tabularx}
\end{center}

Deux différences majeures méritent d'être soulignées.

Premièrement, le \cle{conditionnel}~: le français garde le futur du passé après «~dire~» au passé («~il a dit qu'il viendrait~»). L'anglais utilise \eng{would}, qui ressemble formellement à notre conditionnel mais n'en est pas l'équivalent général~: \eng{would} sert ici de «~futur dans le passé~», conséquence du recul des temps (\eng{will} → \eng{would}). Ne traduisez donc jamais un conditionnel français de politesse («~je voudrais~») par cette logique~: «~je voudrais~» se dit \eng{I would like}, pas \eng{I will like}.

Deuxièmement, la \cle{subordonnée}~: le français exige «~que~» («~il a dit \textbf{qu'}il viendrait~»), tandis que l'anglais peut omettre \eng{that}, surtout à l'oral~: \eng{He said (that) he would come.} À l'écrit du Bac, gardez \eng{that} dans les phrases longues pour plus de clarté.

\begin{attention}[Pièges récapitulatifs des francophones]
\begin{itemize}
\item \eng{He said me to come.} → \eng{He told me to come.} / \eng{He asked me to come.}
\item \eng{He told that he was ill.} → \eng{He told me that he was ill.} (personne obligatoire après \eng{tell})
\item \eng{She suggested me to leave.} → \eng{She suggested that I (should) leave.} ou \eng{She suggested leaving.} (\eng{suggest} ne prend jamais de complément de personne + infinitif)
\item \eng{He denied to take the money.} → \eng{He denied taking the money.} (\eng{deny} + \eng{-ing})
\item \eng{They accused him to cheat.} → \eng{They accused him of cheating.}
\end{itemize}
\end{attention}

\begin{methode}[Choisir le bon verbe introducteur]
\begin{enumerate}
\item Lisez la phrase au style direct et demandez-vous~: quelle est l'\textbf{intention} du locuteur~?
\item Informer simplement~? → \eng{say / tell}. Interroger~? → \eng{ask}. Ordonner~? → \eng{order / tell}. Conseiller~? → \eng{advise}. Promettre~? → \eng{promise}. Nier~? → \eng{deny}. Reconnaître~? → \eng{admit}. Proposer~? → \eng{suggest / offer}. Accuser~? → \eng{accuse}.
\item Appliquez la \textbf{construction} propre au verbe choisi (\eng{that}, \eng{to}-infinitif, \eng{-ing}, préposition).
\item N'oubliez pas le \textbf{recul des temps} si le verbe introducteur est au passé.
\item Relisez~: la phrase rapportée doit être plus courte et plus précise que la version avec \eng{say}.
\end{enumerate}
\end{methode}

\begin{exoresolu}[Transformez avec le verbe entre parenthèses]
Énoncé. Rapporter chaque phrase avec le verbe indiqué~:
1. \eng{“I didn't break the window,” the boy said.} (deny)
2. \eng{“You should apologise to your sister,” the mother said.} (advise)
3. \eng{“I will help you with your homework,” she said.} (promise)
4. \eng{“Why don't we plant trees in the school yard?”} (suggest)
5. \eng{“You stole my phone!” she said to him.} (accuse)
\tcblower
Solution détaillée.
1. \eng{The boy denied breaking the window.} — \eng{deny} + \eng{-ing}~; la négation du style direct (\eng{didn't break}) devient le déni exprimé par \eng{deny} lui-même.
2. \eng{The mother advised her child to apologise to his sister.} — \eng{advise} + personne + \eng{to}-infinitif.
3. \eng{She promised to help me with my homework.} — \eng{promise} + \eng{to}-infinitif, sans complément de personne~; \eng{will} disparaît car la promesse est contenue dans le verbe.
4. \eng{She suggested planting trees in the school yard.} (ou \eng{She suggested that we should plant trees in the school yard.}) — les deux constructions sont correctes.
5. \eng{She accused him of stealing her phone.} — \eng{accuse} + personne + \eng{of} + \eng{-ing}.
\end{exoresolu}

\begin{exercice}[À vous de jouer]
Rapportez avec le verbe le plus précis~:
1. \eng{“Don't walk on the grass,” the guard said to us.} (warn)
2. \eng{“Yes, I took your pen,” my brother said.} (admit)
3. \eng{“I will never lie to you again,” he said.} (promise)
4. \eng{“Please, please lend me your bike,” she said.} (beg)
5. \eng{“I'm sorry I forgot your birthday,” he said.} (apologise)
\end{exercice}

\begin{corrige}[Exercice : À vous de jouer]
1. \eng{The guard warned us not to walk on the grass.}
2. \eng{My brother admitted taking my pen.}
3. \eng{He promised never to lie to me again.}
4. \eng{She begged me to lend her my bike.}
5. \eng{He apologised for forgetting my birthday.}
\end{corrige}

\begin{aretenir}[L'essentiel de la section]
\begin{itemize}
\item \eng{Say (to somebody) that...} mais \eng{tell somebody that...} — jamais l'inverse.
\item Trois constructions à maîtriser~: + \eng{that}, + personne + \eng{to}-infinitif, + \eng{-ing}.
\item \eng{Suggest, recommend, demand} + \eng{that} + \eng{should} + base verbale (= subjonctif français).
\item Le conditionnel français après «~dire~» se rend par \eng{would}~; «~que~» peut être omis en anglais.
\item Varier les verbes introducteurs = écriture plus riche et note plus élevée.
\end{itemize}
\end{aretenir}

\begin{savaistu}[Le verbe introducteur, arme du journaliste]
Les verbes introducteurs précis sont l'outil quotidien du journaliste d'investigation anglophone. Quand un reporter écrit \eng{The minister claimed that...} («~le ministre a affirmé que...~»), il prend ses distances~; \eng{The minister alleged that...} suggère une affirmation non prouvée~; \eng{The minister admitted that...} signale un aveu~; \eng{The minister denied that...} rapporte un démenti. Le choix du verbe oriente le jugement du lecteur~! Au Bac, dans un reportage ou un article, utiliser \eng{claim, allege, admit, deny, insist} au lieu du seul \eng{say} démontre une maîtrise avancée du lexique — et les correcteurs le remarquent toujours.
\end{savaistu}

\coursec{Entraînement guidé et pièges récapitulatifs}

Nous savons désormais choisir le verbe introducteur adapté (\eng{say}, \eng{tell}, \eng{ask}, \eng{warn}, \eng{promise}, \eng{deny}\dots) et nous avons comparé les constructions anglaises et françaises. Il est temps de mettre tout cela en pratique. Cette dernière section est un véritable entraînement au format de l'examen : nous allons transformer des déclarations, des questions, des ordres, puis un dialogue complet, avant de dresser le bilan des cinq pièges qui coûtent le plus de points aux candidats francophones.

\courssub{La méthode R-B-P-M : quatre vérifications avant de valider}

Avant tout exercice, mémorisons une démarche systématique. Chaque transformation du discours direct au discours indirect exige quatre contrôles successifs, dans l'ordre.

\begin{methode}[La checklist R-B-P-M]
Avant de valider une transformation au Bac, vérifie les quatre étapes suivantes :
\begin{enumerate}
\item \cle{R} — \textbf{Reporting verb} : le verbe introducteur est-il correct ? (\eng{say} sans complément de personne, \eng{tell} + complément de personne, \eng{ask} pour les questions et les demandes, \eng{warn}, \eng{promise}, \eng{deny}\dots)
\item \cle{B} — \textbf{Backshift} : le temps du verbe rapporté a-t-il reculé ? (présent → prétérit, prétérit → past perfect, \eng{will} → \eng{would}, \eng{can} → \eng{could}\dots)
\item \cle{P} — \textbf{Pronouns} : les pronoms personnels et possessifs sont-ils adaptés au nouveau locuteur ? (\eng{I} → \eng{he/she}, \eng{my} → \eng{his/her}\dots)
\item \cle{M} — \textbf{Markers} : les marqueurs de temps et de lieu sont-ils transformés ? (\eng{today} → \eng{that day}, \eng{here} → \eng{there}, \eng{tomorrow} → \eng{the next day}\dots)
\end{enumerate}
Si une seule de ces quatre étapes est négligée, la phrase est fautive. Au Bac, relis chaque phrase transformée en te posant les quatre questions R-B-P-M.
\end{methode}

\begin{popfigure}
\begin{tikzpicture}[node distance=0.45cm]
\node[draw=popblue, fill=popblueL, rounded corners, minimum width=2.6cm, minimum height=1.1cm, align=center] (r) {\textbf{1. R}\\Reporting verb};
\node[draw=poporange, fill=poporangeL, rounded corners, minimum width=2.6cm, minimum height=1.1cm, align=center, right=of r] (b) {\textbf{2. B}\\Backshift};
\node[draw=popgreen, fill=popgreenL, rounded corners, minimum width=2.6cm, minimum height=1.1cm, align=center, right=of b] (p) {\textbf{3. P}\\Pronouns};
\node[draw=poppurple, fill=poppurpleL, rounded corners, minimum width=2.6cm, minimum height=1.1cm, align=center, right=of p] (m) {\textbf{4. M}\\Markers};
\draw[-{Stealth}, thick, popdark] (r) -- (b);
\draw[-{Stealth}, thick, popdark] (b) -- (p);
\draw[-{Stealth}, thick, popdark] (p) -- (m);
\end{tikzpicture}
\legende{La checklist visuelle R-B-P-M : quatre étapes à vérifier dans l'ordre pour chaque transformation.}
\end{popfigure}

\begin{exemplebox}[Un exemple complet, étape par étape]
Phrase directe : \eng{“Will you defend human rights?” the activist asked the crowd.} (« Défendrez-vous les droits de l'homme ? » a demandé le militant à la foule.)
\begin{itemize}
\item \textbf{R} : question fermée → \eng{asked} + \eng{if}.
\item \textbf{B} : \eng{will defend} → \eng{would defend}.
\item \textbf{P} : \eng{you} → \eng{they} (la foule).
\item \textbf{M} : pas de marqueur temporel ici ; suppression du point d'interrogation.
\end{itemize}
Phrase rapportée : \eng{The activist asked the crowd if they would defend human rights.} (Le militant a demandé à la foule si elle défendrait les droits de l'homme.)
\end{exemplebox}

\courssub{Exercice 1 : déclarations (backshift simple)}

\begin{exercice}[Exercice 1 — Statements]
Transforme les déclarations suivantes au discours indirect. Commence chaque phrase par l'indication donnée entre parenthèses.
\begin{enumerate}
\item \eng{“I am proud of my country,” the student said.} (The student said\dots)
\item \eng{“We will vote in the next election,” the young people said.} (The young people said\dots)
\item \eng{“I have never seen such a crowd,” the journalist said.} (The journalist said\dots)
\item \eng{“The minister is speaking now,” the reporter said.} (The reporter said\dots)
\item \eng{“I can help you with the form,” the clerk told me.} (The clerk told me\dots)
\item \eng{“We planted trees yesterday,” the pupils said.} (The pupils said\dots)
\end{enumerate}
\end{exercice}

\begin{corrige}[Exercice 1]
\begin{enumerate}
\item \eng{The student said (that) he was proud of his country.} (présent → prétérit ; \eng{I} → \eng{he}, \eng{my} → \eng{his})
\item \eng{The young people said (that) they would vote in the next election.} (\eng{will} → \eng{would})
\item \eng{The journalist said (that) he had never seen such a crowd.} (present perfect → past perfect)
\item \eng{The reporter said (that) the minister was speaking then.} (présent continu → prétérit continu ; \eng{now} → \eng{then})
\item \eng{The clerk told me (that) he could help me with the form.} (\eng{can} → \eng{could})
\item \eng{The pupils said (that) they had planted trees the day before.} (prétérit → past perfect ; \eng{yesterday} → \eng{the day before})
\end{enumerate}
\end{corrige}

\courssub{Exercice 2 : questions fermées et ouvertes}

\begin{exercice}[Exercice 2 — Questions]
Rapporte les questions suivantes. Utilise \eng{if} pour les questions fermées et le mot interrogatif pour les questions ouvertes. Attention à l'ordre des mots : sujet + verbe, sans inversion.
\begin{enumerate}
\item \eng{“Do you know your rights?” the teacher asked us.} (fermée)
\item \eng{“Have you registered to vote?” the official asked the young man.} (fermée)
\item \eng{“Will the meeting start on time?” the journalist asked.} (fermée)
\item \eng{“Where does the mayor live?” the visitor asked.} (ouverte)
\item \eng{“Why did the police stop the march?” the reporter asked the organisers.} (ouverte)
\end{enumerate}
\end{exercice}

\begin{corrige}[Exercice 2]
\begin{enumerate}
\item \eng{The teacher asked us if we knew our rights.} (\eng{do you know} → \eng{we knew} : pas de \eng{did}, pas d'inversion)
\item \eng{The official asked the young man if he had registered to vote.} (present perfect → past perfect)
\item \eng{The journalist asked if the meeting would start on time.} (\eng{will} → \eng{would})
\item \eng{The visitor asked where the mayor lived.} (\eng{does\dots live} → \eng{lived} : ordre sujet + verbe)
\item \eng{The reporter asked the organisers why the police had stopped the march.} (prétérit → past perfect)
\end{enumerate}
\end{corrige}

\begin{attention}[Le piège de l'inversion]
En français, on dit : « Il a demandé où j'allais » — jamais « où allais-je ». L'anglais fait de même dans les questions rapportées : on retrouve l'ordre de la déclaration. \eng{He asked where did I go.} ✗ → \eng{He asked where I went.} ✓. Les auxiliaires \eng{do}, \eng{does}, \eng{did} disparaissent totalement.
\end{attention}

\courssub{Exercice 3 : ordres et demandes}

\begin{exercice}[Exercice 3 — Commands and requests]
Rapporte les ordres et demandes suivants avec la structure \eng{tell/ask + complément + to-infinitive} (ou \eng{not to} pour la négation).
\begin{enumerate}
\item \eng{“Open your books,” the teacher told the class.} (ordre affirmatif)
\item \eng{“Please help me carry this box,” the old woman asked the boy.} (demande)
\item \eng{“Don't touch the documents,” the officer told the visitors.} (ordre négatif)
\item \eng{“Don't be late for the ceremony,” the headmaster warned the students.} (avertissement négatif)
\end{enumerate}
\end{exercice}

\begin{corrige}[Exercice 3]
\begin{enumerate}
\item \eng{The teacher told the class to open their books.}
\item \eng{The old woman asked the boy to help her carry the box.} (\eng{me} → \eng{her})
\item \eng{The officer told the visitors not to touch the documents.} (négation : \eng{not to} + base verbale)
\item \eng{The headmaster warned the students not to be late for the ceremony.}
\end{enumerate}
\end{corrige}

\courssub{Exercice 4 : choisir le bon verbe introducteur}

\begin{exoresolu}[Exercice 4 — Say, tell, ask, warn, promise ou deny ?]
Complète chaque phrase avec le verbe introducteur qui convient, conjugué au prétérit : \eng{say}, \eng{tell}, \eng{ask}, \eng{warn}, \eng{promise}, \eng{deny}.
\begin{enumerate}
\item The candidate \dots\ that he would reduce unemployment.
\item The guide \dots\ the tourists that the road was dangerous at night.
\item The accused man \dots\ that he had stolen the money.
\item My father \dots\ me that hard work always pays.
\item She \dots\ that she was tired, but she didn't \dots\ me why.
\end{enumerate}
\tcblower
\textbf{Solution détaillée.}
\begin{enumerate}
\item \eng{The candidate \textbf{promised} that he would reduce unemployment.} (un engagement pour l'avenir → \eng{promise})
\item \eng{The guide \textbf{warned} the tourists that the road was dangerous at night.} (un danger → \eng{warn} + complément de personne)
\item \eng{The accused man \textbf{denied} that he had stolen the money.} (il affirme le contraire → \eng{deny} ; attention, \eng{deny} ne prend jamais de complément de personne)
\item \eng{My father \textbf{told} me that hard work always pays.} (complément de personne \eng{me} → \eng{tell}, jamais \eng{said me})
\item \eng{She \textbf{said} that she was tired, but she didn't \textbf{tell} me why.} (premier verbe sans complément de personne → \eng{say} ; second avec \eng{me} → \eng{tell})
\end{enumerate}
\end{exoresolu}

\courssub{Exercice 5 : corrige les erreurs typiques des francophones}

\begin{exercice}[Exercice 5 — Error correction]
Chaque phrase contient une erreur classique. Souligne l'erreur et réécris la phrase correctement.
\begin{enumerate}
\item \eng{She asked me where did I live.}
\item \eng{He said me that he was tired.}
\item \eng{The boy said he will come to the meeting.}
\item \eng{My mother told me that she was cooking today.}
\item \eng{The teacher asked me what did I want?}
\end{enumerate}
\end{exercice}

\begin{corrige}[Exercice 5]
\begin{enumerate}
\item \eng{She asked me where I \textbf{lived}.} (pas d'inversion, pas de \eng{did}, backshift)
\item \eng{He \textbf{told} me that he was tired.} (\eng{say} ne prend jamais de complément de personne)
\item \eng{The boy said he \textbf{would} come to the meeting.} (backshift : \eng{will} → \eng{would})
\item \eng{My mother told me that she was cooking \textbf{that day}.} (\eng{today} doit être transformé)
\item \eng{The teacher asked me what I \textbf{wanted}.} (ordre sujet + verbe, et suppression du point d'interrogation : une question rapportée est une déclaration)
\end{enumerate}
\end{corrige}

\courssub{Exercice 6 (type Bac) : rapporter un dialogue complet}

À l'épreuve écrite, on te demandera souvent de rapporter un court dialogue en un seul paragraphe cohérent. Voici d'abord un modèle rédigé, puis ta tâche.

\begin{textebox}[Modèle : reporting a headmaster's announcement]
The headmaster announced that the school would celebrate Human Rights Day the following week. He told the students to prepare speeches and asked them if they were ready. He warned them not to be late and promised that the best speaker would receive a prize.
\end{textebox}

Observons la mécanique du modèle : quatre verbes introducteurs variés (\eng{announced}, \eng{told}, \eng{asked}, \eng{warned}, \eng{promised}), un backshift systématique (\eng{would celebrate}, \eng{were}, \eng{would receive}), un marqueur temporel transformé (\eng{the following week}) et un ordre négatif rapporté avec \eng{not to be late}. C'est exactement ce niveau de précision que le correcteur attend.

\begin{exercice}[Exercice 6 — Type Bac : report the dialogue]
Lis le dialogue suivant, puis rapporte-le en un paragraphe de 4 à 5 phrases au discours indirect. Varie les verbes introducteurs.

\eng{Journalist: “Madam Mayor, why did you create this youth centre?”}

\eng{Mayor: “I created it because young people need a place to meet.”}

\eng{Journalist: “Will you open another centre next year?”}

\eng{Mayor: “Yes, we will. Don't forget to tell your readers that every citizen matters.”}
\end{exercice}

\begin{corrige}[Exercice 6 — Corrigé type]
\eng{The journalist asked the mayor why she had created the youth centre. The mayor explained that she had created it because young people needed a place to meet. The journalist then asked her if she would open another centre the following year. The mayor replied that they would, and she asked the journalist not to forget to tell his readers that every citizen mattered.}

Points notés : backshift complet (\eng{had created}, \eng{needed}, \eng{would open}, \eng{mattered}), question ouverte sans inversion (\eng{why she had created}), question fermée avec \eng{if}, marqueur temporel transformé (\eng{the following year}), ordre négatif rapporté avec \eng{not to forget}, et verbes introducteurs variés (\eng{asked}, \eng{explained}, \eng{replied}).
\end{corrige}

\courssub{Bilan : les 5 pièges à éliminer définitivement}

\begin{attention}[Récapitulatif des erreurs fatales]
\begin{enumerate}
\item \textbf{Inversion conservée} : \eng{He asked where did I go.} ✗ → \eng{He asked where I went.} ✓
\item \textbf{Confusion say/tell} : \eng{She said me the truth.} ✗ → \eng{She told me the truth.} ✓ (\eng{say} = sans complément de personne ; \eng{tell} = toujours avec)
\item \textbf{Oubli du backshift} : \eng{He said he will come.} ✗ → \eng{He said he would come.} ✓
\item \textbf{Marqueurs non transformés} : \eng{She said she was busy today.} ✗ → \eng{She said she was busy that day.} ✓
\item \textbf{Point d'interrogation conservé} : une question rapportée est une affirmation ; elle se termine par un point.
\end{enumerate}
\end{attention}

\begin{popfigure}
\begin{tikzpicture}
\node[draw=popdark, fill=popcream, rounded corners, align=left, text width=13.5cm, inner sep=8pt] {
\textbf{\textcolor{popblue}{Piège 1 — Inversion}} \quad ✗ \eng{where did I go} \quad ✓ \eng{where I went}\\[4pt]
\textbf{\textcolor{popblue}{Piège 2 — say/tell}} \quad ✗ \eng{She said me} \quad ✓ \eng{She told me}\\[4pt]
\textbf{\textcolor{popblue}{Piège 3 — Backshift}} \quad ✗ \eng{he will come} \quad ✓ \eng{he would come}\\[4pt]
\textbf{\textcolor{popblue}{Piège 4 — Marqueurs}} \quad ✗ \eng{today} \quad ✓ \eng{that day}\\[4pt]
\textbf{\textcolor{popblue}{Piège 5 — Ponctuation}} \quad ✗ \eng{\dots what I wanted?} \quad ✓ \eng{\dots what I wanted.}
};
\end{tikzpicture}
\legende{Tableau récapitulatif « Les 5 pièges » : mauvaise réponse contre bonne réponse.}
\end{popfigure}

\begin{center}
\begin{tabularx}{\linewidth}{|X|X|X|}
\hline
\rowcolor{popblueL} \textbf{Piège} & \textbf{Français (calque)} & \textbf{Anglais correct} \\
\hline
Inversion & « Il a demandé où j'allais » & \eng{He asked where I went.} \\
\hline
say/tell & « Elle m'a dit » (dire + moi) & \eng{She told me.} (jamais \eng{said me}) \\
\hline
Backshift & « Il a dit qu'il viendrait » & \eng{He said he would come.} \\
\hline
Marqueurs & « aujourd'hui » → ce jour-là & \eng{today} → \eng{that day} \\
\hline
Ponctuation & point final, pas de « ? » & \eng{She asked what I wanted.} \\
\hline
\end{tabularx}
\end{center}

\begin{savaistu}[Le discours indirect au Bac Tle A]
À l'épreuve écrite d'anglais du Baccalauréat Tle A, les questions de transformation du discours direct au discours indirect rapportent souvent de 2 à 4 points dans la section grammaire. Le critère discriminant — celui qui sépare les bonnes copies des excellentes — est la précision sur les marqueurs spatio-temporels : presque tous les candidats pensent au backshift, mais beaucoup oublient de transformer \eng{today} en \eng{that day} ou \eng{tomorrow} en \eng{the next day}. Applique la checklist R-B-P-M à chaque phrase, et ces points seront les tiens.
\end{savaistu}

\begin{aretenir}[L'essentiel de la leçon]
Rapporter les paroles de quelqu'un, c'est effectuer quatre opérations simultanées : choisir le bon \cle{verbe introducteur}, faire \cle{reculer les temps}, adapter les \cle{pronoms} et transformer les \cle{marqueurs} de temps et de lieu. Les questions rapportées perdent leur inversion et leur point d'interrogation ; les ordres se rapportent avec \eng{tell/ask + complément + (not) to-infinitive}. La méthode R-B-P-M est ton filet de sécurité : vérifie les quatre étapes, et le discours indirect n'aura plus aucun secret pour toi.
\end{aretenir}

\newpage

\coursec{Activité d'intégration}

\coursec{Lesson 38 — Grammar: Revision: Direct and Indirect Speech}

\courssub{Learning Objectives}
À l'issue de cette leçon, tu seras capable de :
\begin{itemize}
\item identifier la structure du discours indirect (\emph{reported speech}) et la distinguer du discours direct ;
\item appliquer systématiquement le \cle{backshift} (recul des temps), les changements de pronoms, possessifs et marqueurs spatio-temporels ;
\item transformer correctement les déclarations, les questions (fermées et \emph{wh-}), les ordres, les interdictions et les suggestions ;
\item choisir le verbe introducteur (\emph{reporting verb}) le plus précis selon la fonction illocutoire (promesse, ordre, avertissement, déni, etc.) ;
\item rédiger un compte rendu journalistique ou un article de presse en anglais formel, en respectant les conventions du discours rapporté ;
\item t'auto-évaluer à l'aide de la grille \cle{R-B-P-M} (\emph{Reporting verb — Backshift — Pronouns — Markers}).
\end{itemize}

\courssub{Context \& Observation : A Press Conference in Bamenda}

\begin{textebox}[The Mayor's Press Conference — Direct Speech Transcript]
\textbf{The Mayor:} “Good morning, everyone. I am delighted to welcome you here today. Our city council has just launched a new programme to protect the rights of every citizen, and we will open a free legal advice centre next month.”

\textbf{The Mayor:} “Every resident must register at the district office before Friday. Do not pay anyone for this service: it is completely free.”

\textbf{Journalist 1:} “Mr Mayor, will the new centre also help young people who cannot find a job?”

\textbf{The Mayor:} “Yes, it will. We are going to create a special desk for young job seekers, because unemployment is the biggest challenge our youth faces today.”

\textbf{Journalist 2:} “Some citizens complained last week that the police had ignored their reports. What are you doing about that?”

\textbf{The Mayor:} “I did not receive any official complaint last week, but I promise you now: I will personally examine every case that is brought to my office.”
\end{textebox}

\begin{savaistu}[Press Conferences in the Anglophone Context]
Dans la tradition démocratique anglophone (\emph{accountability}), la conférence de presse (\emph{press conference}) permet aux élus de rendre des comptes et aux journalistes d'exercer leur rôle de \emph{watchdog} (chien de garde). Au Cameroun anglophone (Bamenda, Buea, Limbe), les médias comme \emph{The Guardian Post}, \emph{The National Times} ou les radios (CBS Radio, Christian Broadcasting Service) rapportent ces événements au discours indirect. Maîtriser cette compétence, c'est accéder à l'information citoyenne et réussir l'épreuve de \emph{Composition} ou de \emph{Grammar} au Baccalauréat.
\end{savaistu}

\courssub{Grammar Focus: The Mechanics of Reported Speech}

\courssub{The Backshift (Le recul des temps)}

\begin{propriete}[Règle du Backshift]
Lorsque le verbe introducteur (\emph{reporting verb}) est au passé (\emph{said, told, asked, promised, denied...}), les temps de la proposition rapportée reculent d'un cran vers le passé. Ce phénomène s'appelle le \cle{backshift}.
\end{propriete}

\begin{center}
\begin{tabularx}{\linewidth}{|X|X|X|}
\hline
\rowcolor{popblueL} \textbf{Direct Speech (Tense)} & \textbf{Reported Speech (Tense)} & \textbf{Exemple} \\
\hline
Present Simple (\emph{I work}) & Past Simple (\emph{he worked}) & \eng{“I live in Buea.”} $\rightarrow$ He said he \textbf{lived} in Buea. \\
\hline
Present Continuous (\emph{I am working}) & Past Continuous (\emph{he was working}) & \eng{“I am reading.”} $\rightarrow$ She said she \textbf{was reading}. \\
\hline
Present Perfect (\emph{I have worked}) & Past Perfect (\emph{he had worked}) & \eng{“I have finished.”} $\rightarrow$ He said he \textbf{had finished}. \\
\hline
Past Simple (\emph{I worked}) & Past Perfect (\emph{he had worked}) & \eng{“I saw him.”} $\rightarrow$ She said she \textbf{had seen} him. \\
\hline
Past Continuous (\emph{I was working}) & Past Perfect Continuous (\emph{he had been working}) & \eng{“I was sleeping.”} $\rightarrow$ He said he \textbf{had been sleeping}. \\
\hline
\emph{will} (\emph{I will go}) & \emph{would} (\emph{he would go}) & \eng{“I will call you.”} $\rightarrow$ She said she \textbf{would} call me. \\
\hline
\emph{can} (\emph{I can swim}) & \emph{could} (\emph{he could swim}) & \eng{“I can help.”} $\rightarrow$ He said he \textbf{could} help. \\
\hline
\emph{may} (\emph{It may rain}) & \emph{might} (\emph{it might rain}) & \eng{“It may rain.”} $\rightarrow$ He said it \textbf{might} rain. \\
\hline
\emph{must} (obligation) (\emph{I must go}) & \emph{had to} (\emph{he had to go}) & \eng{“I must leave.”} $\rightarrow$ She said she \textbf{had to} leave. \\
\hline
\emph{shall} (\emph{Shall I open it?}) & \emph{should} (\emph{he should open it}) & \eng{“Shall I open it?”} $\rightarrow$ He asked if he \textbf{should} open it. \\
\hline
\end{tabularx}
\end{center}

\begin{attention}[Exceptions au Backshift — Pas de recul nécessaire]
Le backshift n'est \textbf{pas} obligatoire (mais reste correct) si :
\begin{enumerate}
\item Le verbe introducteur est au \textbf{présent} (\emph{says, tells, asks}) : \eng{He \textbf{says} he \textbf{is} tired.}
\item On rapporte une \textbf{vérité générale}, une loi scientifique ou un fait permanent** : \eng{The teacher said the Earth \textbf{revolves} around the Sun.}
\item Le discours est rapporté \textbf{immédiatement} et la situation n'a pas changé (marqueurs identiques) : \eng{“I \textbf{am} hungry,” she says.} $\rightarrow$ She says she \textbf{is} hungry.
\end{enumerate}
\emph{Au Bac, applique toujours le backshift si le verbe introducteur est au prétérit, sauf consigne contraire.}
\end{attention}

\courssub{Pronouns, Possessives \& Reflexives}

\begin{propriete}[Adaptation des personnes]
Les pronoms changent pour adopter le point de vue du rapporteur (\emph{reporter}), pas celui du locuteur initial.
\end{propriete}

\begin{center}
\begin{tabularx}{\linewidth}{|X|X|X|}
\hline
\rowcolor{popblueL} \textbf{Direct Speech} & \textbf{Reported Speech} & \textbf{Remarque} \\
\hline
I, me, my, mine, myself & he/she, him/her, his/her, his/hers, himself/herself & Selon le genre du locuteur \\
\hline
we, us, our, ours, ourselves & they, them, their, theirs, themselves & \\
\hline
you, your, yours, yourself/yourselves & I/me/my/mine/myself \textbf{OU} he/she/them... & Dépend de qui \emph{you} désigne \\
\hline
this / these & that / those & Démonstratifs \\
\hline
\end{tabularx}
\end{center}

\eng{“I will bring \textbf{my} friends,” \textbf{he} said.} $\rightarrow$ He said \textbf{he} would bring \textbf{his} friends. \\
\eng{“\textbf{You} must sign \textbf{your} form,” the clerk told \textbf{me}.} $\rightarrow$ The clerk told \textbf{me} I had to sign \textbf{my} form.

\courssub{Spatio-Temporal Markers (Marqueurs d'espace et de temps)}

\begin{propriete}[Décentrage spatio-temporel]
Les mots ancrés dans le \emph{ici-et-maintenant} du locuteur initial doivent être décentrés par rapport au moment du rapport.
\end{propriete}

\begin{center}
\begin{tabularx}{\linewidth}{|X|X|}
\hline
\rowcolor{popblueL} \textbf{Direct Speech} & \textbf{Reported Speech} \\
\hline
now & then / at that moment \\
\hline
today / tonight & that day / that night \\
\hline
yesterday & the day before / the previous day \\
\hline
tomorrow & the next day / the following day \\
\hline
last week / month / year & the week / month / year before / the previous week... \\
\hline
next week / month / year & the following week / month / year \\
\hline
ago & before / previously \\
\hline
here & there \\
\hline
this (adjectif) & that \\
\hline
\end{tabularx}
\end{center}

\eng{“I saw him \textbf{yesterday} \textbf{here}.”} $\rightarrow$ He said he had seen him \textbf{the day before} \textbf{there}.

\courssub{Reporting Structures by Speech Act Type}

\begin{propriete}[Déclarations (Statements)]
\emph{Subject + reporting verb + (that) + clause (backshifted).}
\begin{itemize}
\item Verbes courants : \emph{said, told, stated, announced, declared, claimed, explained, added, replied, answered, mentioned}.
\item \emph{Told} exige un complément d'objet (personne) : \eng{He \textbf{told me} (that) he was late.} \textbf{Pas} \eng{He said me...}
\end{itemize}
\end{propriete}

\begin{propriete}[Questions fermées (Yes/No Questions)]
\emph{Subject + asked / wanted to know / wondered + if / whether + subject + verb (affirmative order, no auxiliary do/does/did).}
\eng{“Do you like coffee?”} $\rightarrow$ She asked \textbf{if/whether I liked} coffee. \\
\eng{“Will you come?”} $\rightarrow$ He wanted to know \textbf{whether I would} come.
\end{propriete}

\begin{propriete}[Questions ouvertes (Wh- Questions)]
\emph{Subject + asked / wanted to know / inquired + wh-word + subject + verb (affirmative order).}
\eng{“Where do you live?”} $\rightarrow$ He asked \textbf{where I lived}. \\
\eng{“What time does it start?”} $\rightarrow$ She asked \textbf{what time it started}.
\end{propriete}

\begin{attention}[Piège majeur : Ordre des mots]
\emph{Jamais} d'inversion sujet/verbe, \emph{jamais} d'auxiliaire \emph{do/does/did} dans la question indirecte.
\begin{itemize}
\item \textbf{Faux :} \eng{He asked where \textbf{did I live}.} / \eng{He asked where \textbf{lived I}.}
\item \textbf{Vrai :} \eng{He asked where \textbf{I lived}.}
\end{itemize}
\emph{Jamais} de point d'interrogation à la fin d'une phrase rapportée (c'est une phrase déclarative).
\end{attention}

\begin{propriete}[Ordres, Interdictions, Demandes, Conseils (Imperative / Modal structures)]
\emph{Subject + reporting verb + Object + (not) to-infinitive.}
\begin{center}
\begin{tabularx}{\linewidth}{|X|X|X|}
\hline
\rowcolor{popblueL} \textbf{Fonction} & \textbf{Verbes introducteurs} & \textbf{Structure} \\
\hline
Ordre & \emph{ordered, told, commanded, instructed} & \emph{ordered him \textbf{to leave}} \\
\hline
Interdiction & \emph{forbade, banned, warned ... not to} & \emph{warned us \textbf{not to pay}} \\
\hline
Demande / Requête & \emph{asked, begged, requested, urged} & \emph{asked me \textbf{to help}} \\
\hline
Conseil & \emph{advised, recommended, urged} & \emph{advised him \textbf{to see} a doctor} \\
\hline
Suggestion & \emph{suggested, proposed} & \emph{suggested \textbf{going} / suggested \textbf{that he (should) go}} \\
\hline
\end{tabularx}
\end{center}
\eng{“Don't touch that wire!”} $\rightarrow$ He \textbf{warned me not to touch} that wire. \\
\eng{“Please, sit down.”} $\rightarrow$ She \textbf{asked me to sit} down.
\end{propriete}

\courssub{Reporting Verbs: Precision \& Variety (Au-delà de \emph{Say/Tell})}

\begin{aretenir}[Classification fonctionnelle des verbes introducteurs]
Choisir le bon verbe évite \emph{he said that...} répétitif et montre la nuance (register).
\begin{center}
\begin{tabularx}{\linewidth}{|X|X|X|}
\hline
\rowcolor{popblueL} \textbf{Fonction illocutoire} & \textbf{Verbes (Variété formelle)} & \textbf{Structure typique} \\
\hline
Affirmer / Déclarer & \emph{stated, declared, announced, claimed, asserted, maintained} & V + that-clause \\
\hline
Promettre / S'engager & \emph{promised, pledged, vowed, undertook, committed to} & V + to-inf / V + that-clause \\
\hline
Nier / Contester & \emph{denied, refuted, rejected, contradicted, disputed} & V + -ing / V + that-clause \\
\hline
Avertir / Mettre en garde & \emph{warned, cautioned, alerted, admonished} & V + obj + (not) to-inf / V + that-clause \\
\hline
Ordonner / Exiger & \emph{ordered, commanded, directed, instructed, demanded} & V + obj + to-inf / V + that-clause (subjunctive) \\
\hline
Questionner / S'enquérir & \emph{asked, inquired, questioned, wondered, wanted to know} & V + if/wh-clause \\
\hline
Répondre / Réagir & \emph{replied, responded, answered, retorted, acknowledged} & V + that-clause / V + to-inf \\
\hline
Suggérer / Proposer & \emph{suggested, proposed, recommended, put forward} & V + -ing / V + that-clause (should) \\
\hline
Accepter / Refuser & \emph{agreed, consented, accepted / refused, declined, rejected} & V + to-inf \\
\hline
\end{tabularx}
\end{center}
\end{aretenir}

\begin{exemplebox}[Exemples de variation]
\begin{itemize}
\item \eng{“I will definitely build the bridge,” the Minister said.} $\rightarrow$ The Minister \textbf{pledged} to build the bridge. / The Minister \textbf{undertook} that the bridge \textbf{would} be built.
\item \eng{“I didn't steal the money,” the suspect said.} $\rightarrow$ The suspect \textbf{denied stealing} the money. / The suspect \textbf{denied that he had stolen} the money.
\item \eng{“You really should see a doctor,” my friend said.} $\rightarrow$ My friend \textbf{advised me to see} a doctor. / My friend \textbf{recommended that I (should) see} a doctor.
\end{itemize}
\end{exemplebox}

\courssub{Comparative Approach: English vs. French — Frequent Errors}

\begin{attention}[Checklist des erreurs francophones — \cle{R-B-P-M} Étendu]
\begin{enumerate}
\item \textbf{R — Reporting Verb} : \eng{He said me to go} $\rightarrow$ \textbf{Faux}. \eng{He \textbf{told} me to go} / \eng{He \textbf{ordered} me to go}. \emph{Say} ne prend pas de COI personne.
\item \textbf{B — Backshift} : \eng{He said he \textbf{will} come} $\rightarrow$ \textbf{Faux} (si verbe introducteur au passé). \eng{He said he \textbf{would} come}. Double recul : \eng{“I \textbf{did} it”} $\rightarrow$ He said he \textbf{had done} it (pas \emph{did}).
\item \textbf{P — Pronouns} : \eng{The Mayor said \textbf{I} am delighted} $\rightarrow$ \textbf{Faux}. \eng{The Mayor said \textbf{he} was delighted}. Attention à \emph{our} $\rightarrow$ \emph{their}, \emph{my} $\rightarrow$ \emph{his}.
\item \textbf{M — Markers} : \eng{He said he would do it \textbf{tomorrow}} $\rightarrow$ \textbf{Faux}. \eng{He said he would do it \textbf{the next day}}.
\item \textbf{Word Order (Questions)} : \eng{He asked where \textbf{was the station}} $\rightarrow$ \textbf{Faux}. \eng{He asked where \textbf{the station was}}.
\item \textbf{Punctuation} : \eng{He asked where I lived?} $\rightarrow$ \textbf{Faux}. Point final obligatoire.
\item \textbf{\emph{Suggest} + subjonctif/should} : \eng{He suggested me to go} $\rightarrow$ \textbf{Faux}. \eng{He suggested \textbf{going}} / \eng{He suggested \textbf{that I (should) go}}.
\item \textbf{\emph{Recommend/Propose}} : Même structure que \emph{suggest}. Pas de \emph{to-infinitif} direct après le verbe.
\end{enumerate}
\end{attention}

\courssub{Guided Practice: The Press Conference Task (Group Work)}

\begin{methode}[Procédure : De la transcription à l'article de presse]
\begin{enumerate}
\item \textbf{Role-Play (10 min).} Groupes de 4 : 1 Maire, 2 Journalistes, 1 Rédacteur en chef (\emph{Editor}). Le Maire lit le texte ; Journalistes posent les questions. Tous notent : temps, pronoms, marqueurs (\emph{today, here, now, last week, next month, Friday}).
\item \textbf{Grammatical Spotting (5 min).} Surlignez dans le texte : 1 promesse, 1 ordre, 1 interdiction, 1 question fermée, 1 question \emph{wh-}, 1 négation au prétérit. Associez chaque réplique au verbe introducteur précis (\emph{welcomed, announced, ordered, warned, asked, denied, promised}).
\item \textbf{Drafting (15 min).} Rédigez l'article (8--10 lignes). Contraintes : \textbf{Backshift complet}, pronoms/marqueurs adaptés, \textbf{minimum 3 verbes introducteurs différents}, structure journalistique (Headline, Lead, Body, Conclusion).
\item \textbf{Peer Review (10 min).} Échangez avec un autre groupe. Corrigez avec la grille \cle{R-B-P-M} ci-dessous. Annotez en marge.
\item \textbf{Plenary (5 min).} Lecture du meilleur article. Classe identifie verbes et backshifts.
\end{enumerate}
\end{methode}

\begin{methode}[La Grille R-B-P-M — Outil d'auto-correction]
Pour valider une phrase ou un texte au discours indirect, cochez les 4 cases :
\begin{enumerate}
\item \textbf{R — Reporting Verb} : Varié ? Précis ? (\emph{promised} vs \emph{said}, \emph{denied} vs \emph{said}) ? Syntaxe correcte (\emph{told me} vs \emph{said me}) ?
\item \textbf{B — Backshift} : Tous les temps ont-ils reculé ? (\emph{am$\rightarrow$was, will$\rightarrow$would, has launched$\rightarrow$had launched, did not receive$\rightarrow$had not received}) ? Modaux traités (\emph{can$\rightarrow$could, must$\rightarrow$had to}) ?
\item \textbf{P — Pronouns/Determiners} : \emph{I$\rightarrow$he/she, we$\rightarrow$they, my$\rightarrow$his/her, our$\rightarrow$their, this$\rightarrow$that, your$\rightarrow$my/our/their} ?
\item \textbf{M — Markers} : \emph{now$\rightarrow$then, today$\rightarrow$that day, here$\rightarrow$there, yesterday$\rightarrow$the day before, tomorrow$\rightarrow$the next day, last week$\rightarrow$the week before, next month$\rightarrow$the following month, ago$\rightarrow$before} ?
\end{enumerate}
\end{methode}

\courssub{Model Answer \& Detailed Commentary}

\begin{exoresolu}[Model Article: “Mayor Pledges Free Legal Centre for All Citizens”]
Rédige un article de 8 à 10 lignes rapportant la conférence de presse.
\tcblower
\textbf{Model Article:}

\eng{\textbf{Mayor Pledges Free Legal Centre for All Citizens}}

\eng{Speaking at City Hall yesterday morning, the Mayor of Bamenda stated that he was delighted to welcome the press and announced that the city council had just launched a new programme to protect the rights of every citizen. He pledged that a free legal advice centre would open the following month. He ordered every resident to register at the district office before Friday and warned the public not to pay anyone for the service, stressing that it was completely free. Asked whether the centre would also assist unemployed young people, he replied that it would and added that a special desk for young job seekers would be created, as unemployment was the biggest challenge the youth faced. When questioned about alleged police inaction, the Mayor denied having received any official complaint the week before, but undertook to personally examine every case brought to his office.}

\textbf{Commentaire linguistique détaillé :}
\begin{itemize}
\item \textbf{Verbes introducteurs (R)} : \emph{stated, announced, pledged, ordered, warned, stressed, replied, added, denied, undertook} (10 verbes distincts). Choix fonctionnels : \emph{pledged/undertook} pour l'engagement fort ; \emph{denied} pour la contestation ; \emph{warned} pour l'interdiction.
\item \textbf{Backshift (B)} : \emph{am$\rightarrow$was} ; \emph{has launched$\rightarrow$had launched} ; \emph{will open$\rightarrow$would open} ; \emph{is$\rightarrow$was} ; \emph{did not receive$\rightarrow$had not received} (double recul) ; \emph{faces$\rightarrow$faced} ; \emph{promise$\rightarrow$undertook/promised}.
\item \textbf{Pronoms (P)} : \emph{I$\rightarrow$he, our$\rightarrow$the city council's/their, we$\rightarrow$he/the council, you$\rightarrow$the press/the public, your$\rightarrow$their} (dans \emph{your reports} $\rightarrow$ \emph{their reports}).
\item \textbf{Marqueurs (M)} : \emph{today$\rightarrow$yesterday morning} (contexte) ; \emph{here$\rightarrow$at City Hall} (précision) ; \emph{next month$\rightarrow$the following month} ; \emph{last week$\rightarrow$the week before} ; \emph{now$\rightarrow$(omis/intégré dans le verbe \emph{undertook})} ; \emph{Friday} inchangé (date calendaire fixe, mais \emph{before Friday} rapporté tel quel car antérieur au moment du discours).
\item \textbf{Questions rapportées} : Fermée $\rightarrow$ \emph{asked whether... would} ; Ouverte $\rightarrow$ \emph{questioned about... denied having received} (structure \emph{deny + -ing} pour fait passé).
\item \textbf{Registre journalistique} : Participe présent initial (\emph{Speaking at...}), titre au présent simple (\emph{Headline style}), ton neutre, passif possible (\emph{was created}), connecteurs logiques (\emph{as, when, but}).
\end{itemize}
\end{exoresolu}

\courssub{Individual Practice: Sentence Transformation}

\begin{exercice}[Transformation Drill — Bac Style]
Transformez les phrases suivantes en discours indirect. Utilisez le verbe introducteur suggéré entre parenthèses. Appliquez la grille R-B-P-M.
\begin{enumerate}
\item \eng{“I have never visited the new museum,” the tourist said.} (\emph{claimed})
\item \eng{“Don't forget to submit your application before Friday,” the secretary told the candidates.} (\emph{reminded})
\item \eng{“Why didn't you attend the meeting yesterday?” the boss asked me.} (\emph{wanted to know})
\item \eng{“I will pay the fine if I have to,” the driver said.} (\emph{agreed})
\item \eng{“Let's organize a clean-up campaign next Saturday,” the class prefect proposed.} (\emph{suggested})
\item \eng{“The water supply has been cut off since Monday,” the residents complained.} (\emph{reported})
\item \eng{“You must not use your phones during the exam,” the invigilator warned us.} (\emph{forbade})
\item \eng{“Did the delegation arrive safely?” the Minister asked the protocol officer.} (\emph{inquired})
\end{enumerate}
\end{exercice}

\begin{corrige}[Exercice Transformation Drill]
\begin{enumerate}
\item The tourist \textbf{claimed (that) he had never visited} the new museum.
\item The secretary \textbf{reminded the candidates to submit} their application before Friday. (\emph{remind sb to do sth})
\item The boss \textbf{wanted to know why I had not attended} the meeting the day before. (\emph{wh- order, backshift past$\rightarrow$past perfect, yesterday$\rightarrow$the day before})
\item The driver \textbf{agreed to pay} the fine if he had to. (\emph{agree + to-inf} ; \emph{had to} reste \emph{had to} car déjà modal passé / \emph{must} $\rightarrow$ \emph{had to}).
\item The class prefect \textbf{suggested organizing} a clean-up campaign the following Saturday. / \textbf{suggested (that) they (should) organize}... (\emph{suggest + -ing / that-clause} ; \emph{next$\rightarrow$following}).
\item The residents \textbf{reported (that) the water supply had been cut off} since Monday. (\emph{Passive voice backshift: has been$\rightarrow$had been} ; \emph{since Monday} inchangé si point de repère fixe).
\item The invigilator \textbf{forbade us to use} our phones during the exam. / \textbf{warned us not to use}... (\emph{forbid sb to do} / \emph{warn sb not to do} ; \emph{your$\rightarrow$our}).
\item The Minister \textbf{inquired whether the delegation had arrived} safely. (\emph{Closed question$\rightarrow$if/whether} ; \emph{arrived$\rightarrow$had arrived}).
\end{enumerate}
\end{corrige}

\courssub{Summary \& Self-Assessment}

\begin{aretenir}[Synthèse pour le Bac — Direct vs Indirect Speech]
\begin{itemize}
\item \textbf{Structure globale} : [Reporting Verb (Past)] + [Conjonction (\emph{that/if/wh-})] + [Sujet] + [Verbe Backshifté] + [Compléments adaptés].
\item \textbf{Backshift} : Présent $\rightarrow$ Passé ; Passé/Perfect $\rightarrow$ Plus-que-parfait ; \emph{will/can/may/must} $\rightarrow$ \emph{would/could/might/had to}.
\item \textbf{Marqueurs} : \emph{Now/Here/This} $\rightarrow$ \emph{Then/There/That} ; \emph{Today/Yesterday/Tomorrow} $\rightarrow$ \emph{That day/The day before/The next day}.
\item \textbf{Questions} : Ordre affirmatif (Sujet + Verbe), pas de \emph{do/does/did}, pas de \emph{?}.
\item \textbf{Imperatif/Modaux} : \emph{Verb + Object + (not) to-infinitive} (sauf \emph{suggest/recommend/propose} $\rightarrow$ \emph{-ing / that...should}).
\item \textbf{Variété lexicale} : Bannir \emph{said/told} systématique $\rightarrow$ \emph{stated, promised, denied, warned, inquired, pledged, undertook...}
\end{itemize}
\end{aretenir}

\begin{experience}[Activité orale : “Press Briefing Simulation”]
En binôme : L'élève A est un porte-parole du Ministère de la Jeunesse annonçant un nouveau programme \emph{YouthConnekt} (3 mesures, 1 date, 1 lieu). L'élève B est journaliste (3 questions : 1 fermée, 1 \emph{wh-}, 1 sur un problème passé). A répond. B rédige le \emph{lead} (premier paragraphe) de l'article en 5 min. Échange de rôles. Correction croisée R-B-P-M.
\end{experience}

\courssub{Homework / Devoirs}
\begin{enumerate}
\item Relire l'article modèle et apprendre par cœur la liste des 10 verbes introducteurs variés utilisés.
\item Faire l'exercice de transformation (ci-dessus) si non fini en classe.
\item Rédiger un court article (60--80 mots) rapportant une annonce fictive du Principal de ton lycée sur les nouvelles règles de discipline (port de l'uniforme, horaires, téléphone). Utiliser \textbf{5 verbes introducteurs différents} et respecter la grille R-B-P-M. À rendre pour la prochaine séance.
\end{enumerate}

\newpage

\coursec{Méthodes et exercices résolus}

\courssub{Méthodes et exercices résolus}

\begin{propriete}[Tableau de correspondance des temps (Backshift) et marqueurs spatio-temporels]
\textbf{Règle fondamentale :} Lorsque le verbe introducteur est au \textbf{passé} (\eng{said, told, asked, explained, wondered...}), les temps du discours direct reculent d'un cran dans le discours indirect. Cette règle est \textbf{systématique} en anglais standard (sauf vérité générale ou fait encore actuel au moment du rapport).

\begin{center}
\begin{tabularx}{\linewidth}{|X|X|X|}
\hline
\rowcolor{popblueL}
\textbf{Discours Direct} & \textbf{Discours Indirect (Backshift)} & \textbf{Exemple} \\
\hline
Present Simple (\eng{I work}) & Past Simple (\eng{he worked}) & \eng{He said he \textbf{worked} hard.} \\
\hline
Present Continuous (\eng{I am working}) & Past Continuous (\eng{he was working}) & \eng{He said he \textbf{was working}.} \\
\hline
Present Perfect (\eng{I have worked}) & Past Perfect (\eng{he had worked}) & \eng{He said he \textbf{had worked}.} \\
\hline
Past Simple (\eng{I worked}) & Past Perfect (\eng{he had worked}) & \eng{He said he \textbf{had worked}.} \\
\hline
Past Continuous (\eng{I was working}) & Past Perfect Continuous (\eng{he had been working}) & \eng{He said he \textbf{had been working}.} \\
\hline
\eng{will} (\eng{I will go}) & \eng{would} (\eng{he would go}) & \eng{He said he \textbf{would} go.} \\
\hline
\eng{can} (\eng{I can swim}) & \eng{could} (\eng{he could swim}) & \eng{He said he \textbf{could} swim.} \\
\hline
\eng{may} (\eng{It may rain}) & \eng{might} (\eng{it might rain}) & \eng{He said it \textbf{might} rain.} \\
\hline
\eng{must} (\eng{I must go}) & \eng{had to} (\eng{he had to go}) & \eng{He said he \textbf{had to} go.} \\
\hline
\eng{shall} (futur/offre) & \eng{should} / \eng{would} & \eng{He said he \textbf{should/would} help.} \\
\hline
\end{tabularx}
\end{center}

\textbf{Exceptions (Pas de backshift) :}
\begin{itemize}
    \item \textbf{Vérité générale / scientifique :} \eng{The teacher said water \textbf{boils} at 100°C.}
    \item \textbf{Fait encore vrai au moment du rapport :} \eng{He said he \textbf{loves} his job.} (Il l'aime toujours).
    \item \textbf{Verbe introducteur au présent :} \eng{He \textbf{says} he \textbf{is} tired.}
    \item \textbf{Modalités \eng{could, would, should, might, ought to, used to} :} inchangés.
\end{itemize}

\textbf{Marqueurs spatio-temporels (changement obligatoire si verbe introducteur au passé) :}
\begin{center}
\begin{tabularx}{\linewidth}{|X|X|}
\hline
\rowcolor{popblueL}
\textbf{Direct} & \textbf{Indirect} \\
\hline
\eng{now} & \eng{then} / \eng{at that moment} \\
\hline
\eng{today} & \eng{that day} \\
\hline
\eng{tomorrow} & \eng{the next day} / \eng{the following day} \\
\hline
\eng{yesterday} & \eng{the day before} / \eng{the previous day} \\
\hline
\eng{next week/month/year} & \eng{the following week/month/year} \\
\hline
\eng{last week/month/year} & \eng{the previous week/month/year} \\
\hline
\eng{ago} & \eng{before} \\
\hline
\eng{here} & \eng{there} \\
\hline
\eng{this / these} & \eng{that / those} \\
\hline
\end{tabularx}
\end{center}
\end{propriete}

\begin{methode}[Transformer une phrase affirmative du direct à l'indirect : le recul des temps (backshift)]
\textbf{Objectif :} Passer correctement du discours direct au discours indirect pour une déclaration, en appliquant la règle du \cle{backshift} (recul des temps) lorsque le verbe introducteur est au passé.
\begin{enumerate}
    \item \textbf{Identifier le verbe introducteur} (ex. \eng{said}, \eng{told me}, \eng{explained}). S'il est au \textbf{passé} (\eng{said}, \eng{told}), le recul des temps est \textbf{obligatoire} en anglais standard (sauf vérité générale). S'il est au présent (\eng{says}), les temps ne changent pas.
    \item \textbf{Repérer le temps du verbe cité} dans le discours direct.
    \item \textbf{Appliquer le tableau de correspondance} (voir \emph{Propriété} ci-dessus) : Present $\to$ Past ; Past $\to$ Past Perfect ; Present Perfect $\to$ Past Perfect ; Will $\to$ Would ; Can $\to$ Could ; Must $\to$ Had to.
    \item \textbf{Adapter les pronoms} (I $\to$ he/she ; my $\to$ his/her ; we $\to$ they, etc.) et les déterminants (this $\to$ that ; these $\to$ those).
    \item \textbf{Supprimer les guillemets}, ajouter la conjonction \eng{that} (souvent omise à l'oral) et remettre la phrase à la forme affirmative (sujet + verbe).
    \item \textbf{Vérifier l'accord des temps} : si le fait rapporté est encore vrai \textbf{au moment où l'on parle}, on \textbf{peut} garder le présent (ex. \eng{He said the Earth \textbf{is} round}).
\end{enumerate}
\textbf{Structure type :} \eng{Subject + reporting verb (past) + (that) + subject + verb (backshifted) + complements.}
\end{methode}

\begin{exoresolu}[Transformation de déclarations : backshift et pronoms]
\textbf{Énoncé :} Transformez les phrases suivantes du discours direct au discours indirect. Le verbe introducteur est au passé.
\begin{enumerate}
    \item \eng{The Minister of Youth Affairs said, “We \textbf{will} launch a new employment programme for graduates in Buea next month.”}
    \item \eng{My sister told me, “I \textbf{have just bought} a new smartphone at the market in Douala.”}
    \item \eng{The teacher explained, “Water \textbf{boils} at 100 degrees Celsius.”}
    \item \eng{The coach shouted, “You \textbf{must} train harder if you \textbf{want} to win the Mount Cameroon Race.”}
\end{enumerate}

\tcblower
\textbf{Solution détaillée :}

\textbf{1.} \eng{The Minister of Youth Affairs said (that) they \textbf{would} launch a new employment programme for graduates in Buea \textbf{the following month}.}
\begin{itemize}
    \item \textit{Raisonnement :} \eng{said} (passé) $\to$ backshift obligatoire. \eng{We} $\to$ \eng{they} (3\textsuperscript{e} personne). \eng{will} $\to$ \eng{would}. \eng{next month} $\to$ \eng{the following month} (marqueur temporel éloigné).
\end{itemize}

\textbf{2.} \eng{My sister told me (that) she \textbf{had just bought} a new smartphone at the market in Douala.}
\begin{itemize}
    \item \textit{Raisonnement :} \eng{told me} (passé) $\to$ backshift. \eng{I} $\to$ \eng{she}. \eng{have bought} (Present Perfect) $\to$ \eng{had bought} (Past Perfect). \eng{just} reste bien placé avant le participe passé. \eng{at the market in Douala} : lieu inchangé car le locuteur y est peut-être encore, mais \eng{here} serait devenu \eng{there}.
\end{itemize}

\textbf{3.} \eng{The teacher explained (that) water \textbf{boils} at 100 degrees Celsius.}
\begin{itemize}
    \item \textit{Raisonnement :} \eng{explained} (passé) mais \textbf{vérité scientifique universelle}. Le présent \eng{boils} \textbf{ne change pas}. Pas de changement de pronom (sujet \eng{water} inchangé).
\end{itemize}

\textbf{4.} \eng{The coach shouted (that) we \textbf{had to} train harder if we \textbf{wanted} to win the Mount Cameroon Race.}
\begin{itemize}
    \item \textit{Raisonnement :} \eng{shouted} (passé) $\to$ backshift. \eng{You} (adressé aux joueurs dont fait partie le rapporteur) $\to$ \eng{we}. \eng{must} (obligation) $\to$ \eng{had to} (pas de \eng{musted}). \eng{want} (Present) $\to$ \eng{wanted} (Past) dans la subordonnée \eng{if}.
\end{itemize}
\textbf{Conclusion :} Le recul des temps affecte le verbe principal et les verbes des subordonnées. \eng{Must} n'a pas de passé $\to$ \eng{had to}.
\end{exoresolu}

\begin{methode}[Rapporter une question (yes/no et Wh-) : ordre des mots et auxiliaires]
\textbf{Objectif :} Transformer une question directe en question indirecte (rapportée) sans faire l'erreur classique de garder l'ordre interrogatif ou l'auxiliaire \eng{do/does/did}.
\begin{enumerate}
    \item \textbf{Identifier le type de question} : \cle{Yes/No question} (commence par auxiliaire/modal) ou \cle{Wh-question} (commence par \eng{who, what, where, when, why, how}).
    \item \textbf{Choisir le verbe introducteur} : \eng{ask}, \eng{wonder}, \eng{want to know}, \eng{inquire}. Souvent suivi de \eng{if} ou \eng{whether} (pour Yes/No) ou du mot interrogatif (pour Wh-).
    \item \textbf{Appliquer le backshift} sur le verbe (si verbe introducteur au passé) : \eng{do/does} $\to$ \eng{did} (disparaît à la forme affirmative) ; \eng{did} $\to$ \eng{had done} ; \eng{is/are} $\to$ \eng{was/were} ; \eng{will} $\to$ \eng{would}.
    \item \textbf{Remettre à l'ordre \textbf{affirmatif} (Sujet + Verbe)} : \textbf{Jamais} d'auxiliaire \eng{do/does/did} en discours indirect. \textbf{Jamais} d'inversion sujet-verbe.
    \item \textbf{Adapter pronoms et marqueurs spatio-temporels} (today $\to$ that day, here $\to$ there, etc.).
    \item \textbf{Ponctuation :} Point final (pas de point d'interrogation), pas de guillemets.
\end{enumerate}
\textbf{Structures types :}
\begin{itemize}
    \item Yes/No : \eng{He asked (me) if/whether + sujet + verbe (backshifted).}
    \item Wh- : \eng{He asked (me) [wh-word] + sujet + verbe (backshifted).}
\end{itemize}
\end{methode}

\begin{exoresolu}[Transformation de questions directes en indirectes]
\textbf{Énoncé :} Rapportez les questions suivantes au discours indirect (verbe introducteur au passé).
\begin{enumerate}
    \item \eng{The journalist asked the Minister, “When \textbf{will} the new bridge over the Wouri River \textbf{be inaugurated}?”}
    \item \eng{My friend asked me, “Do you \textbf{speak} English fluently?”}
    \item \eng{The tourist wondered, “How much \textbf{does} a taxi ride to the airport \textbf{cost}?”}
    \item \eng{The student asked the teacher, “Did you \textbf{correct} our essays yesterday?”}
\end{enumerate}

\tcblower
\textbf{Solution détaillée :}

\textbf{1.} \eng{The journalist asked the Minister when the new bridge over the Wouri River \textbf{would be inaugurated}.}
\begin{itemize}
    \item \textit{Raisonnement :} Question \eng{Wh-} (\eng{when}). Verbe introducteur \eng{asked} (passé). Backshift : \eng{will be inaugurated} (Future Passive) $\to$ \eng{would be inaugurated} (Conditional Passive). \textbf{Ordre affirmatif} : \eng{the bridge would be inaugurated} (pas \eng{would the bridge be}).
\end{itemize}

\textbf{2.} \eng{My friend asked me \textbf{if/whether} I \textbf{spoke} English fluently.}
\begin{itemize}
    \item \textit{Raisonnement :} Question \eng{Yes/No} (auxiliaire \eng{Do}). Introducteur : \eng{asked me if/whether}. Backshift : \eng{Do you speak} (Present Simple) $\to$ \eng{I spoke} (Past Simple). \textbf{Attention :} L'auxiliaire \eng{do} disparaît, le verbe principal prend la marque du passé (\eng{spoke}). \eng{you} $\to$ \eng{I}.
\end{itemize}

\textbf{3.} \eng{The tourist wondered \textbf{how much} a taxi ride to the airport \textbf{cost}.}
\begin{itemize}
    \item \textit{Raisonnement :} Question \eng{Wh-} (\eng{how much}). Verbe \eng{wondered} (passé). Backshift : \eng{does cost} (Present Simple emphatique/question) $\to$ \eng{cost} (Past Simple). \textbf{Ordre affirmatif} : \eng{a taxi ride cost} (pas \eng{did a taxi ride cost}). Sujet \eng{a taxi ride} inchangé (3\textsuperscript{e} pers).
\end{itemize}

\textbf{4.} \eng{The student asked the teacher \textbf{if/whether} he/she \textbf{had corrected} their essays \textbf{the day before}.}
\begin{itemize}
    \item \textit{Raisonnement :} Question \eng{Yes/No} (\eng{Did}). Introducteur \eng{asked if/whether}. Backshift : \eng{Did you correct} (Past Simple) $\to$ \eng{had corrected} (Past Perfect). \eng{you} $\to$ \eng{he/she} (le prof). \eng{our} $\to$ \eng{their} (les essais des élèves). \eng{yesterday} $\to$ \eng{the day before} (ou \eng{the previous day}).
\end{itemize}
\textbf{Conclusion :} La règle d'or : \textbf{ORDRE AFFIRMATIF} (Sujet + Verbe) et \textbf{PAS DE DO/DOES/DID}.
\end{exoresolu}

\begin{methode}[Rapporter un ordre, une demande ou un conseil : l'infinitif]
\textbf{Objectif :} Transformer l'impératif (ordre, conseil, demande) en discours indirect en utilisant la structure \cle{verbe introducteur + objet + infinitif (avec to)}.
\begin{enumerate}
    \item \textbf{Identifier la fonction illocutoire} : Ordre (\eng{Order}), Demande polie (\eng{Request}), Conseil (\eng{Advice}), Avertissement (\eng{Warn}), Interdiction (\eng{Forbid}).
    \item \textbf{Choisir le verbe introducteur adapté} :
    \begin{itemize}
        \item Ordre : \eng{tell}, \eng{order}, \eng{command}.
        \item Demande : \eng{ask}, \eng{request}, \eng{beg}.
        \item Conseil : \eng{advise}, \eng{recommend}, \eng{urge}.
        \item Interdiction/Négation : \eng{tell} + \textbf{not} + infinitif / \eng{forbid} + objet + \eng{to} + infinitif / \eng{warn} + objet + \eng{against} + \eng{-ing}.
    \end{itemize}
    \item \textbf{Structure unique} : \eng{Reporting verb + Object (me, him, us, the students...) + TO + Base Verbale}.
    \item \textbf{Gérer la négation} : En direct : \eng{“Don't go!”} $\to$ Indirect : \eng{He told me \textbf{not to go}}. (Le \eng{not} se place \textbf{avant} \eng{to}).
    \item \textbf{Adapter pronoms et marqueurs} comme d'habitude.
\end{enumerate}
\textbf{Attention :} \eng{Suggest} ne s'utilise \textbf{PAS} avec l'infinitif (\eng{*He suggested me to go} $\times$). \eng{Suggest} + \eng{that} + \eng{should} / subjonctif / \eng{-ing} (\eng{He suggested going} / \eng{He suggested that I (should) go}).
\end{methode}

\begin{exoresolu}[Transformation d'impératifs (ordres, demandes, interdictions)]
\textbf{Énoncé :} Rapportez les phrases suivantes au discours indirect en choisissant le verbe introducteur le plus approprié.
\begin{enumerate}
    \item \eng{The policeman said to the driver, “Show me your driving licence.”}
    \item \eng{The mother said to her son, “Don't play near the bendskin station.”}
    \item \eng{My friend said, “Please, help me carry this bag of cocoa.”}
    \item \eng{The doctor said to the patient, “You should stop smoking.”} (Conseil avec modal)
\end{enumerate}

\tcblower
\textbf{Solution détaillée :}

\textbf{1.} \eng{The policeman \textbf{ordered} the driver \textbf{to show} him his driving licence.}
\begin{itemize}
    \item \textit{Raisonnement :} Impératif \eng{Show} = Ordre. Verbe introducteur : \eng{ordered} (plus précis que \eng{told}). Structure : \eng{ordered} + \textbf{objet} (\eng{the driver}) + \textbf{to} + \textbf{base verbale} (\eng{show}). \eng{me} $\to$ \eng{him} (le policier). \eng{your} $\to$ \eng{his} (celui du chauffeur).
\end{itemize}

\textbf{2.} \eng{The mother \textbf{told} her son \textbf{not to play} near the bendskin station.}
\begin{itemize}
    \item \textit{Raisonnement :} Impératif négatif \eng{Don't play}. Verbe introducteur : \eng{told} (standard pour ordre/conseil parental). Structure : \eng{told} + \textbf{objet} (\eng{her son}) + \textbf{not to} + \textbf{base verbale} (\eng{play}). \textbf{Règle :} \eng{not} avant \eng{to}. Alternative : \eng{warned her son against playing} / \eng{forbade her son to play}.
\end{itemize}

\textbf{3.} \eng{My friend \textbf{asked} me \textbf{to help} him carry that bag of cocoa.}
\begin{itemize}
    \item \textit{Raisonnement :} \eng{Please} + Impératif = Demande polie (\eng{Request}). Verbe introducteur : \eng{asked}. Structure : \eng{asked} + \textbf{objet} (\eng{me}) + \textbf{to} + \textbf{base verbale} (\eng{help}). \eng{this} $\to$ \eng{that}. \eng{me} (objet de \eng{help}) $\to$ \eng{him} (mon ami).
\end{itemize}

\textbf{4.} \eng{The doctor \textbf{advised} the patient \textbf{to stop} smoking.}
\begin{itemize}
    \item \textit{Raisonnement :} \eng{You should stop} = Conseil (\eng{Advice}). Verbe introducteur : \eng{advised}. Structure : \eng{advised} + \textbf{objet} (\eng{the patient}) + \textbf{to} + \textbf{base verbale} (\eng{stop}). Le modal \eng{should} disparaît, remplacé par l'infinitif \eng{to stop}. \eng{stop smoking} (gérondif après \eng{stop} sens “arrêter de”) conservé.
\end{itemize}
\textbf{Conclusion :} L'infinitif \eng{to + BV} est la clé. N'oubliez jamais l'\textbf{objet} (me, him, them) après \eng{tell, ask, advise, order, warn}.
\end{exoresolu}

\begin{methode}[Choisir le bon verbe introducteur et gérer les “Reporting Verbs” patterns]
\textbf{Objectif :} Varier les verbes introducteurs (\eng{say, tell, ask} ne suffisent pas au Bac) et maîtriser leur construction grammaticale spécifique (patterns).
\begin{enumerate}
    \item \textbf{Analyser l'intention du locuteur} : Informer ? Ordonner ? Promettre ? Refuser ? S'excuser ? Accuser ? Suggérer ?
    \item \textbf{Sélectionner le verbe précis} dans la liste des \cle{Reporting Verbs} classés par structure (voir tableau ci-dessous).
    \item \textbf{Appliquer le pattern (structure) imposé par ce verbe} :
    \begin{itemize}
        \item \eng{Verb + that-clause (backshift)} : \eng{claim, explain, promise, deny, admit, complain}.
        \item \eng{Verb + Object + to-infinitive} : \eng{tell, ask, advise, warn, order, beg, persuade, remind}.
        \item \eng{Verb + to-infinitive (pas d'objet)} : \eng{promise, offer, agree, refuse, threaten, decide}.
        \item \eng{Verb + -ing} : \eng{deny, admit, suggest, recommend, regret, avoid}.
        \item \eng{Verb + Object + that-clause (should/subjunctive)} : \eng{advise, recommend, propose} (formel).
        \item \eng{Verb + Preposition + -ing} : \eng{apologize for, insist on, accuse of, blame for, congratulate on}.
    \end{itemize}
    \item \textbf{Vérifier la concordance des temps} (backshift si verbe introducteur au passé).
    \item \textbf{Attention aux “Faux Amis” de structure} : \eng{Suggest} $\neq$ \eng{suggérer} (fr) + infinitif. \eng{Explain} $\neq$ \eng{expliquer} (fr) + \eng{to me}. \eng{Explain \textbf{to me} that...} / \eng{Explain \textbf{me} that...} ($\times$).
\end{enumerate}
\end{methode}

\begin{exoresolu}[Variété des verbes introducteurs et patterns grammaticaux]
\textbf{Énoncé :} Rapportez les phrases suivantes en utilisant le verbe introducteur indiqué entre parenthèses (au passé). Respectez la structure imposée par chaque verbe.
\begin{enumerate}
    \item \eng{“I didn't steal the money from the till,” the cashier said. (\textbf{deny})}
    \item \eng{“I will definitely pay you back next week,” he told me. (\textbf{promise})}
    \item \eng{“You should really see a doctor about that cough,” she told me. (\textbf{advise})}
    \item \eng{“Let's go to the beach this weekend,” proposed my brother. (\textbf{suggest})}
    \item \eng{“I'm sorry I broke your phone,” said Alice. (\textbf{apologize})}
\end{enumerate}

\tcblower
\textbf{Solution détaillée :}

\textbf{1.} \eng{The cashier \textbf{denied stealing} the money from the till.}
\begin{itemize}
    \item \textit{Raisonnement :} \eng{Deny} = nier. Pattern : \eng{Verb + -ing} (gérondif). Pas d'objet. Backshift : \eng{didn't steal} (Past) $\to$ \eng{stealing} (gérondif “présent” pour antériorité) ou \eng{having stolen} (plus formel). \eng{denied stealing} est standard.
\end{itemize}

\textbf{2.} \eng{He \textbf{promised to pay} me back the following week.}
\begin{itemize}
    \item \textit{Raisonnement :} \eng{Promise} = promettre. Pattern : \eng{Verb + to-infinitive} (sujet de la promesse = sujet du verbe introducteur). \eng{will pay} $\to$ \eng{to pay} (infinitif). \eng{next week} $\to$ \eng{the following week}. \eng{you} $\to$ \eng{me}.
\end{itemize}

\textbf{3.} \eng{She \textbf{advised me to see} a doctor about that cough.}
\begin{itemize}
    \item \textit{Raisonnement :} \eng{Advise} = conseiller. Pattern : \eng{Verb + Object + to-infinitive}. \eng{should see} (modal conseil) $\to$ \eng{to see} (infinitif). Objet \eng{me} obligatoire après \eng{advised}.
\end{itemize}

\textbf{4.} \eng{My brother \textbf{suggested going} to the beach that weekend / \textbf{suggested that we (should) go} to the beach that weekend.}
\begin{itemize}
    \item \textit{Raisonnement :} \eng{Suggest} = suggérer. \textbf{Interdit :} \eng{*suggested us to go}. Patterns corrects : 1. \eng{suggest + -ing} (\eng{going}) ; 2. \eng{suggest + that + should + BV} / \eng{that + subjonctif} (\eng{we go}). \eng{this weekend} $\to$ \eng{that weekend}.
\end{itemize}

\textbf{5.} \eng{Alice \textbf{apologized for breaking} my phone.}
\begin{itemize}
    \item \textit{Raisonnement :} \eng{Apologize} = s'excuser. Pattern : \eng{Verb + Preposition (for) + -ing}. \eng{broke} (Past) $\to$ \eng{breaking} (gérondif) / \eng{having broken}. \eng{your} $\to$ \eng{my}.
\end{itemize}
\textbf{Conclusion :} La connaissance du \textbf{pattern} (construction) de chaque verbe est cruciale. \eng{Deny/Suggest/Admit/Regret} + \eng{-ing} ; \eng{Promise/Offer/Refuse} + \eng{to V} ; \eng{Advise/Tell/Ask/Warn} + \eng{Obj + to V} ; \eng{Apologize/Insist/Accuse} + \eng{Prep + -ing}.
\end{exoresolu}

\begin{methode}[Comparaison systématique Français $\leftrightarrow$ Anglais : déjouer les pièges de traduction]
\textbf{Objectif :} Éviter les calques syntaxiques du français qui sont des erreurs grammaticales graves en anglais (faux amis structurels).
\begin{enumerate}
    \item \textbf{\eng{Say vs Tell} :} Fr : “Il a dit...” $\to$ \eng{He \textbf{said}...} (pas d'objet) OU \eng{He \textbf{told} \textbf{me}...} (objet \textbf{obligatoire}). \eng{*He said me} = $\times$ \textbf{Erreur majeure}.
    \item \textbf{\eng{Ask vs Tell/Order} :} Fr : “Il m'a demandé de partir.” $\to$ \eng{He \textbf{asked} \textbf{me} \textbf{to} leave} (demande). \eng{He \textbf{told} \textbf{me} \textbf{to} leave} (ordre). Ne pas utiliser \eng{ask} pour un ordre.
    \item \textbf{\eng{Suggest / Recommend} :} Fr : “Il m'a suggéré de partir.” $\to$ \eng{He \textbf{suggested} \textbf{(that) I (should) leave} / \textbf{leaving}}. \textbf{JAMAIS} \eng{*suggested me to leave}.
    \item \textbf{\eng{Explain / Describe / Propose} :} Fr : “Il m'a expliqué...” $\to$ \eng{He \textbf{explained} \textbf{to me} that...}. \textbf{JAMAIS} \eng{*explained me that...}. Idem pour \eng{propose, say, announce, mention}.
    \item \textbf{Concordance des temps (Backshift) :} Fr : “Il a dit qu'il \textbf{viendra}” (Futur). En $\to$ \eng{He said he \textbf{would} come} (Conditionnel). L'anglais \textbf{reculle} presque toujours au passé ; le français garde souvent le futur/présent.
    \item \textbf{Pronoms “On” / “Nous” :} Fr : “On m'a dit...” $\to$ \eng{\textbf{People} say / \textbf{They} say / \textbf{I was told}...} (Passif fréquent). “Nous” (général) $\to$ \eng{People / One / You}.
    \item \textbf{Marqueurs spatio-temporels :} Fr : “Il a dit : ‘Je viens \textbf{ici} \textbf{demain}’.” $\to$ \eng{He said he was coming \textbf{there} \textbf{the next day}}. L'anglais \textbf{exige} le changement (\eng{here$\to$there, now$\to$then, this$\to$that}) ; le français le tolère souvent à l'oral.
\end{enumerate}
\end{methode}

\begin{exoresolu}[Correction d'erreurs typiques de francophones (Traduction/Transformation)]
\textbf{Énoncé :} Les phrases suivantes contiennent des erreurs classiques de francophones. Identifiez l'erreur, expliquez-la et donnez la phrase correcte au discours indirect (verbe introducteur au passé).
\begin{enumerate}
    \item \eng{The principal said the students that they must work harder.}
    \item \eng{My father asked me if did I want to go to the village.}
    \item \eng{She suggested me to apply for the scholarship.}
    \item \eng{He explained me that he will come tomorrow.}
    \item \eng{The thief denied to have stolen the phone.}
\end{enumerate}

\tcblower
\textbf{Solution détaillée :}

\textbf{1.} \eng{The principal \textbf{told} the students (that) they \textbf{had to} work harder.}
\begin{itemize}
    \item \textit{Erreur :} \eng{*said the students}. \textbf{Règle :} \eng{Say} ne prend \textbf{jamais} d'objet personne direct. \eng{Tell} prend un objet (\eng{told the students}). Backshift : \eng{must} $\to$ \eng{had to}.
\end{itemize}

\textbf{2.} \eng{My father asked me \textbf{whether/if I wanted} to go to the village.}
\begin{itemize}
    \item \textit{Erreur :} \eng{*if did I want}. \textbf{Règle :} Ordre \textbf{AFFIRMATIF} en discours indirect (Sujet + Verbe). Pas d'auxiliaire \eng{did}. Backshift : \eng{do you want} (Present) $\to$ \eng{I wanted} (Past).
\end{itemize}

\textbf{3.} \eng{She \textbf{suggested that I (should) apply} / \textbf{suggested applying} for the scholarship.}
\begin{itemize}
    \item \textit{Erreur :} \eng{*suggested me to apply}. \textbf{Règle :} \eng{Suggest} + \eng{that-clause} (+ \eng{should}/subjonctif) OU \eng{-ing}. \textbf{Jamais} d'infinitif, \textbf{jamais} d'objet direct (\eng{me}).
\end{itemize}

\textbf{4.} \eng{He \textbf{explained to me} that he \textbf{would come} \textbf{the next day}.}
\begin{itemize}
    \item \textit{Erreurs :} 1. \eng{*explained me} $\to$ \eng{explained \textbf{to} me}. 2. \eng{will come} $\to$ \eng{would come} (Backshift obligatoire). 3. \eng{tomorrow} $\to$ \eng{the next day} (Marqueur temporel).
\end{itemize}

\textbf{5.} \eng{The thief \textbf{denied stealing} / \textbf{denied having stolen} the phone.}
\begin{itemize}
    \item \textit{Erreur :} \eng{*denied to have stolen}. \textbf{Règle :} \eng{Deny} + \eng{-ing} (Gérondif). Pas d'infinitif.
\end{itemize}
\textbf{Conclusion :} Ces 5 erreurs (\eng{say/tell}, ordre interrogatif, \eng{suggest}, \eng{explain to}, \eng{deny + -ing}) sont les \textbf{plus sanctionnées au Baccalauréat}. Les maîtriser garantit des points faciles.
\end{exoresolu}

\begin{attention}[Les 5 erreurs qui coûtent des points au Bac]
\begin{enumerate}
    \item \textbf{\eng{*He said me / *He asked me that...} (Confusion \eng{Say} vs \eng{Tell} / \eng{Ask}).}
    \cle{Règle} : \eng{Say} (pas d'objet) / \eng{Tell} + \textbf{Objet} / \eng{Ask} + \textbf{Objet} + \eng{if/whether/to}.
    \item \textbf{Garder l'ordre interrogatif ou l'auxiliaire \eng{Do/Did} en discours indirect.}
    \cle{Règle} : \eng{He asked where I \textbf{lived}} (pas \eng{*where did I live} / \eng{*where I did live}).
    \item \textbf{\eng{*Suggest / Recommend / Propose + Objet + Infinitif} (Calque du français “suggérer à qqn de”).}
    \cle{Règle} : \eng{Suggest + that-clause (should) / -ing}. \eng{Advise / Tell / Ask + Objet + to-infinitif}.
    \item \textbf{Oublier le \eng{to} après \eng{Explain / Describe / Propose / Say / Announce + Objet}.}
    \cle{Règle} : \eng{Explain \textbf{to} me that...} / \eng{Propose \textbf{to} me that...}. \eng{Tell} est le \textbf{seul} verbe courant de parole prenant un objet direct sans préposition.
    \item \textbf{Oublier le \textbf{Backshift} (recul des temps) ou les marqueurs spatio-temporels (\eng{here/now/this}).}
    \cle{Règle} : Verbe introducteur au passé $\to$ \textbf{Backshift systématique} (sauf vérité générale). \eng{Here $\to$ There, Now $\to$ Then, Today $\to$ That day, Tomorrow $\to$ The next day, Yesterday $\to$ The day before, Next week $\to$ The following week, This/These $\to$ That/Those}.
\end{enumerate}
\textbf{Astuce de relecture :} Pour chaque phrase rapportée, cochez mentalement : \fbox{Verbe introducteur ?} \fbox{Backshift ?} \fbox{Pronoms ?} \fbox{Marqueurs ?} \fbox{Structure du verbe (Pattern) ?} \fbox{Ordre affirmatif ?}
\end{attention}

\begin{savaistu}[Discours rapporté dans l'anglais camerounais et nigérian]
Dans l'anglais standard (British English), le \cle{backshift} est la norme après un verbe introducteur au passé. Cependant, dans l'anglais camerounais (CamE) et l'anglais nigérian (NigE), influencés par les langues locales et le français, on observe souvent une \textbf{absence de backshift} dans le langage courant : \eng{“He said he \textbf{will} come tomorrow”} (au lieu de \eng{would}). De même, les marqueurs spatio-temporels ne sont pas toujours décalés : \eng{“She said she is coming \textbf{here} \textbf{now}”}. Pour le Baccalauréat et l'écrit formel, \textbf{l'anglais standard avec backshift et décalage des marqueurs est exigé}.
\end{savaistu}

\newpage

\coursec{Exercices}

\courssub{Je vérifie mes connaissances}

\begin{exercice}[QCM : la bonne forme rapportée]
Entoure la bonne réponse dans chaque cas.

\begin{enumerate}
\item He asked me \ldots
\begin{enumerate}
\item where did I live \item where I lived \item where do I live \item where I live
\end{enumerate}
\item She said that she \ldots the letter the day before.
\begin{enumerate}
\item writes \item has written \item had written \item will write
\end{enumerate}
\item The teacher told us \ldots late.
\begin{enumerate}
\item to not be \item not to be \item don't be \item not be
\end{enumerate}
\item “I will visit Buea tomorrow,” Mary said. — Mary said she \ldots visit Buea the next day.
\begin{enumerate}
\item will \item would \item would have \item will have
\end{enumerate}
\item He asked me \ldots I had ever seen Mount Cameroon.
\begin{enumerate}
\item that \item what \item if \item did
\end{enumerate}
\item “Don't touch the cocoa beans!” the farmer said. — The farmer told the children \ldots
\begin{enumerate}
\item don't touch the cocoa beans \item to not touch the cocoa beans \item not to touch the cocoa beans \item that they don't touch the cocoa beans
\end{enumerate}
\end{enumerate}
\end{exercice}

\begin{exercice}[Vrai ou faux ? Justifie]
Indique si chaque affirmation est vraie ou fausse, puis justifie ta réponse en une phrase (en français).

\begin{enumerate}
\item \eng{That} est obligatoire après \eng{said}.
\item On garde l'inversion sujet-verbe dans les questions rapportées (\eng{He asked me where did I live}).
\item \eng{Must} devient toujours \eng{had to} au discours indirect.
\item \eng{Say} et \eng{tell} s'emploient de la même façon : on peut dire \eng{He said me the truth}.
\item \eng{Tomorrow} devient \eng{the next day} ou \eng{the following day} au discours indirect.
\item Une question fermée (réponse \eng{yes/no}) se rapporte avec \eng{if} ou \eng{whether}.
\end{enumerate}
\end{exercice}

\begin{exercice}[Vocabulaire : les verbes introducteurs]
Complète chaque définition avec le verbe introducteur qui convient, choisi dans la liste : \eng{warn — deny — promise — admit — suggest — refuse}.

\begin{enumerate}
\item \eng{\ldots} : to say that you did something wrong. (avouer)
\item \eng{\ldots} : to say that something is not true. (nier)
\item \eng{\ldots} : to tell someone that something bad may happen. (avertir)
\item \eng{\ldots} : to say that you will certainly do something. (promettre)
\item \eng{\ldots} : to say no to a request or an offer. (refuser)
\item \eng{\ldots} : to give an idea or a plan for someone to consider. (suggérer)
\end{enumerate}
\end{exercice}

\begin{exercice}[Conjugaison : le recul des temps]
Réécris le verbe entre parenthèses au temps qui convient, comme si la phrase était rapportée après \eng{He said that\ldots}.

\begin{enumerate}
\item “I \emph{(work)} in a bank in Douala.”
\item “We \emph{(finish)} our exams last week.”
\item “She \emph{(never / see)} a white man before.”
\item “They \emph{(build)} a new school next year.”
\item “I \emph{(can)} speak three languages.”
\item “The match \emph{(start)} at four o'clock.”
\end{enumerate}
\end{exercice}

\courssub{Je m'entraîne}

\begin{exercice}[Traduction vers l'anglais]
Traduis les phrases suivantes en anglais, au discours indirect.

\begin{enumerate}
\item Il a dit qu'il était fatigué et qu'il voulait rentrer à la maison.
\item Elle m'a demandé si j'avais déjà visité Bamenda.
\item Le professeur nous a dit de ne pas oublier nos dictionnaires.
\item Ma mère m'a demandé où j'avais mis son téléphone.
\item Il a promis qu'il m'aiderait le lendemain.
\item Le chauffeur a dit que le car partait dans dix minutes.
\end{enumerate}
\end{exercice}

\begin{exercice}[Transformation : discours direct vers indirect]
Mets les phrases suivantes au discours indirect en commençant par les mots indiqués.

\begin{enumerate}
\item “I am preparing my lesson,” the student said. — The student said \ldots
\item “Do you like Cameroonian music?” she asked me. — She asked me \ldots
\item “Where did you buy this shirt?” he asked his brother. — He asked his brother \ldots
\item “Open your books at page 45,” the teacher said. — The teacher told us \ldots
\item “I will call you tomorrow,” Paul said. — Paul said \ldots
\item “Don't walk on the grass!” the guard said to the children. — The guard told the children \ldots
\item “We have won the match,” the captain said. — The captain said \ldots
\item “Can you help me carry this bag?” the old woman asked the boy. — The old woman asked the boy \ldots
\end{enumerate}
\end{exercice}

\begin{textebox}[At the police station]
Last Saturday, a young man was taken to the police station in Yaoundé. The inspector \ldots\ him why he \ldots\ (run) away when he saw the officers. The young man first \ldots\ nothing, but later he \ldots\ that he \ldots\ (steal) a mobile phone at the market the day before. The inspector \ldots\ him that stealing \ldots\ (be) a serious crime and \ldots\ him \ldots\ (not / do) it again. Finally, the young man \ldots\ that he \ldots\ (return) the phone to its owner the following day.
\end{textebox}

\begin{exercice}[Texte à trous : verbes introducteurs et formes verbales]
Complète le texte « At the police station » ci-dessus avec le verbe introducteur et la forme verbale qui conviennent. Verbes introducteurs à utiliser une seule fois chacun : \eng{said — told — asked — warned — admitted — promised}.
\end{exercice}

\begin{exercice}[Appariements : phrase directe et phrase rapportée]
Associe chaque phrase au discours direct (1--6) à sa version rapportée (a--f).

\textbf{Discours direct :}
\begin{enumerate}
\item “I have lost my voter's card.”
\item “Will you vote in the next election?”
\item “Don't cross the road here.”
\item “I can explain the problem.”
\item “Why are you late?”
\item “We are cleaning the classroom now.”
\end{enumerate}

\textbf{Discours indirect :}
\begin{enumerate}[label=\alph*.]
\item She asked me if I would vote in the next election.
\item He said he had lost his voter's card.
\item The teacher asked the boy why he was late.
\item She told me not to cross the road there.
\item He said he could explain the problem.
\item They said they were cleaning the classroom then.
\end{enumerate}
\end{exercice}

\courssub{Je me prépare au Bac}

\begin{textebox}[A dialogue with the chief]
Last week, some young people from a village near Bafoussam went to see their traditional chief. They told him that the road to their quarter was in a very bad state and that the bendskin riders refused to go there when it rained. The chief listened carefully. Then he said that he understood their problem and that he had already written to the council the previous month. He asked the young people if they were ready to take part in community work. They answered that they were ready and that they would bring their friends the following Saturday. The chief smiled and told them not to forget their tools. He added that together they could change their village.
\end{textebox}

\begin{exercice}[Compréhension et langue — type Bac]
Lis le texte « A dialogue with the chief » ci-dessus, puis réponds aux questions.

\textbf{A. Comprehension questions.} Answer in complete sentences, in your own words.
\begin{enumerate}
\item Why did the young people go to see the chief?
\item What had the chief done before their visit?
\item What did the chief ask the young people to do?
\end{enumerate}

\textbf{B. Language work.}
\begin{enumerate}
\item Find in the text three sentences in reported speech and rewrite them in direct speech.
\item Report the following sentence as the chief said it: \eng{“I will speak to the mayor tomorrow.”} (Begin: \eng{The chief said that\ldots})
\item Give the French translation of: \eng{“He asked the young people if they were ready to take part in community work.”}
\end{enumerate}
\end{exercice}

\begin{textebox}[Texte à traduire]
Hier, mon père m'a dit qu'il avait rencontré le directeur de mon école. Il m'a demandé si je travaillais bien en classe. J'ai répondu que je faisais de mon mieux et que j'avais réussi mes derniers examens. Mon père m'a dit de ne pas perdre mon temps et m'a promis qu'il m'achèterait un ordinateur si j'obtenais le baccalauréat.
\end{textebox}

\begin{exercice}[Traduction — type Bac]
Traduis en anglais le passage « Texte à traduire » ci-dessus.
\end{exercice}

\begin{exercice}[Rédaction guidée — type Bac]
Tu es journaliste pour le journal de ton lycée. Tu as assisté à un dialogue entre un policier et un conducteur de moto-taxi (\eng{bendskin}) à Douala. Rédige un court article d'environ \textbf{150 mots} dans lequel tu rapportes ce dialogue \textbf{au discours indirect}.

Consignes :
\begin{enumerate}
\item Commence par une phrase d'introduction (où, quand, qui).
\item Rapporte au moins \textbf{six} échanges : deux déclarations, deux questions (une fermée, une ouverte) et deux ordres ou conseils.
\item Utilise au moins \textbf{quatre verbes introducteurs différents} (\eng{said, told, asked, warned, ordered, advised\ldots}).
\item Termine par une phrase de conclusion (la fin de l'incident).
\end{enumerate}

Voici le début du dialogue que tu as entendu, pour t'inspirer :

Policeman: “Stop! Where is your driving licence?”
Rider: “I have forgotten it at home, sir.”
Policeman: “Why are you not wearing a helmet?”
Rider: “I will buy one tomorrow, I promise.”
Policeman: “Don't ride without a helmet again. It is dangerous.”
Rider: “I understand, sir. I will respect the rules.”
\end{exercice}



\newpage

\coursec{Corrigés des exercices}

\begin{corrige}[Exercice 1 — QCM : la bonne forme rapportée]
\begin{enumerate}
\item \textbf{b} — \eng{where I lived} : dans une question rapportée, on revient à l'ordre sujet + verbe (pas d'inversion) et le temps recule (\eng{live} $\rightarrow$ \eng{lived}).
\item \textbf{c} — \eng{had written} : le present perfect recule au past perfect ; \eng{the day before} confirme le recul.
\item \textbf{b} — \eng{not to be} : la négation se place \emph{avant} l'infinitif (\eng{tell somebody not to do something}).
\item \textbf{b} — \eng{would} : \eng{will} recule à \eng{would}, et \eng{tomorrow} devient \eng{the next day}.
\item \textbf{c} — \eng{if} : question fermée (réponse \eng{yes/no}) rapportée avec \eng{if} ou \eng{whether}.
\item \textbf{c} — \eng{not to touch the cocoa beans} : ordre négatif rapporté = \eng{tell + complément + not to + infinitif}.
\end{enumerate}
\end{corrige}

\begin{corrige}[Exercice 2 — Vrai ou faux ? Justifie]
\begin{enumerate}
\item \textbf{Faux.} \eng{That} est facultatif après \eng{said} : \eng{He said (that) he was tired.} « Il a dit qu'il était fatigué. »
\item \textbf{Faux.} Dans une question rapportée, on revient à l'ordre de la phrase affirmative : \eng{He asked me where I lived.} L'inversion disparaît et il n'y a plus d'auxiliaire \eng{do/did}.
\item \textbf{Faux.} \eng{Must} devient généralement \eng{had to}, mais il peut rester inchangé quand il exprime un conseil ou une déduction : \eng{He said I must be careful.}
\item \textbf{Faux.} \eng{Say} ne prend pas de complément de personne direct : on dit \eng{He said the truth to me} ou \eng{He told me the truth}. Jamais \eng{He said me}.
\item \textbf{Vrai.} \eng{Tomorrow} $\rightarrow$ \eng{the next day / the following day}.
\item \textbf{Vrai.} \eng{“Are you ready?”} $\rightarrow$ \eng{He asked me if/whether I was ready.}
\end{enumerate}
\end{corrige}

\begin{corrige}[Exercice 3 — Vocabulaire : les verbes introducteurs]
\begin{enumerate}
\item \eng{admit} (avouer) \item \eng{deny} (nier) \item \eng{warn} (avertir) \item \eng{promise} (promettre) \item \eng{refuse} (refuser) \item \eng{suggest} (suggérer)
\end{enumerate}
Rappel des constructions : \eng{admit/deny/suggest + V-ing} (\eng{He admitted stealing the phone.}) ; \eng{promise/refuse + to + infinitif} (\eng{She promised to help me.}) ; \eng{warn + somebody + (not) to + infinitif} (\eng{They warned us not to go there.}).
\end{corrige}

\begin{corrige}[Exercice 4 — Conjugaison : le recul des temps]
\begin{enumerate}
\item He said that he \textbf{worked} in a bank in Douala. (present simple $\rightarrow$ past simple)
\item He said that they \textbf{had finished} their exams the week before. (past simple $\rightarrow$ past perfect)
\item He said that she \textbf{had never seen} a white man before. (present perfect $\rightarrow$ past perfect)
\item He said that they \textbf{would build} a new school the following year. (\eng{will} $\rightarrow$ \eng{would})
\item He said that he \textbf{could} speak three languages. (\eng{can} $\rightarrow$ \eng{could})
\item He said that the match \textbf{started} at four o'clock. (present simple $\rightarrow$ past simple)
\end{enumerate}
\end{corrige}

\begin{corrige}[Exercice 5 — Traduction vers l'anglais]
\begin{enumerate}
\item \eng{He said (that) he was tired and (that) he wanted to go home.}
\item \eng{She asked me if I had ever visited Bamenda.}
\item \eng{The teacher told us not to forget our dictionaries.}
\item \eng{My mother asked me where I had put her phone.}
\item \eng{He promised (that) he would help me the next day.}
\item \eng{The driver said (that) the bus was leaving in ten minutes.}
\end{enumerate}
Points de vigilance : « si » interrogatif = \eng{if} (jamais \eng{that}) ; « de ne pas oublier » = \eng{not to forget} ; « le lendemain » = \eng{the next day}.
\end{corrige}

\begin{corrige}[Exercice 6 — Transformation : discours direct vers indirect]
\begin{enumerate}
\item The student said (that) he was preparing his lesson. (\eng{am preparing} $\rightarrow$ \eng{was preparing})
\item She asked me if I liked Cameroonian music. (question fermée : \eng{if} + ordre affirmatif)
\item He asked his brother where he had bought that shirt. (past simple $\rightarrow$ past perfect ; \eng{this} $\rightarrow$ \eng{that})
\item The teacher told us to open our books at page 45. (ordre : \eng{tell + complément + to + infinitif})
\item Paul said (that) he would call me the next day. (\eng{will} $\rightarrow$ \eng{would} ; \eng{tomorrow} $\rightarrow$ \eng{the next day})
\item The guard told the children not to walk on the grass. (ordre négatif : \eng{not to + infinitif})
\item The captain said (that) they had won the match. (present perfect $\rightarrow$ past perfect)
\item The old woman asked the boy if he could help her carry that bag. (\eng{can} $\rightarrow$ \eng{could})
\end{enumerate}
\end{corrige}

\begin{corrige}[Exercice 7 — Texte à trous : verbes introducteurs et formes verbales]
The inspector \textbf{asked} him why he \textbf{had run} away when he saw the officers. The young man first \textbf{said} nothing, but later he \textbf{admitted} that he \textbf{had stolen} a mobile phone at the market the day before. The inspector \textbf{told} him that stealing \textbf{was} a serious crime and \textbf{warned} him \textbf{not to do} it again. Finally, the young man \textbf{promised} that he \textbf{would return} the phone to its owner the following day.

Justifications : \eng{asked} introduit une question rapportée (\eng{why he had run}, past perfect car l'action est antérieure) ; \eng{admitted that he had stolen} (past perfect + \eng{the day before}) ; \eng{warned him not to do} (avertissement négatif) ; \eng{promised that he would return} (\eng{will} $\rightarrow$ \eng{would}).
\end{corrige}

\begin{corrige}[Exercice 8 — Appariements : phrase directe et phrase rapportée]
1 -- b ; 2 -- a ; 3 -- d ; 4 -- e ; 5 -- c ; 6 -- f.

Observe les indices : recul des temps (\eng{have lost} $\rightarrow$ \eng{had lost} ; \eng{can} $\rightarrow$ \eng{could} ; \eng{are cleaning} $\rightarrow$ \eng{were cleaning}), question fermée avec \eng{if} (\eng{Will you vote\ldots} $\rightarrow$ \eng{if I would vote}), question ouverte sans inversion (\eng{why he was late}), ordre négatif avec \eng{not to}, et changement des marqueurs spatio-temporels (\eng{here} $\rightarrow$ \eng{there} ; \eng{now} $\rightarrow$ \eng{then}).
\end{corrige}

\begin{corrige}[Exercice 9 — Compréhension et langue — type Bac]
Barème indicatif (sur 20) : A. Compréhension 6 pts (2 pts par réponse correcte et rédigée) ; B.1 6 pts (2 pts par phrase correctement retransformée) ; B.2 4 pts ; B.3 4 pts.

\textbf{A. Comprehension.}
\begin{enumerate}
\item \eng{They went to see the chief because the road to their quarter was in a very bad state and the bendskin riders refused to go there when it rained.}
\item \eng{Before their visit, the chief had already written to the council.}
\item \eng{He asked them if they were ready to take part in community work and told them not to forget their tools.}
\end{enumerate}

\textbf{B. Language work.}
\begin{enumerate}
\item Exemples possibles : \eng{They told him that the road to their quarter was in a very bad state} $\rightarrow$ \eng{“The road to our quarter is in a very bad state,” they told him.} ; \eng{He said that he had already written to the council the previous month} $\rightarrow$ \eng{“I wrote to the council last month,” he said.} ; \eng{He told them not to forget their tools} $\rightarrow$ \eng{“Don't forget your tools,” he told them.}
\item \eng{The chief said that he would speak to the mayor the next day.} (\eng{will} $\rightarrow$ \eng{would} ; \eng{tomorrow} $\rightarrow$ \eng{the next day})
\item « Il a demandé aux jeunes s'ils étaient prêts à participer aux travaux communautaires. »
\end{enumerate}
\end{corrige}

\begin{corrige}[Exercice 10 — Traduction — type Bac]
Barème indicatif (sur 10) : 2 pts par phrase correctement traduite ; retirer 0,5 pt par erreur de temps, de verbe introducteur ou d'ordre des mots.

\eng{Yesterday, my father told me that he had met the headmaster of my school. He asked me if I was working hard in class. I answered that I was doing my best and that I had passed my last examinations. My father told me not to waste my time and promised me that he would buy me a computer if I passed the baccalauréat.}

Points de vigilance : \eng{told me that he had met} (past perfect, action antérieure) ; \eng{asked me if I was working} (question fermée, sans inversion) ; \eng{told me not to waste} (négation avant l'infinitif) ; \eng{promised me that he would buy} (\eng{will} $\rightarrow$ \eng{would}) ; « le directeur de mon école » = \eng{the headmaster of my school}.
\end{corrige}

\begin{corrige}[Exercice 11 — Rédaction guidée — modèle de production]
Barème indicatif (sur 20) : respect de la consigne et de la longueur 4 pts ; exactitude du discours indirect (recul des temps, pronoms, marqueurs) 8 pts ; variété des verbes introducteurs 4 pts ; organisation (introduction, conclusion) et qualité de la langue 4 pts.

\eng{Last Monday morning, at the Akwa junction in Douala, I witnessed a dialogue between a policeman and a bendskin rider. The policeman ordered the rider to stop and asked him where his driving licence was. The young man admitted that he had forgotten it at home. Then the policeman asked him why he was not wearing a helmet. The rider promised that he would buy one the following day. The policeman warned him not to ride without a helmet again and explained that it was very dangerous. He also advised him to always carry his documents with him. The rider answered that he understood and that he would respect the rules. Finally, the policeman let him go with a warning, and the rider promised to be a more careful citizen in the future.}

Points de vigilance : ce modèle contient une introduction (où, quand, qui), six échanges rapportés, sept verbes introducteurs (\eng{ordered, asked, admitted, promised, warned, advised, answered}) et une conclusion. Dans ta propre rédaction, vérifie : le recul des temps (\eng{have forgotten} $\rightarrow$ \eng{had forgotten}), l'ordre des mots dans les questions rapportées (\eng{where his driving licence was}), et \eng{not to} pour les ordres négatifs (\eng{warned him not to ride}).
\end{corrige}

\newpage

\coursec{Fiche bilan}

\begin{popfigure}
\begin{center}
\begin{center}\tcbox[colback=popblue,colframe=popblue,coltext=white,arc=9pt,boxsep=4pt,left=6pt,right=6pt]{\bfseries\small Direct \& Indirect Speech}\end{center}
\vspace{-2pt}
\begin{tcbraster}[raster columns=3,raster equal height=rows,raster column skip=6pt,raster row skip=6pt,enhanced,blankest]
\begin{tcolorbox}[enhanced,colback=poporange,colframe=poporange,coltext=white,boxrule=0.9pt,arc=3pt,left=4pt,right=4pt,top=3pt,bottom=3pt,boxsep=1pt,halign=left]\bfseries\scriptsize Backshift des temps\end{tcolorbox}
\begin{tcolorbox}[enhanced,colback=popgreen,colframe=popgreen,coltext=white,boxrule=0.9pt,arc=3pt,left=4pt,right=4pt,top=3pt,bottom=3pt,boxsep=1pt,halign=left]\bfseries\scriptsize Pronoms \& possessifs\end{tcolorbox}
\begin{tcolorbox}[enhanced,colback=poppurple,colframe=poppurple,coltext=white,boxrule=0.9pt,arc=3pt,left=4pt,right=4pt,top=3pt,bottom=3pt,boxsep=1pt,halign=left]\bfseries\scriptsize Marqueurs spatio-temporels\end{tcolorbox}
\begin{tcolorbox}[enhanced,colback=poppink,colframe=poppink,coltext=white,boxrule=0.9pt,arc=3pt,left=4pt,right=4pt,top=3pt,bottom=3pt,boxsep=1pt,halign=left]\bfseries\scriptsize Questions : if / whether, wh-\end{tcolorbox}
\begin{tcolorbox}[enhanced,colback=popgold,colframe=popgold,coltext=white,boxrule=0.9pt,arc=3pt,left=4pt,right=4pt,top=3pt,bottom=3pt,boxsep=1pt,halign=left]\bfseries\scriptsize Ordres : tell / ask + to-inf.\end{tcolorbox}
\begin{tcolorbox}[enhanced,colback=popdark,colframe=popdark,coltext=white,boxrule=0.9pt,arc=3pt,left=4pt,right=4pt,top=3pt,bottom=3pt,boxsep=1pt,halign=left]\bfseries\scriptsize Verbes introducteurs\end{tcolorbox}
\end{tcbraster}

\end{center}
\legende{Carte mentale de la leçon : rapporter les paroles en anglais}
\end{popfigure}

\begin{bilan}
\textbf{1. Lexique clé (English / Français)}
\begin{center}
\begin{tabularx}{\linewidth}{|X|X|}
\hline
\rowcolor{popblueL} \textbf{English} & \textbf{Français} \\ \hline
direct / quoted speech & discours direct \\ \hline
indirect / reported speech & discours indirect \\ \hline
reporting verb (say, tell, ask) & verbe introducteur \\ \hline
backshift of tenses & recul (concordance) des temps \\ \hline
statement & déclaration, affirmation \\ \hline
to quote / to report & citer / rapporter \\ \hline
inverted commas / quotation marks & guillemets \\ \hline
\end{tabularx}
\end{center}

\textbf{2. Règles à connaître par cœur}
\begin{center}
\begin{tabularx}{\linewidth}{|l|X|X|}
\hline
\rowcolor{popblueL} \textbf{Direct} & \textbf{Indirect (backshift)} & \textbf{Exemple} \\ \hline
present simple & past simple & “I vote.” $\rightarrow$ He said he voted. \\ \hline
present continuous & past continuous & “I am working.” $\rightarrow$ She said she was working. \\ \hline
present perfect & past perfect & “I have voted.” $\rightarrow$ He said he had voted. \\ \hline
past simple & past perfect & “We won.” $\rightarrow$ They said they had won. \\ \hline
will & would & “I will help.” $\rightarrow$ She said she would help. \\ \hline
can / may & could / might & “I can come.” $\rightarrow$ He said he could come. \\ \hline
must (obligation) & had to & “You must vote.” $\rightarrow$ He said I had to vote. \\ \hline
\end{tabularx}
\end{center}

\begin{itemize}
\item \textbf{Pas de backshift} si le verbe introducteur est au présent (\eng{he says}) ou si le fait rapporté est encore vrai / une vérité générale : \eng{He said the earth is round.}
\item \textbf{Pronoms} : on adapte à la nouvelle situation (\eng{I} $\rightarrow$ \eng{he/she} ; \eng{my} $\rightarrow$ \eng{his/her} ; \eng{we} $\rightarrow$ \eng{they}).
\item \textbf{Marqueurs} : \eng{now} $\rightarrow$ \eng{then} ; \eng{today} $\rightarrow$ \eng{that day} ; \eng{yesterday} $\rightarrow$ \eng{the day before} ; \eng{tomorrow} $\rightarrow$ \eng{the next day} ; \eng{here} $\rightarrow$ \eng{there} ; \eng{this} $\rightarrow$ \eng{that} ; \eng{ago} $\rightarrow$ \eng{before}.
\item \textbf{Questions} : ordre sujet + verbe, pas d'inversion, pas de point d'interrogation. Fermée : \eng{asked if/whether} ; ouverte : \eng{asked wh- + sujet + verbe}.
\item \textbf{Ordres / demandes} : \eng{told / ordered / asked + objet + (not) to + BV} : \eng{“Sit down!” $\rightarrow$ She told me to sit down.}
\item \textbf{say vs tell} : \eng{say (something to somebody)} ; \eng{tell somebody something}. Jamais \eng{say me}.
\end{itemize}

\textbf{3. Méthode express (4 étapes)}
\begin{enumerate}
\item Repérer le type de phrase : déclaration, question, ordre.
\item Choisir le verbe introducteur : \eng{said / told / asked} (+ objet).
\item Appliquer le backshift et adapter pronoms + marqueurs spatio-temporels.
\item Vérifier : pas de guillemets, pas d'inversion dans les questions, pas de \eng{that} après \eng{ask}.
\end{enumerate}
\end{bilan}

\begin{aretenir}[Checklist avant le Bac]
\begin{itemize}
\item[$\square$] Je connais le tableau du backshift (present $\rightarrow$ past, will $\rightarrow$ would, can $\rightarrow$ could).
\item[$\square$] Je sais quand le backshift n'est \emph{pas} obligatoire (vérité encore valable, verbe au présent).
\item[$\square$] Je transforme correctement les pronoms et les possessifs (\eng{I/my} $\rightarrow$ \eng{he/his}).
\item[$\square$] Je connais les changements de marqueurs : \eng{today} $\rightarrow$ \eng{that day}, \eng{tomorrow} $\rightarrow$ \eng{the next day}.
\item[$\square$] Je rapporte une question fermée avec \eng{if} ou \eng{whether}, sans inversion.
\item[$\square$] Je rapporte une question ouverte avec \eng{wh- + sujet + verbe}.
\item[$\square$] Je rapporte un ordre avec \eng{tell/ask/order + objet + to-infinitive}.
\item[$\square$] Je rapporte un ordre négatif avec \eng{not to} : \eng{He told me not to go.}
\item[$\square$] Je distingue \eng{say} et \eng{tell} sans erreur.
\item[$\square$] Je supprime guillemets et point d'interrogation au discours indirect.
\item[$\square$] J'évite les calques du français : pas de \eng{say that if}, pas de futur après \eng{he said that} quand \eng{would} est requis.
\item[$\square$] Je sais utiliser des verbes introducteurs variés : \eng{promise, advise, warn, suggest, apologize for, invite}.
\end{itemize}
\end{aretenir}

\findecours

\end{document}
